% Lesson 25 — Listening: Texts on keeping informed about health problems related to drugs, tobacco, alcohol abuse, and aging — Anglais Tle SH — Cours signé OnBuch+
% Programme officiel MINESEC. Compiler avec : tectonic EN25-listening-texts-on-keeping-informed-about-health-proble.tex
\newcommand{\DOCMATIERE}{Anglais}
\newcommand{\DOCNIVEAU}{Tle SH}
\newcommand{\DOCMODULE}{Chapter 3 : Environment, Well-being and Health}
\newcommand{\DOCLECON}{EN25}
\newcommand{\DOCTITRE}{Lesson 25 — Listening: Texts on keeping informed about health problems related to drugs, tobacco, alcohol abuse, and aging}
% OnBuch+ — préambule « cours signés » ENRICHI (compilé avec tectonic / XeLaTeX)
% Système visuel « Pop » d'OnBuch+ : fond crème (jamais blanc), bordures encre
% épaisses, ombres « dures » décalées, titres Archivo Black en capitales.
% Version MATHÉMATIQUES : graphes (pgfplots/tikz), boîtes pédagogiques
% (définition, propriété, méthode, attention, le savais-tu, exercice résolu,
% exercice, TP, bilan), page de garde et signature OnBuch+ ancrée en bas.
%
% Macros attendues dans main.tex AVANT \input{preamble} :
%   \DOCMATIERE  \DOCNIVEAU  \DOCMODULE  \DOCLECON (numéro)  \DOCTITRE
\documentclass[11pt]{article}

\usepackage{fontspec}
\usepackage[english,french]{babel} % français = langue principale ; l'anglais via \eng{..} / textebox
\usepackage[a4paper,margin=2cm,top=3.4cm,bottom=2.8cm,headheight=1.4cm,footskip=1.3cm]{geometry}
\usepackage{amsmath,amssymb,amsfonts}
\usepackage{mathtools} % \xmapsto, \coloneqq... souvent utilisés par les modèles
\usepackage{cancel} % simplifications barrées (bilans de forces)
% \abs{z} (module, valeur absolue) et \norm{u} (norme), utilisés par les modèles
\providecommand{\abs}{}\let\abs\relax\DeclarePairedDelimiter{\abs}{\lvert}{\rvert}
\providecommand{\norm}{}\let\norm\relax\DeclarePairedDelimiter{\norm}{\lVert}{\rVert}
\usepackage[table]{xcolor} % [table] : fournit aussi \rowcolors (alternance de lignes)
\usepackage{pagecolor}
\usepackage{tikz}
\usetikzlibrary{arrows.meta,positioning,calc,shapes.geometric,shapes.misc,shapes.arrows,decorations.pathmorphing,decorations.pathreplacing,decorations.markings,patterns,shadows,mindmap,trees,fit,backgrounds,matrix,intersections,angles,quotes}
\usepackage{pgfplots}
\usepackage[version=4]{mhchem} % équations nucléaires \ce{^{235}_{92}U}
\usepackage[european,siunitx]{circuitikz} % schémas électriques
\pgfplotsset{compat=1.18}
\usepgfplotslibrary{fillbetween,groupplots} % aires entre courbes (calcul intégral)
\usepackage{enumitem}
\usepackage{fancyhdr}
% [most] charge toutes les bibliothèques tcolorbox (listings, minted, raster,
% documentation...) et gonfle inutilement la mémoire TeX sur un document long
% (~300 boîtes) : on ne charge que ce qui est réellement utilisé.
\usepackage[skins,breakable]{tcolorbox}
\usepackage{graphicx}
\usepackage{colortbl}
\usepackage{booktabs}
\usepackage{tabularx}
\usepackage{diagbox} % case d'en-tête diagonale des tableaux à double entrée
% Colonnes extensibles centrée / gauche / droite (C, L, R), souvent utilisées par les modèles
\newcolumntype{C}{>{\centering\arraybackslash}X}
\newcolumntype{L}{>{\raggedright\arraybackslash}X}
\newcolumntype{R}{>{\raggedleft\arraybackslash}X}
\usepackage{array}
\usepackage{multicol}
\usepackage{multirow}
\usepackage{siunitx}
% parse-numbers=false : accepte aussi \SI{2,5 \times 10^{-3}}{...}
\sisetup{locale=FR,output-decimal-marker={,},group-minimum-digits=4,parse-numbers=false}
\DeclareSIUnit\ohm{\ensuremath{\symup{\Omega}}} % U+2126 absent de la police
\DeclareSIUnit\division{div} % réglages d'oscilloscope (ms/div, V/div)
\DeclareSIUnit\tr{tr} % tours (vitesse de rotation, ex. \SI{1500}{\tr\per\minute})
\DeclareSIUnit\molar{M} % concentration molaire (ex. \SI{3}{\molar}, \SI{1}{\micro\molar})
\DeclareSIUnit\an{an} % année (le modèle confond parfois avec \year, un compteur TeX) — ex. \SI{5}{\kilo\meter\per\an}
\DeclareSIUnit\FCFA{FCFA} % franc CFA (ex. \SI{500}{\FCFA\per\kilo\gram})
\DeclareSIUnit\degreeBrix{\textdegree Brix} % teneur en sucre d'un sirop/jus (réfractométrie)
\DeclareSIUnit\month{mois} % le modèle utilise parfois \month, un compteur TeX, comme unité de durée
\usepackage{qrcode}
% \degree : commande fréquemment utilisée par les modèles pour un angle ou une
% température en dehors de \SI{}{} (non fournie par les paquets chargés).
\providecommand{\degree}{^{\circ}}
% \qed (fin de démonstration), utilisé par les modèles sans amsthm
\providecommand{\qed}{\hfill$\square$}
% \barre : double barre des valeurs interdites dans un tableau de variations (tkz-tab)
\providecommand{\barre}{\ensuremath{\|}}
% environnement proof (amsthm non chargé) utilisé par les modèles
\ifdefined\proof\else\newenvironment{proof}[1][Démonstration]{\par\noindent\textit{#1.}\ }{\qed\par}\fi
\usepackage{lastpage}
\usepackage{hyperref}
\hypersetup{hidelinks}

% ============================================================
% Palette « Pop » OnBuch+ — mêmes hex que la classe P (plus_ui.dart)
% ============================================================
\definecolor{poporange}{HTML}{F59321}
\definecolor{popdark}{HTML}{E07A0C}
\definecolor{popcream}{HTML}{FFF6E8}
\definecolor{popink}{HTML}{1C1714}
\definecolor{poppurple}{HTML}{7A5AE0}
\definecolor{popgreen}{HTML}{2E9E6A}
\definecolor{poppink}{HTML}{E0457B}
\definecolor{popblue}{HTML}{2F7DE1}
\definecolor{popgold}{HTML}{C98A12}
\definecolor{popmuted}{HTML}{8A7F76}
\definecolor{popline}{HTML}{EADFD2}
\definecolor{popgray}{HTML}{8A7F76}
\colorlet{popgrey}{popgray}
% Alias tolérants (le modèle nomme parfois les couleurs différemment)
\colorlet{popwhite}{white}
\colorlet{popblack}{popink}
\colorlet{popred}{poppink}
\colorlet{popyellow}{popgold}
% Teintes claires (fonds de tableaux/encadrés) — le modèle nomme parfois
% les couleurs avec un suffixe L (ex. popblueL) sans les avoir définies.
\colorlet{poporangeL}{poporange!20}
\colorlet{popdarkL}{popdark!20}
\colorlet{poppurpleL}{poppurple!20}
\colorlet{popgreenL}{popgreen!20}
\colorlet{poppinkL}{poppink!20}
\colorlet{popblueL}{popblue!20}
\colorlet{popgoldL}{popgold!20}
\colorlet{popredL}{popred!20}
\colorlet{popgrayL}{popgray!20}
% Teintes claires pour fonds de tableaux / zones
\colorlet{popblueL}{popblue!10!white}
\colorlet{poporangeL}{poporange!14!white}
\colorlet{popgreenL}{popgreen!12!white}
\colorlet{poppurpleL}{poppurple!12!white}
\colorlet{poppinkL}{poppink!12!white}
\colorlet{popgoldL}{popgold!14!white}

\pagecolor{popcream}

% ============================================================
% Typographie
% ============================================================
\defaultfontfeatures{Ligatures=TeX}
\setmainfont{PlusJakartaSans-Medium.ttf}[Path=../../../fonts/,BoldFont=PlusJakartaSans-Bold.ttf]
% Symboles, API et emojis absents de la police principale : repli sur DejaVu Sans (caractères actifs)
\newfontface\symfont{DejaVuSans.ttf}[Path=../../../fonts/]
\makeatletter
\newcommand\symactive[1]{\catcode#1=13 \begingroup\lccode`\~=#1 \lowercase{\endgroup\def~}{{\symfont\char#1}}}
\newcommand\symblank[1]{\catcode#1=13 \begingroup\lccode`\~=#1 \lowercase{\endgroup\def~}{}}
\newcommand\symhyph[1]{\catcode#1=13 \begingroup\lccode`\~=#1 \lowercase{\endgroup\def~}{\mbox{-}}}
\newcount\symc
\newcommand\symrange[2]{\symc=#1 \@whilenum\symc<#2 \do{\expandafter\symactive\expandafter{\the\symc}\advance\symc by 1 }}
\newcommand\symblankrange[2]{\symc=#1 \@whilenum\symc<#2 \do{\expandafter\symblank\expandafter{\the\symc}\advance\symc by 1 }}
\makeatother
\symrange{"0250}{"02FF}\symactive{"2713}\symactive{"2714}\symactive{"2717}\symactive{"2718}\symactive{"2610}\symactive{"2611}\symactive{"2612}
\symhyph{"2011}\symhyph{"2E1A}\symblankrange{"1F300}{"1FAFF}\symblank{"FE0F}
% Maths en sans-serif (Fira Math) avec chiffres et lettres droites de Plus
% Jakarta, pour que les formules se fondent dans le texte.
\usepackage[math-style=ISO,bold-style=ISO]{unicode-math}
\setmathfont{FiraMath-Regular.otf}[Path=../../../fonts/]
\setmathfont{PlusJakartaSans-Medium.ttf}[Path=../../../fonts/,range={up/num,up/{Latin,latin},\mathup}]
\usepackage{icomma} % virgule décimale sans espace en mode math
% Caractères mathématiques Unicode absents des polices de texte (Plus Jakarta,
% Archivo Black) : rendus par leur équivalent mathématique, en texte comme en
% formule — sinon ils s'affichent en carré vide (ex. « Arithmétique dans ℤ »).
\usepackage{newunicodechar}
\newunicodechar{ℝ}{\ensuremath{\mathbb{R}}}
\newunicodechar{ℤ}{\ensuremath{\mathbb{Z}}}
\newunicodechar{ℂ}{\ensuremath{\mathbb{C}}}
\newunicodechar{ℕ}{\ensuremath{\mathbb{N}}}
\newunicodechar{ℚ}{\ensuremath{\mathbb{Q}}}
\newunicodechar{𝔻}{\ensuremath{\mathbb{D}}}
\newunicodechar{α}{\ensuremath{\alpha}}
\newunicodechar{β}{\ensuremath{\beta}}
\newunicodechar{γ}{\ensuremath{\gamma}}
\newunicodechar{δ}{\ensuremath{\delta}}
\newunicodechar{ε}{\ensuremath{\varepsilon}}
\newunicodechar{θ}{\ensuremath{\theta}}
\newunicodechar{λ}{\ensuremath{\lambda}}
\newunicodechar{μ}{\ensuremath{\mu}}
\newunicodechar{σ}{\ensuremath{\sigma}}
\newunicodechar{τ}{\ensuremath{\tau}}
\newunicodechar{φ}{\ensuremath{\varphi}}
\newunicodechar{ψ}{\ensuremath{\psi}}
\newunicodechar{ω}{\ensuremath{\omega}}
\newunicodechar{Σ}{\ensuremath{\Sigma}}
% majuscules produites par \MakeUppercase dans les titres (α-aminés -> Α)
\newunicodechar{Α}{\ensuremath{\alpha}}
\newunicodechar{Β}{\ensuremath{\beta}}
\newunicodechar{Γ}{\ensuremath{\Gamma}}
\newunicodechar{Δ}{\ensuremath{\Delta}}
\newunicodechar{Ε}{\ensuremath{\varepsilon}}
\newunicodechar{Θ}{\ensuremath{\Theta}}
\newunicodechar{Λ}{\ensuremath{\Lambda}}
\newunicodechar{Μ}{\ensuremath{\mu}}
\newunicodechar{Τ}{\ensuremath{\tau}}
\newunicodechar{Φ}{\ensuremath{\Phi}}
\newunicodechar{Ψ}{\ensuremath{\Psi}}
\newunicodechar{Ω}{\ensuremath{\Omega}}
\newunicodechar{↦}{\ensuremath{\mapsto}}
\newunicodechar{∈}{\ensuremath{\in}}
\newunicodechar{∉}{\ensuremath{\notin}}
\newunicodechar{∧}{\ensuremath{\wedge}}
\newunicodechar{⇔}{\ensuremath{\Leftrightarrow}}
\newunicodechar{⟺}{\ensuremath{\Longleftrightarrow}}
\newunicodechar{⇒}{\ensuremath{\Rightarrow}}
\newunicodechar{⟹}{\ensuremath{\Longrightarrow}}
\newunicodechar{∘}{\ensuremath{\circ}}
\newunicodechar{∪}{\ensuremath{\cup}}
\newunicodechar{∩}{\ensuremath{\cap}}
\newunicodechar{≡}{\ensuremath{\equiv}}
\newunicodechar{⊂}{\ensuremath{\subset}}
\newunicodechar{⊥}{\ensuremath{\perp}}
\newunicodechar{∃}{\ensuremath{\exists}}
\newunicodechar{∀}{\ensuremath{\forall}}
\newunicodechar{∛}{\ensuremath{\sqrt[3]{\,}}}
\newunicodechar{∜}{\ensuremath{\sqrt[4]{\,}}}
\newunicodechar{⃗}{\ensuremath{{}^{\to}}}
\newunicodechar{ⁿ}{\ifmmode^{n}\else\textsuperscript{\ensuremath{n}}\fi}
\newunicodechar{ˣ}{\ifmmode^{x}\else\textsuperscript{\ensuremath{x}}\fi}
\newunicodechar{ᵇ}{\ifmmode^{b}\else\textsuperscript{\ensuremath{b}}\fi}
\newunicodechar{ʳ}{\ifmmode^{r}\else\textsuperscript{\ensuremath{r}}\fi}
\newunicodechar{ᵘ}{\ifmmode^{u}\else\textsuperscript{\ensuremath{u}}\fi}
\newunicodechar{ʸ}{\ifmmode^{y}\else\textsuperscript{\ensuremath{y}}\fi}
\newunicodechar{ᵃ}{\ifmmode^{a}\else\textsuperscript{\ensuremath{a}}\fi}
\newunicodechar{ᵉ}{\ifmmode^{e}\else\textsuperscript{\ensuremath{e}}\fi}
\newunicodechar{ᵏ}{\ifmmode^{k}\else\textsuperscript{\ensuremath{k}}\fi}
\newunicodechar{ᵖ}{\ifmmode^{p}\else\textsuperscript{\ensuremath{p}}\fi}
\newunicodechar{ᴺ}{\ifmmode^{N}\else\textsuperscript{\ensuremath{N}}\fi}
\newunicodechar{ᶜ}{\ifmmode^{c}\else\textsuperscript{\ensuremath{c}}\fi}
\newunicodechar{ʰ}{\ifmmode^{h}\else\textsuperscript{\ensuremath{h}}\fi}
\newunicodechar{ᵠ}{\ifmmode^{q}\else\textsuperscript{\ensuremath{q}}\fi}
\newunicodechar{ₙ}{\ifmmode_{n}\else\textsubscript{\ensuremath{n}}\fi}
\newunicodechar{ₐ}{\ifmmode_{a}\else\textsubscript{\ensuremath{a}}\fi}
\newunicodechar{ᵢ}{\ifmmode_{i}\else\textsubscript{\ensuremath{i}}\fi}
\newunicodechar{ᵣ}{\ifmmode_{r}\else\textsubscript{\ensuremath{r}}\fi}
\newunicodechar{ₖ}{\ifmmode_{k}\else\textsubscript{\ensuremath{k}}\fi}
\newunicodechar{ᵦ}{\ifmmode_{\beta}\else\textsubscript{\ensuremath{\beta}}\fi}
\newunicodechar{ₓ}{\ifmmode_{x}\else\textsubscript{\ensuremath{x}}\fi}
\newfontfamily\popdisplay{ArchivoBlack-Regular.ttf}[Path=../../../fonts/]
\newfontfamily\poplogo{SpaceGrotesk-Bold.ttf}[Path=../../../fonts/]
\newfontfamily\popsemi{PlusJakartaSans-SemiBold.ttf}[Path=../../../fonts/]
\newfontfamily\popxbold{PlusJakartaSans-ExtraBold.ttf}[Path=../../../fonts/]
\newfontfamily\popxboldls{PlusJakartaSans-ExtraBold.ttf}[Path=../../../fonts/,LetterSpace=3.0]

\setlist[enumerate]{leftmargin=1.7em,itemsep=5pt,topsep=4pt}
\setlist[itemize]{leftmargin=1.7em,itemsep=4pt,topsep=4pt,label={\color{poporange}\textbullet}}
\setlength{\parindent}{0pt}
\setlength{\parskip}{5pt}
\renewcommand{\arraystretch}{1.25}

% pgfplots : style Pop par défaut
\pgfplotsset{
  every axis/.append style={axis line style={popink,line width=1pt},tick style={popink},
    grid style={popline,line width=0.5pt},label style={font=\small},tick label style={font=\footnotesize}},
  popcurve/.style={poporange,line width=1.8pt,no markers,smooth},
}
% Même style disponible dans un \draw TikZ simple (hors pgfplots)
\tikzset{popcurve/.style={poporange,line width=1.8pt}}
\tikzset{popgrid/.style={popline,very thin}}
\colorlet{popcurve}{poporange} % le nom du style est parfois utilisé comme couleur

% ============================================================
% Pill / squiggle
% ============================================================
\newcommand{\poppill}[3]{%
  \tikz[baseline=-0.55ex]\node[draw=popink,line width=1.2pt,rounded corners=2.6pt,%
    fill=#2,inner xsep=6.5pt,inner ysep=3pt]{%
    \popxboldls\fontsize{7}{7}\selectfont\color{#3}\MakeUppercase{#1}};%
}
\newcommand{\popsquiggle}[1]{%
  \tikz[baseline=-0.4ex]{
    \draw[#1,line width=2.6pt,line cap=round]
      plot [smooth] coordinates {(0,0.09) (0.34,0.01) (0.68,0.21) (1.02,0.01) (1.36,0.10)};
  }%
}
% Mot-clé surligné (marqueur orange sous le texte)
\newcommand{\cle}[1]{{\popxbold\color{popdark}#1}}

% ============================================================
% En-tête / pied
% ============================================================
\pagestyle{fancy}
\fancyhf{}
\renewcommand{\headrulewidth}{0pt}
\renewcommand{\footrulewidth}{0pt}
\lhead{\raisebox{-0.15ex}{\poplogo\fontsize{13}{13}\selectfont\color{popink}OnBuch\color{poporange}+}}
\rhead{\poppill{\DOCMATIERE\ \textbullet\ \DOCNIVEAU\ \textbullet\ Leçon \DOCLECON}{white}{popink}}
\lfoot{\popxbold\fontsize{7.6}{7.6}\selectfont\color{popmuted}\MakeUppercase{Cours signé OnBuch+}\ \textbullet\ {\popsemi\DOCTITRE}}
\rfoot{\tikz[baseline=-0.6ex]\node[rounded corners=8.5pt,fill=popink,minimum height=17pt,minimum width=17pt,inner xsep=5pt,inner sep=0]{\popxbold\fontsize{8}{8}\selectfont\color{popcream}\thepage\,/\,\pageref*{LastPage}};}

% ============================================================
% Titres
% ============================================================
\newcommand{\chaptitre}[2]{%
  \begin{tcolorbox}[enhanced,colback=poporange,colframe=popink,boxrule=2.6pt,arc=11pt,
      drop shadow=popink,left=15pt,right=15pt,top=12pt,bottom=11pt,before skip=3pt,after skip=18pt]
    \poppill{#2}{popink}{popcream}\par\vspace{9pt}
    {\raggedright\hyphenpenalty=10000\exhyphenpenalty=10000\popdisplay\fontsize{22}{24}\selectfont\color{popcream}\MakeUppercase{#1}\par}
  \end{tcolorbox}
}
% Section numérotée : pastille orange avec le numéro + titre Archivo
\newcounter{popsec}
\newcommand{\coursec}[1]{%
  \stepcounter{popsec}\setcounter{popsub}{0}%
  \par\vspace{16pt}\noindent
  \begin{minipage}{\linewidth}
  \tikz[baseline=-0.7ex]\node[draw=popink,line width=1.6pt,fill=poporange,rounded corners=4pt,minimum size=22pt,inner sep=0]{\popdisplay\fontsize{12}{12}\selectfont\color{popcream}\thepopsec};%
  \hspace{8pt}{\popdisplay\fontsize{15}{17}\selectfont\color{popink}\MakeUppercase{#1}}\par
  \vspace{4pt}\noindent{\color{poporange}\rule{36pt}{3.6pt}}
  \end{minipage}\par\nopagebreak\vspace{8pt}
}
\newcounter{popsub}
\newcommand{\courssub}[1]{%
  \stepcounter{popsub}%
  \par\vspace{9pt}\noindent{\popxbold\fontsize{11.5}{13}\selectfont\color{popink}{\color{poporange}\thepopsec.\thepopsub}\hspace{6pt}#1}\par\nopagebreak\vspace{4pt}
}

% ============================================================
% Boîtes pédagogiques — même recette : fond blanc, bordure épaisse colorée,
% ombre dure encre, badge plein. Titre optionnel [#1].
% ============================================================
\tcbset{popbase/.style={breakable,enhanced,colback=white,boxrule=2.3pt,arc=9pt,
  shadow={3pt}{-3pt}{0pt}{popink},left=14pt,right=14pt,top=11pt,bottom=11pt,before skip=11pt,after skip=15pt,
  fontupper=\popsemi\fontsize{10.6}{14.5}\selectfont\color{popink},
  fontlower=\popsemi\fontsize{10.6}{14.5}\selectfont\color{popink}}}
% Coche dessinée (le glyphe ✓ n'existe pas dans les polices du document)
\newcommand{\popcheck}[1]{\tikz[baseline=-0.5ex]\draw[#1,line width=1.6pt,line cap=round,line join=round](-0.13,0.0)--(-0.04,-0.09)--(0.13,0.1);}
\providecommand{\coche}{\popcheck{popgreen}}
\providecommand{\resultat}[1]{\par\smallskip\noindent\fcolorbox{popgreen}{popgreenL}{\parbox{\dimexpr\linewidth-2\fboxsep-2\fboxrule\relax}{\textbf{Résultat.} #1}}\par\smallskip}
\newcommand{\popboxhead}[3]{\poppill{#1}{#2}{white}\ifx\relax#3\relax\else\hspace{6pt}{\popxbold\fontsize{10.6}{12}\selectfont\color{popink}#3}\fi\par\vspace{7pt}}

\newtcolorbox{aretenir}[1][]{popbase,colframe=popgreen,before upper={\popboxhead{À retenir}{popgreen}{#1}}}
% Texte en anglais (document de lecture, dialogue, extrait) : cadre
% d'étude. Un titre optionnel (source, auteur) s'affiche dans l'en-tête.
\newtcolorbox{textebox}[1][]{popbase,colframe=popblue,before upper={\popboxhead{Text}{popblue}{#1}\begin{otherlanguage}{english}},after upper={\end{otherlanguage}}}
% Marques ✓ / ✗ (forme correcte / fautive) souvent utilisées par les modèles
\providecommand{\cross}{\textcolor{poppink}{\ensuremath{\times}}}
\providecommand{\xmark}{\cross}\providecommand{\crossmark}{\cross}\providecommand{\cmark}{\ensuremath{\checkmark}}
% Ligne de réponse / justification à compléter (inventée par les modèles)
\providecommand{\justification}[1]{\par\noindent\textit{Justification~:} \dotfill\par}
% \textipa (package tipa non chargé) : la police ne couvre pas l'API -> simple passage
\providecommand{\textipa}[1]{#1}
% Soulignement (ulem non chargé) : \uline, \ul
\providecommand{\uline}[1]{\underline{#1}}\providecommand{\ul}[1]{\underline{#1}}
% \glossaire{\item ...} inventée par les modèles : liste de glossaire
\providecommand{\glossaire}[1]{\begin{itemize}#1\end{itemize}}
% \alert (beamer) inventée par les modèles : texte mis en évidence
\providecommand{\alert}[1]{\textbf{\textcolor{poppink}{#1}}}
% variantes ASCII d'ordinaux écrites en macro
\providecommand{\iere}{\textsuperscript{re}}\providecommand{\ieme}{\textsuperscript{e}}\providecommand{\ere}{\textsuperscript{re}}\providecommand{\eme}{\textsuperscript{e}}\providecommand{\er}{\textsuperscript{er}}\providecommand{\ie}{\textsuperscript{e}}
% Ordinaux français écrits en macro par les modèles : 1\ière, 3\ième
\newcommand{\ière}{\textsuperscript{re}}\newcommand{\ième}{\textsuperscript{e}}
% \exemple (inventée par les modèles) : étiquette « Exemple : »
\providecommand{\exemple}{\textbf{Exemple~:}\ }
% \square utilisable aussi hors mode math (cases à cocher)
\AtBeginDocument{\let\squareMath\square\renewcommand{\square}{\ensuremath{\squareMath}}}
% Mot ou phrase en anglais à l'intérieur d'un paragraphe français
\newcommand{\eng}[1]{\ifmmode\text{\textit{#1}}\else\begingroup\selectlanguage{english}\itshape #1\endgroup\fi}
\let\esp\eng % alias toléré
% flèches d'intonation inventées par les modèles
\providecommand{\textsearrow}{}\renewcommand{\textsearrow}{\ensuremath{\searrow}}\providecommand{\textnearrow}{}\renewcommand{\textnearrow}{\ensuremath{\nearrow}}
% \fr{..} : mot français (inventé par les modèles)
\providecommand{\fr}[1]{\textit{#1}}
\newtcolorbox{exemplebox}[1][]{popbase,colframe=poporange,before upper={\popboxhead{Exemple}{poporange}{#1}}}
\newtcolorbox{definition}[1][]{popbase,colframe=popblue,before upper={\popboxhead{Définition}{popblue}{#1}}}
\newtcolorbox{propriete}[1][]{popbase,colframe=poppurple,before upper={\popboxhead{Propriété}{poppurple}{#1}}}
\newtcolorbox{methode}[1][]{popbase,colframe=popgold,before upper={\popboxhead{Méthode}{popgold}{#1}}}
\newtcolorbox{attention}[1][]{popbase,colframe=poppink,before upper={\popboxhead{Attention}{poppink}{#1}}}
\newtcolorbox{experience}[1][]{popbase,colframe=popblue,colback=popblueL,before upper={\popboxhead{Expérience}{popblue}{#1}}}
% « Le savais-tu ? » : carte violette pleine (contexte, culture, Cameroun)
\newtcolorbox{savaistu}[1][]{popbase,colback=poppurple,colframe=popink,
  fontupper=\popsemi\fontsize{10.4}{14.2}\selectfont\color{white},
  before upper={\poppill{Le savais-tu\,?}{white}{poppurple}\ifx\relax#1\relax\else\hspace{6pt}{\popxbold\color{white}#1}\fi\par\vspace{7pt}}}
% Exercice résolu : énoncé en haut, \tcblower puis la solution sur fond crème
\newcounter{popexo}\newcounter{popexr}
\newtcolorbox{exoresolu}[1][]{popbase,colframe=poporange,colbacklower=popcream,
  segmentation style={popink,line width=1.2pt,dashed},
  before upper={\stepcounter{popexr}\popboxhead{Exercice résolu \thepopexr}{poporange}{#1}},
  before lower={\poppill{Solution}{popink}{popcream}\par\vspace{6pt}}}
% Exercice (non résolu en place) — la correction est dans la partie Corrigés
\newtcolorbox{exercice}[1][]{popbase,colframe=popink,
  before upper={\stepcounter{popexo}\popboxhead{Exercice \thepopexo}{popink}{#1}}}
% Corrigé
\newtcolorbox{corrige}[1][]{popbase,colframe=popgreen,colback=popgreenL,
  before upper={\popboxhead{Corrigé}{popgreen}{#1}}}
% Objectifs / prérequis / bilan
\newtcolorbox{objectifs}{popbase,colframe=popink,colback=poporangeL,
  before upper={\popboxhead{Objectifs de la leçon}{popink}{}}}
\newtcolorbox{prerequis}{popbase,colframe=popink,colback=popblueL,
  before upper={\popboxhead{Prérequis}{popblue}{}}}
\newtcolorbox{bilan}{popbase,colframe=popink,colback=white,boxrule=2.8pt,
  before upper={\popboxhead{Fiche bilan}{poporange}{L'essentiel à connaître pour l'examen}}}
% Figure encadrée avec légende
\newtcolorbox{popfigure}[1][]{enhanced,breakable=false,colback=white,colframe=popline,boxrule=1.4pt,arc=8pt,
  left=10pt,right=10pt,top=10pt,bottom=8pt,before skip=10pt,after skip=12pt,halign=center,
  lower separated=false,halign lower=center,
  fontlower=\popsemi\fontsize{9}{11.5}\selectfont\color{popmuted},
  #1}
\newcommand{\legende}[1]{\par\vspace{6pt}{\popsemi\fontsize{9}{11.5}\selectfont\color{popmuted}#1}}

% ============================================================
% Signature OnBuch+
% ============================================================
\newcommand{\poplogoinline}[1]{{\poplogo\fontsize{#1}{#1}\selectfont\color{popink}OnBuch\color{poporange}+}}
\newcommand{\signatureonbuch}{%
  \begin{tcolorbox}[enhanced,colback=popink,colframe=popink,boxrule=2.2pt,arc=10pt,
      drop shadow=poporange,left=14pt,right=14pt,top=10pt,bottom=10pt,before skip=0pt,after skip=0pt]
    \tikz[baseline=-0.7ex]\node[circle,fill=popgreen,draw=popcream,line width=1.3pt,minimum size=22pt,inner sep=0]{\popcheck{white}};%
    \hspace{9pt}\begin{minipage}[c]{0.62\linewidth}
      {\popxbold\fontsize{12}{13}\selectfont\color{popcream}Signé\ \textbullet\ {\poplogo OnBuch\color{poporange}+}}\par\vspace{2pt}
      {\raggedright\popsemi\fontsize{8}{10.5}\selectfont\color{popline}\MakeUppercase{Équipe pédagogique OnBuch+}\\\MakeUppercase{Conforme au programme officiel MINESEC}\par}
    \end{minipage}\hfill
    {\poplogo\fontsize{22}{22}\selectfont\color{popcream}OnBuch\color{poporange}+}
  \end{tcolorbox}%
}

% ============================================================
% Page de garde
% #1 titre · #2 matière · #3 niveau · #4 module · #5 n° leçon · #6 durée ·
% #7 description / objectifs (texte ou itemize)
% ============================================================
\newcommand{\pagegarde}[7]{%
  \thispagestyle{empty}
  \newgeometry{a4paper,margin=1.7cm,top=1.6cm,bottom=1.4cm}%
  \begin{tikzpicture}[remember picture,overlay]
    \fill[poporange!18] ([xshift=-3.2cm,yshift=-3.6cm]current page.north east) circle (4.6cm);
    \fill[poppurple!14] ([xshift=2.4cm,yshift=7.6cm]current page.south west) circle (3.8cm);
  \end{tikzpicture}%
  {\poplogo\fontsize{30}{30}\selectfont\color{popink}OnBuch\color{poporange}+}\hfill
  \raisebox{0.6ex}{\poppill{Cours signé \textbullet\ Premium}{popink}{popcream}}\par
  \vspace{1pt}\popsquiggle{poppurple}\par
  \vspace{8pt}
  \begin{tcolorbox}[enhanced,colback=poporange,colframe=popink,boxrule=2.8pt,arc=13pt,
      drop shadow=popink,left=18pt,right=18pt,top=12pt,bottom=13pt,after skip=10pt]
    \poppill{#2 \textbullet\ #3}{popink}{popcream}\hspace{5pt}\poppill{#4}{white}{popink}\par\vspace{8pt}
    {\popdisplay\fontsize{14}{14}\selectfont\color{popink}LEÇON #5}\par\vspace{4pt}
    {\raggedright\hyphenpenalty=10000\exhyphenpenalty=10000\popdisplay\fontsize{25}{27}\selectfont\color{popcream}\MakeUppercase{#1}\par}
    \vspace{8pt}
    {\popxbold\fontsize{9.5}{9.5}\selectfont\color{popink}\MakeUppercase{Durée conseillée : #6}}
  \end{tcolorbox}
  \begin{tcolorbox}[enhanced,colback=white,colframe=popink,boxrule=2.2pt,arc=10pt,
      drop shadow=popink,left=16pt,right=16pt,top=10pt,bottom=10pt,after skip=8pt]
    \poppill{Dans cette leçon}{poppurple}{white}\par\vspace{5pt}
    \popsemi\fontsize{9.3}{12.6}\selectfont\color{popink}#7
  \end{tcolorbox}
  \noindent\begin{minipage}[t]{0.487\linewidth}
    \begin{tcolorbox}[enhanced,colback=popgreen,colframe=popink,boxrule=2.2pt,arc=10pt,
        drop shadow=popink,left=11pt,right=11pt,top=10pt,bottom=10pt,height=3.1cm,valign=center]
      \tcbox[colback=white,colframe=popink,boxrule=1.3pt,arc=5pt,boxsep=0pt,nobeforeafter,
        left=3pt,right=3pt,top=3pt,bottom=3pt,valign=center]{\qrcode[height=1.9cm]{https://play.google.com/store/apps/details?id=com.luvvix.onbuch&hl=fr}}%
      \hspace{8pt}\begin{minipage}[c]{0.5\linewidth}
        {\popdisplay\fontsize{11}{12.5}\selectfont\color{white}\MakeUppercase{Télécharge OnBuch}}\par\vspace{4pt}
        {\raggedright\popsemi\fontsize{8.4}{11}\selectfont\color{white}Gratuit \textbullet\ Google Play\\Tuteur IA, annales, concours, résultats\par}
      \end{minipage}
    \end{tcolorbox}
  \end{minipage}\hfill
  \begin{minipage}[t]{0.487\linewidth}
    \begin{tcolorbox}[enhanced,colback=poppurple,colframe=popink,boxrule=2.2pt,arc=10pt,
        drop shadow=popink,left=11pt,right=11pt,top=10pt,bottom=10pt,height=3.1cm,valign=center]
      \tikz[baseline=-0.7ex]\node[circle,fill=white,draw=popink,line width=1.3pt,minimum size=1.6cm,inner sep=0]{\popdisplay\fontsize{24}{24}\selectfont\color{poppurple}+};%
      \hspace{8pt}\begin{minipage}[c]{0.6\linewidth}
        {\popdisplay\fontsize{11}{12.5}\selectfont\color{white}\MakeUppercase{Passe à OnBuch+}}\par\vspace{4pt}
        {\raggedright\popsemi\fontsize{8.4}{11}\selectfont\color{white}Tous les cours signés, Tuteur IA illimité, fiches PDF — onglet OnBuch+ de l'app\par}
      \end{minipage}
    \end{tcolorbox}
  \end{minipage}
  \vspace{8pt}
  \popcontenu
  \vfill
  \signatureonbuch
  \restoregeometry
  \newpage
}

% Rangée « Ce cours contient » (page de garde)
\newcommand{\poptile}[3]{% #1 couleur #2 gros texte #3 légende
  \begin{tcolorbox}[enhanced,colback=#1,colframe=popink,boxrule=2pt,arc=9pt,shadow={2.5pt}{-2.5pt}{0pt}{popink},
      left=6pt,right=6pt,top=9pt,bottom=9pt,halign=center,valign=center,height=1.75cm,width=\linewidth,nobeforeafter]
    {\popdisplay\fontsize{12.5}{13}\selectfont\color{white}\MakeUppercase{#2}}\par\vspace{3pt}
    {\centering\popsemi\fontsize{7.8}{9.5}\selectfont\color{white}#3\par}
  \end{tcolorbox}}
\newcommand{\popcontenu}{%
  \par\noindent{\popxboldls\fontsize{7.5}{7.5}\selectfont\color{popmuted}CE COURS SIGNÉ CONTIENT}\par\vspace{7pt}
  \noindent\begin{minipage}[t]{0.235\linewidth}\poptile{popblue}{Cours}{complet, illustré, pas à pas}\end{minipage}\hfill
  \begin{minipage}[t]{0.235\linewidth}\poptile{poporange}{Méthodes}{et exercices résolus}\end{minipage}\hfill
  \begin{minipage}[t]{0.235\linewidth}\poptile{poppink}{Exercices}{type examen + corrigés}\end{minipage}\hfill
  \begin{minipage}[t]{0.235\linewidth}\poptile{popgreen}{Fiche bilan}{l'essentiel à retenir}\end{minipage}\par}

% Sommaire
\newtcolorbox{sommaire}{enhanced,colback=white,colframe=popink,boxrule=2.2pt,arc=10pt,
  drop shadow=popink,left=18pt,right=18pt,top=13pt,bottom=13pt,before skip=6pt,after skip=20pt,
  before upper={\poppill{Sommaire}{poppurple}{white}\par\vspace{9pt}\popsemi\fontsize{10.6}{14.5}\selectfont\color{popink}}}

% Signature de fin de cours — poussée tout en bas de la dernière page
\newcommand{\findecours}{%
  \par\vfill\nopagebreak
  \begin{center}\popsquiggle{poporange}\end{center}\vspace{4pt}
  \signatureonbuch
}
\AtBeginDocument{\def\llbracket{\mathopen{[\mkern-3.5mu[}}\def\rrbracket{\mathclose{]\mkern-3.5mu]}}} % intervalles d'entiers (glyphe ⟦ absent de la police)
% Glyphes absents de Fira Math : redéfinis à partir de glyphes disponibles
\AtBeginDocument{%
  \def\setminus{\mathbin{\backslash}}\let\smallsetminus\setminus
  \def\vdots{\mathord{\vbox{\baselineskip3.2pt\lineskiplimit0pt\kern4pt\hbox{.}\hbox{.}\hbox{.}}}}%
  \def\ddots{\mathinner{\mkern1mu\raise6.5pt\hbox{.}\mkern2mu\raise3.6pt\hbox{.}\mkern2mu\raise0.7pt\hbox{.}\mkern1mu}}%
  \def\triangle{\mathord{\scalebox{0.9}{$\symup{\Delta}$}}}%
  \def\nabla{\mathord{\raisebox{\depth}{\scalebox{1}[-1]{$\symup{\Delta}$}}}}%
  \def\checkmark{\popcheck{popdark}}%
  \def\star{\mathbin{\ast}}\def\bigstar{\mathord{\ast}}%
  \def\diamond{\mathbin{\scalebox{0.6}{$\Diamond$}}}%
  \def\bigcup{\mathop{\mathchoice{\raisebox{-0.3em}{\scalebox{1.7}{$\cup$}}}{\scalebox{1.25}{$\cup$}}{\cup}{\cup}}\displaylimits}%
  \def\bigcap{\mathop{\mathchoice{\raisebox{-0.3em}{\scalebox{1.7}{$\cap$}}}{\scalebox{1.25}{$\cap$}}{\cap}{\cap}}\displaylimits}%
  \def\lesseqqgtr{\mathrel{\lessgtr}}\def\bigoplus{\mathop{\oplus}\displaylimits}%
}
\providecommand{\ssi}{si et seulement si} % abréviation fréquente
\AtBeginDocument{\providecommand{\wideparen}{\overparen}} % arc de cercle

\begin{document}

\pagegarde{Lesson 25 — Listening: Texts on keeping informed about health problems related to drugs, tobacco, alcohol abuse, and aging}{Anglais}{Tle SH}{Chapter 3 : Environment, Well-being and Health}{EN25}{2 h}{Cette leçon de compréhension orale (listening) entraîne les élèves de Terminale A à comprendre des textes sur les problèmes de santé liés à la drogue, au tabac, à l'alcool et au vieillissement. Elle fournit un script d'écoute original, le lexique thématique, une méthode d'écoute en trois étapes et des questions de type Bac.
\vspace{4pt}
\begin{multicols}{2}\small
\begin{itemize}
\item À la fin de cette leçon, je sais identifier le thème général d'un enregistrement ou d'un texte lu sur la santé (drogue, tabac, alcool, vieillissement).
\item À la fin de cette leçon, je sais repérer les mots-clés, les chiffres, les dates et les noms propres dans un script d'écoute.
\item À la fin de cette leçon, je sais utiliser le lexique anglais des addictions et du vieillissement (drug abuse, smoking, liver, elderly...).
\item À la fin de cette leçon, je sais répondre à des questions de compréhension de types variés : vrai/faux, QCM, réponses courtes.
\item À la fin de cette leçon, je sais distinguer les informations essentielles des détails secondaires dans un texte écouté.
\item À la fin de cette leçon, je sais réemployer le lexique de la santé dans de courtes phrases personnelles.
\end{itemize}\end{multicols}}

\begin{sommaire}
\begin{multicols}{2}
\begin{enumerate}
\item Mise en route et anticipation
\item Script d'écoute et glossaire
\item Compréhension : questions variées
\item Lexique thématique : addictions et vieillissement
\item Méthode d'écoute en 3 étapes et exemple traité
\item Production guidée : message de sensibilisation
\item Activité d'intégration
\item Méthodes et exercices résolus
\item Exercices
\item Corrigés des exercices
\item Fiche bilan
\end{enumerate}
\end{multicols}
\end{sommaire}

\begin{objectifs}
\begin{itemize}
\item À la fin de cette leçon, je sais identifier le thème général d'un enregistrement ou d'un texte lu sur la santé (drogue, tabac, alcool, vieillissement).
\item À la fin de cette leçon, je sais repérer les mots-clés, les chiffres, les dates et les noms propres dans un script d'écoute.
\item À la fin de cette leçon, je sais utiliser le lexique anglais des addictions et du vieillissement (drug abuse, smoking, liver, elderly...).
\item À la fin de cette leçon, je sais répondre à des questions de compréhension de types variés : vrai/faux, QCM, réponses courtes.
\item À la fin de cette leçon, je sais distinguer les informations essentielles des détails secondaires dans un texte écouté.
\item À la fin de cette leçon, je sais réemployer le lexique de la santé dans de courtes phrases personnelles.
\item À la fin de cette leçon, je sais appliquer une méthode d'écoute en trois étapes (avant, pendant, après l'écoute).
\item À la fin de cette leçon, je sais produire un court message de sensibilisation (awareness message) sur les dangers des addictions.
\end{itemize}
\end{objectifs}

\begin{prerequis}
\begin{itemize}
\item Le vocabulaire de base du corps humain et des maladies (classe de Première : body, disease, pain, treatment).
\item La compréhension de textes lus à voix haute et les questions WH- (who, what, where, when, why, how).
\item Les nombres, pourcentages et dates en anglais (fifty per cent, two thousand, in 2023).
\item Les modaux de conseil et d'obligation (should, must, have to) utiles pour les messages de prévention.
\item La distinction fait/opinion dans un texte (leçon de début d'année sur les connecteurs).
\end{itemize}
\end{prerequis}

\newpage

\coursec{Mise en route et anticipation}

\courssub{Situation d'accroche : dans un taxi à Yaoundé}

Imaginez la scène : vous êtes assis à l'arrière d'un taxi jaune, coincé dans les embouteillages de l'avenue Kennedy à Yaoundé. Le chauffeur monte le son de la radio. La CRTV diffuse un reportage en anglais sur la santé des jeunes Camerounais : \eng{“Alcohol and tobacco consumption among young people is rising sharply in our big cities...”} Vous entendez des chiffres, des noms de maladies, des témoignages de médecins et de lycéens... mais tout va très vite ! À la fin du reportage, un ami vous demande : \eng{“So, what did they say exactly?”} — et vous réalisez que vous n'avez retenu que quelques mots isolés.

C'est exactement la compétence exigée au Baccalauréat : comprendre un texte lu ou un enregistrement sur un sujet de société, et en retenir l'essentiel. La problématique de cette leçon est donc la suivante : \textbf{comment capter le thème, les mots-clés, les chiffres et les noms propres d'un texte oral sur les problèmes de santé liés à la drogue, au tabac, à l'alcool et au vieillissement ?} Pour y répondre, nous allons d'abord activer ce que vous savez déjà, apprendre à prédire le contenu avant l'écoute, et découvrir cinq mots-clés indispensables.

\courssub{Brainstorming : What health problems do you know?}

\begin{experience}[Brainstorming collectif — 5 minutes]
Le professeur écrit au tableau la question : \eng{What health problems do you know?} Chaque élève cite une maladie ou un problème de santé, en anglais si possible. Réponses attendues :
\begin{itemize}
\item \eng{malaria} « le paludisme » — \eng{AIDS} « le sida » — \eng{cancer} « le cancer »
\item \eng{diabetes} « le diabète » — \eng{cholera} « le choléra » — \eng{typhoid} « la typhoïde »
\item \eng{high blood pressure} « l'hypertension » — \eng{obesity} « l'obésité »
\end{itemize}
Le professeur oriente ensuite la discussion : \eng{Today, we will focus on health problems caused by drugs, tobacco, alcohol abuse, and on the health of elderly people.}
\end{experience}

Ce chapitre s'intitule \eng{Environment, Well-being and Health} « Environnement, bien-être et santé ». L'idée directrice : notre santé dépend de notre environnement (air, eau, nourriture) mais aussi de nos \cle{comportements} — fumer, boire, se droguer — et du vieillissement naturel de notre corps.

\courssub{Discussion : drugs, tobacco and alcohol among young Cameroonians}

Avant d'écouter, parlons de ce que nous voyons autour de nous. Répondez en anglais, avec des phrases courtes :

\begin{itemize}
\item \eng{Do young people smoke in your neighbourhood? Where do they buy cigarettes?}
\item \eng{Why do some students drink alcohol or take drugs? Is it peer pressure?} « la pression du groupe »
\item \eng{What can schools and families do to protect young people?}
\end{itemize}

\begin{exemplebox}[Phrases utiles pour la discussion]
\eng{Many young people start smoking because of peer pressure.} « Beaucoup de jeunes commencent à fumer à cause de la pression du groupe. »\par
\eng{Cigarettes are sold everywhere, even near schools.} « On vend des cigarettes partout, même près des écoles. »\par
\eng{Drugs destroy the brain and the future of young people.} « La drogue détruit le cerveau et l'avenir des jeunes. »\par
\eng{Schools should organise awareness campaigns.} « Les écoles devraient organiser des campagnes de sensibilisation. »
\end{exemplebox}

\begin{savaistu}[Tobacco kills — WHO figures]
Selon l'Organisation mondiale de la santé (\eng{WHO — World Health Organization}), le tabac tue plus de 8 millions de personnes par an dans le monde, dont plus d'un million par le \cle{tabagisme passif} (\eng{second-hand smoke}), c'est-à-dire la fumée respirée involontairement par les non-fumeurs. C'est pourquoi de nombreux pays interdisent de fumer dans les lieux publics (\eng{smoking is banned in public places}).
\end{savaistu}

\courssub{Prédiction à partir du titre : What will the recording be about?}

\begin{methode}[Anticiper avant d'écouter]
Avant d'écouter un enregistrement, ne restez jamais passif. Suivez ces étapes :
\begin{enumerate}
\item \textbf{Lisez le titre} : il annonce presque toujours le thème général.
\item \textbf{Lisez toutes les questions} (vrai/faux, QCM) : elles révèlent les détails à repérer — chiffres, noms propres, lieux, causes.
\item \textbf{Observez l'image} éventuelle : qui parle ? où ? de quoi ?
\item \textbf{Formulez une prédiction} : \eng{I think the recording will be about...}
\item \textbf{Préparez mentalement le vocabulaire} du thème : quels mots vais-je probablement entendre ?
\end{enumerate}
\end{methode}

Appliquons tout de suite. Le titre de notre enregistrement est : \eng{“Keeping informed: drugs, tobacco, alcohol and aging”}.

\begin{exemplebox}[Prédictions possibles]
\eng{I think the recording will be about the dangers of smoking and drinking.}\par
\eng{I think a doctor will give advice to young people.}\par
\eng{I think we will hear statistics about diseases and elderly people.}
\end{exemplebox}

\begin{attention}[Le piège n°1 en listening]
Ne cherchez \textbf{pas} à comprendre chaque mot ! C'est l'erreur classique des élèves francophones : dès qu'un mot est inconnu, ils bloquent, paniquent et perdent la suite. En compréhension orale, visez d'abord le \cle{sens global} (\eng{gist}), puis seulement les détails demandés par les questions. Un mot inconnu est rarement essentiel : le contexte suffit souvent à le deviner.
\end{attention}

\courssub{Objectifs d'écoute : que faut-il repérer ?}

Pendant l'écoute, votre cerveau doit fonctionner comme un filtre. Voici les quatre cibles prioritaires :

\begin{center}
\begin{tabularx}{\linewidth}{|l|X|X|}
\hline
\rowcolor{popblueL} \textbf{Cible} & \textbf{English} & \textbf{Exemple à repérer} \\
\hline
Thème & topic / subject & \eng{health problems of young people} \\
\hline
Mots-clés & key words & \eng{smoking, liver, addiction, elderly} \\
\hline
Chiffres & figures / statistics & \eng{8 million deaths, 20\% of students} \\
\hline
Noms propres & proper nouns & \eng{WHO, CRTV, Yaoundé, Dr. Ngo Bell} \\
\hline
\end{tabularx}
\end{center}

\begin{aretenir}[Les mots-outils des questions]
Les questions de compréhension commencent presque toujours par : \eng{Who} (qui), \eng{What} (quoi/quel), \eng{Where} (où), \eng{When} (quand), \eng{Why} (pourquoi), \eng{How many / How much} (combien). Soulignez ces mots dans les questions avant l'écoute : ils vous disent \emph{quel type d'information} écouter.
\end{aretenir}

\courssub{Carte mentale du thème}

\begin{popfigure}
\begin{center}
\begin{tikzpicture}[mindmap, grow cyclic,
  every node/.style={concept, align=center, text=white},
  root concept/.append style={concept color=popblue, font=\bfseries\large},
  level 1/.append style={level distance=3.4cm, sibling angle=90},
  level 1 concept/.append style={font=\bfseries}]
\node[root concept, align=center]{HEALTH\\PROBLEMS}
  child[concept color=poporange]{ node[align=center]{DRUGS\\ \footnotesize cannabis, cocaine} }
  child[concept color=poppurple]{ node[align=center]{TOBACCO\\ \footnotesize cigarettes, smoke} }
  child[concept color=popgreen]{ node[align=center]{ALCOHOL\\ \footnotesize beer, wine, whisky} }
  child[concept color=poppink]{ node[align=center]{AGING\\ \footnotesize elderly people} };
\end{tikzpicture}
\end{center}
\legende{Carte mentale du thème : les quatre grandes causes de problèmes de santé étudiées dans cette leçon.}
\end{popfigure}

\courssub{Pré-enseignement des 5 mots-clés indispensables}

Avant toute écoute, le professeur « pré-enseigne » les mots sans lesquels le texte serait incompréhensible. Observez d'abord ces cinq mots dans des phrases authentiques :

\begin{textebox}[Five key words in context — read before listening]
\eng{Drug abuse is a serious problem in many big cities. When a person cannot stop taking a substance, doctors speak of addiction. Alcohol, for example, slowly destroys the liver, and the damage is often invisible for years. Elderly people are also at risk: their bodies are weaker, and even small quantities of alcohol or tobacco can be dangerous for them. That is why health workers organise awareness campaigns in schools and markets, to inform the population before it is too late.}
\end{textebox}

\begin{definition}[Les 5 mots-clés de la leçon]
\begin{itemize}
\item \cle{abuse} (nom, « a-biouze ») : \emph{usage excessif et dangereux} d'une substance. \eng{drug abuse} « l'abus de drogues », \eng{alcohol abuse} « l'abus d'alcool ». Attention : ne signifie pas « abuser de quelqu'un » au sens de tromper.
\item \cle{addiction} (nom, « a-dik-cheune ») : \emph{dépendance}, impossibilité d'arrêter. Adjectif : \eng{addicted to} « dépendant de ». \eng{He is addicted to tobacco.} « Il est dépendant du tabac. »
\item \cle{liver} (nom, « li-veur ») : le \emph{foie}, l'organe que l'alcool détruit. \eng{Alcohol damages the liver.} « L'alcool endommage le foie. »
\item \cle{elderly} (adjectif, « èl-deur-li ») : \emph{âgé}. \eng{elderly people} ou \eng{the elderly} « les personnes âgées ». Plus poli que \eng{old people}.
\item \cle{awareness} (nom, « a-wère-nèsse ») : la \emph{sensibilisation}, la conscience d'un problème. \eng{awareness campaign} « campagne de sensibilisation ». Verbe : \eng{to raise awareness} « sensibiliser ».
\end{itemize}
\end{definition}

\begin{center}
\begin{tabularx}{\linewidth}{|l|l|X|}
\hline
\rowcolor{popblueL} \textbf{English} & \textbf{Français} & \textbf{Exemple en contexte} \\
\hline
abuse & abus, usage excessif & \eng{Drug abuse is rising among students.} \\
\hline
addiction & dépendance & \eng{Smoking leads to addiction.} \\
\hline
liver & foie & \eng{Too much alcohol destroys the liver.} \\
\hline
elderly (people) & (personnes) âgé(es) & \eng{The elderly need special care.} \\
\hline
awareness & sensibilisation & \eng{We must raise awareness about AIDS.} \\
\hline
\end{tabularx}
\end{center}

\begin{attention}[Faux amis et pièges de prononciation]
\begin{itemize}
\item \eng{abuse} (nom) se prononce « a-biousse », mais le \emph{verbe} \eng{to abuse} se prononce « a-biouze ». Même orthographe, deux prononciations !
\item \eng{the liver} = le foie ; ne confondez pas avec \eng{to live} « vivre » ni avec \eng{delivery} « livraison ».
\item On dit \eng{addicted \textbf{to} drugs} (préposition \eng{to}), jamais \eng{addicted of drugs}.
\item \eng{elderly} est un adjectif : on ne dit pas \eng{an elderly}, mais \eng{an elderly person} ou \eng{the elderly} « les personnes âgées ».
\end{itemize}
\end{attention}

\begin{exoresolu}[Vérifiez vos 5 mots-clés]
Énoncé — Complétez chaque phrase avec le mot-clé qui convient : \eng{abuse, addiction, liver, elderly, awareness}.
\begin{enumerate}
\item \eng{The government has launched an \_\_\_\_\_\_\_\_ campaign against smoking in schools.}
\item \eng{My grandfather is 80 years old; \_\_\_\_\_\_\_\_ people like him must not drink too much.}
\item \eng{Alcohol \_\_\_\_\_\_\_\_ is the cause of many road accidents in Douala.}
\item \eng{The doctor said that whisky had seriously damaged his \_\_\_\_\_\_\_\_.}
\item \eng{After two years of smoking, the boy developed a real \_\_\_\_\_\_\_\_ to cigarettes.}
\end{enumerate}
\tcblower
Solution détaillée :
\begin{enumerate}
\item \eng{awareness} — une campagne de \emph{sensibilisation} : \eng{an awareness campaign}.
\item \eng{elderly} — \eng{elderly people} « les personnes âgées » (grand-père de 80 ans).
\item \eng{abuse} — \eng{alcohol abuse} « l'abus d'alcool » (nom composé fréquent).
\item \eng{liver} — le whisky attaque le \emph{foie}.
\item \eng{addiction} — \eng{an addiction to cigarettes} « une dépendance aux cigarettes » (préposition \eng{to}).
\end{enumerate}
\end{exoresolu}

Vous êtes maintenant prêts : le thème est activé, la prédiction est faite, les cinq mots-clés sont connus. Passons à l'écoute proprement dite.

\coursec{Mise en route, anticipation et script d'écoute}

\courssub{Mise en route et anticipation}

\begin{experience}[Activation des connaissances et prédictions]
Avant d'écouter, observez le titre de l'émission \eng{“Health Watch”} et les images projetées (hôpital, jeunes qui fument, personne âgée, médicaments/tramadol). En binôme, répondez oralement en anglais :
\begin{enumerate}
\item What do you think \eng{“Health Watch”} is about? (Topic)
\item What health problems do you expect to hear about? (Vocabulary prediction: \eng{drugs, tobacco, alcohol, aging, cancer, liver, addiction})
\item Who might be speaking? (Roles: \eng{reporter, doctor, presenter})
\item What is the purpose of such a programme? (Function: \eng{warn, inform, advise, prevent})
\end{enumerate}
Notez 3 mots-clés anglais que vous vous attendez à entendre. Partagez au tableau. Cette phase d'\textbf{anticipation} (prédiction) est la clé de la réussite en compréhension de l'oral au Baccalauréat.
\end{experience}

\courssub{Première écoute : compréhension globale (gist)}

\begin{methode}[Stratégie : Écoute globale (Gist Listening)]
L'objectif n'est \textbf{pas} de comprendre chaque mot, mais de saisir le \textbf{sujet principal}, le \textbf{ton} (alarmiste, informatif, conseillé), la \textbf{structure} (reportage, interview, conseils) et le \textbf{public visé}.
\begin{enumerate}
\item \textbf{Ne prenez pas de notes détaillées.} Notez seulement 2 ou 3 mots-clés qui vous frappent.
\item \textbf{Repérez :} Qui parle ? (Un speaker, un journaliste, un médecin). Où ? (Yaoundé Central Hospital). De quoi ? (Dangers drogues/tabac/alcool, vieillissement).
\item \textbf{Question guide :} \eng{What is the main message of this radio programme?}
\end{enumerate}
\end{methode}

\noindent L'enseignant lit (ou diffuse) le script suivant \textbf{une première fois} à vitesse normale. Écoutez sans écrire.

\begin{textebox}[Script d'écoute — Health Watch, CRTV Radio]
Good morning, listeners. This is Health Watch on CRTV. Today, our reporter takes us to Yaoundé Central Hospital, where Doctor Ngono warns young Cameroonians about the dangers of drug, tobacco and alcohol abuse. According to recent figures, nearly thirty per cent of young adults in urban areas smoke regularly, and alcohol consumption among teenagers has doubled in ten years. Doctor Ngono explains that tobacco damages the lungs and causes cancer, while alcohol destroys the liver. Hard drugs such as cannabis and tramadol lead to addiction, madness and sometimes death. The doctor also reminds us that elderly people need special care: ageing weakens the body and the memory. Her advice is simple: say no to drugs, do sport, eat well, and keep informed about your health. Stay healthy, stay informed!
\end{textebox}

\begin{exoresolu}[Vérification de la compréhension globale]
\textbf{Énoncé :} Après cette première écoute, cochez la bonne réponse pour résumer le reportage.
\begin{enumerate}
\item Le reportage parle principalement de : 
  \begin{itemize}
  \item[a)] L'ouverture d'un nouvel hôpital à Yaoundé.
  \item[b)] Les dangers des addictions et l'importance de la prévention santé.
  \item[c)] L'histoire de la radio CRTV.
  \end{itemize}
\item Le ton du message est :
  \begin{itemize}
  \item[a)] Humoristique et léger.
  \item[b)] Informatif, préventif et solennel.
  \item[c)] Politique et polémique.
  \end{itemize}
\end{enumerate}
\tcblower
\textbf{Solution détaillée :}
\begin{enumerate}
\item \textbf{b)} Le texte cite \eng{dangers}, \eng{warns}, \eng{advice}, \eng{prevention} (implicite par \eng{keep informed}). Les statistiques (30\%, doubled) et les conséquences médicales (cancer, liver destruction, addiction) confirment le thème préventif.
\item \textbf{b)} Le vocabulaire (\eng{warns}, \eng{reminds}, \eng{advice}, \eng{Stay healthy}) et la structure (constat $\rightarrow$ explication $\rightarrow$ conseil) indiquent un ton solennel et préventif, typique d'une campagne de santé publique.
\end{enumerate}
\end{exoresolu}

\courssub{Deuxième écoute : prise de notes ciblée (détails chiffrés et lexicaux)}

\begin{methode}[Stratégie : Écoute sélective (Scanning Listening)]
Lors de la \textbf{deuxième lecture} (ou diffusion), concentrez-vous sur les \textbf{informations précises} : chiffres, noms propres, noms de maladies, noms de substances. Utilisez un tableau de prise de notes comme ci-dessous.
\end{methode}

\begin{center}
\begin{tabularx}{\linewidth}{|X|X|X|}
\hline
\rowcolor{popblueL} \textbf{Catégorie} & \textbf{Informations à relever (Anglais)} & \textbf{Traduction / Notes (Français)} \\
\hline
\textbf{Source / Lieu} & CRTV, Yaoundé Central Hospital, Dr Ngono & Radio nationale, hôpital central, médecin expert \\
\hline
\textbf{Chiffres / Statistiques} & nearly 30\%, young adults, urban areas, doubled in 10 years & ~30\% jeunes adultes urbains fument ; conso. alcool $\times$2 en 10 ans \\
\hline
\textbf{Substances citées} & tobacco, alcohol, cannabis, tramadol & Tabac, alcool, cannabis, tramadol (antalgique opioïde détourné) \\
\hline
\textbf{Conséquences santé} & damages lungs, causes cancer, destroys liver, addiction, madness, death & Abîme poumons, cancer, détruit foie, dépendance, troubles psychiques, mort \\
\hline
\textbf{Population vulnérable} & elderly people, ageing weakens body and memory & Personnes âgées, vieillissement affaiblit corps et mémoire \\
\hline
\textbf{Conseils (Advice)} & say no to drugs, do sport, eat well, keep informed & Refuser drogues, sport, bien manger, s'informer \\
\hline
\end{tabularx}
\end{center}

\noindent Après la deuxième écoute, l'enseignant laisse 2 minutes pour compléter le tableau, puis lance une \textbf{mise en commun orale} : un élève vient écrire au tableau les chiffres, un autre les substances, etc.

\courssub{Vérification collective des prédictions}

Revenez aux hypothèses notées lors de la \textbf{Mise en route} (Section 1).
\begin{itemize}
\item Quelles prédictions étaient \textbf{justes} ? (Ex: “Il sera question de tramadol”, “Il y aura des statistiques”).
\item Quelles prédictions étaient \textbf{fausses} ou \textbf{non mentionnées} ? (Ex: “On parlera du sida”, “Le ministre interviendra”).
\item Quels \textbf{indices sonores} (mots-clés entendus, intonation) ont permis de valider ou d'invalider vos hypothèses ?
\end{itemize}
Cette métacognition renforce vos stratégies d'écoute pour le Baccalauréat.

\courssub{Glossaire des mots difficiles du script}

Voici le lexique essentiel extrait du script, classé par catégorie grammaticale, avec la prononciation approximative (notation francisée) et la traduction contextuelle. \textbf{Apprenez ces mots par cœur} : ils reviennent dans les sujets d'examen.

\begin{definition}[\cle{Abuse} (nom) — \eng{/əˈbjuːs/} « a-biousse »]
\textbf{Sens :} Utilisation excessive, mauvaise utilisation, abus (d'une substance, d'un pouvoir, de la confiance).
\textbf{Collocations fréquentes :} \eng{drug abuse} (toxicomanie / abus de drogues), \eng{alcohol abuse} (abus d'alcool / alcoolisme), \eng{substance abuse} (toxicomanie), \eng{child abuse} (maltraitance d'enfants).
\textbf{Attention :} Le verbe \eng{to abuse} (se conjugue) signifie « abuser de », « maltraiter », « insulter ».
\eng{He abused his authority.} = Il a abusé de son autorité.
\end{definition}

\begin{definition}[\cle{Addiction} (nom) — \eng{/əˈdɪkʃn/} « a-dic-chone »]
\textbf{Sens :} Dépendance (physique ou psychologique) à une substance ou une activité. État de ne plus pouvoir s'en passer.
\textbf{Structure :} \eng{addiction to} [+ nom / gérondif].
\eng{His addiction to tramadol ruined his life.} = Sa dépendance au tramadol a ruiné sa vie.
\textbf{Adjectif :} \eng{addicted} (dépendant). \eng{to be addicted to} = être accro à / être dépendant de.
\eng{He is addicted to social media.} = Il est accro aux réseaux sociaux.
\textbf{Nom de la personne :} \eng{an addict} (un toxicomane, un accro).
\end{definition}

\begin{center}
\begin{tabularx}{\linewidth}{|l|l|X|}
\hline
\rowcolor{popblueL} \textbf{Mot (English)} & \textbf{Nature} & \textbf{Traduction / Contexte (Français)} \\
\hline
\eng{listeners} & n. pl. & auditeurs \\
\eng{reporter} & n. & reporter, journaliste \\
\eng{warns} & v. (3e pers. sg) & met en garde (de \eng{to warn}) \\
\eng{dangers} & n. pl. & dangers \\
\eng{figures} & n. pl. & chiffres, statistiques \\
\eng{young adults} & n. pl. & jeunes adultes \\
\eng{urban areas} & n. pl. & zones urbaines \\
\eng{regularly} & adv. & régulièrement \\
\eng{consumption} & n. & consommation \\
\eng{teenagers} & n. pl. & adolescents \\
\eng{doubled} & v. (past/part.) & a doublé (de \eng{to double}) \\
\eng{damages} & v. (3e pers. sg) & abîme, endommage (de \eng{to damage}) \\
\eng{lungs} & n. pl. & poumons \\
\eng{causes} & v. (3e pers. sg) & cause, provoque \\
\eng{cancer} & n. & cancer \\
\eng{destroys} & v. (3e pers. sg) & détruit (de \eng{to destroy}) \\
\eng{liver} & n. & foie \\
\eng{hard drugs} & n. pl. & drogues dures \\
\eng{cannabis} & n. & cannabis \\
\eng{tramadol} & n. & tramadol (antalgique opioïde détourné) \\
\eng{lead to} & v. (phrase) & mènent à, entraînent (de \eng{to lead to}) \\
\eng{madness} & n. & folie, démence \\
\eng{elderly people} & n. pl. & personnes âgées, aînés \\
\eng{special care} & n. & soins particuliers / adaptés \\
\eng{ageing} & n. / adj. & vieillissement / vieillissant \\
\eng{weakens} & v. (3e pers. sg) & affaiblit (de \eng{to weaken}) \\
\eng{advice} & n. (invariable) & conseil(s) — \textbf{jamais \eng{advices}} \\
\eng{keep informed} & loc. v. & se tenir informé \\
\hline
\end{tabularx}
\end{center}

\courssub{Relecture silencieuse et consolidation lexicale}

Lisez maintenant le script en silence (dans la \texttt{textebox} ci-dessus) pendant 3 minutes. Soulignez les mots du glossaire, notez la structure des phrases (\eng{Subject + Verb + Object}), et observez les connecteurs logiques (\eng{while}, \eng{also}, \eng{according to}).

\begin{exemplebox}[Exemples tirés du script — Traduction et analyse]
\begin{itemize}
\item \eng{According to recent figures, nearly 30 per cent of young adults in urban areas smoke regularly.}
      \newline « Selon des chiffres récents, près de 30 pour cent des jeunes adultes en zones urbaines fument régulièrement. »
      \newline \textbf{Structure :} \eng{According to} (selon) + nom ; \eng{nearly} (presque) + chiffre + \eng{per cent} (invariable, s'écrit en deux mots en anglais britannique).
\item \eng{Alcohol consumption among teenagers has doubled in ten years.}
      \newline « La consommation d'alcool chez les adolescents a doublé en dix ans. »
      \textbf{Grammaire :} \eng{Present Perfect} (\eng{has doubled}) pour un fait passé avec conséquence présente / évolution sur une période non terminée.
\item \eng{Tobacco damages the lungs and causes cancer, while alcohol destroys the liver.}
      \newline « Le tabac abîme les poumons et cause le cancer, \textbf{tandis que} l'alcool détruit le foie. »
      \textbf{Connecteur :} \eng{while} = « alors que / tandis que » (contraste / opposition).
\item \eng{Hard drugs such as cannabis and tramadol lead to addiction, madness and sometimes death.}
      \newline « Les drogues dures \textbf{telles que} le cannabis et le tramadol mènent à la dépendance, la folie et parfois la mort. »
      \textbf{Structure :} \eng{such as} = « tels que / comme » (exemples). Sujet pluriel \eng{drugs} $\rightarrow$ verbe \eng{lead} (sans \eng{s}).
\item \eng{Her advice is simple: say no to drugs, do sport, eat well, and keep informed about your health.}
      \newline « Son conseil est simple : dites non à la drogue, faites du sport, mangez bien, et tenez-vous informés sur votre santé. »
      \textbf{Grammaire :} \eng{Advice} (indénombrable) $\rightarrow$ \eng{is} (singulier). \eng{Imperative} (impératif) pour les conseils.
\end{itemize}
\end{exemplebox}

\begin{attention}[Pièges majeurs pour francophones — Faux-amis et erreurs fréquentes]
\begin{enumerate}
\item \textbf{\eng{Sensible} $\neq$ Sensible (FR).}
      \begin{itemize}
      \item \eng{Sensible} (AN) = \textbf{Raisonnable, sensé, prudent}. \eng{Be sensible!} = Sois raisonnable !
      \item \textbf{Sensible (FR)} = Qui réagit fortement, délicat $\rightarrow$ \eng{Sensitive}. \eng{My skin is sensitive.} = Ma peau est sensible.
      \end{itemize}
\item \textbf{\eng{Advice} (Conseil) est \textbf{indénombrable}.}
      \begin{itemize}
      \item \eng{A piece of advice} (un conseil) — \textbf{Jamais} \eng{an advice} ni \eng{advices}.
      \item \eng{The doctor gave me good advice.} (Pas \eng{good advices}).
      \end{itemize}
\item \textbf{\eng{Drug} vs \eng{Medicine / Medication}.}
      \begin{itemize}
      \item \eng{Drug} : connotation souvent illicite (drogue) ou substance active (médicament en pharma).
      \item \eng{Medicine / Medication} : médicament (remède prescrit). \eng{Take your medicine.} = Prends ton médicament.
      \end{itemize}
\item \textbf{\eng{Elderly} vs \eng{Old}.}
      \begin{itemize}
      \item \eng{Elderly people} / \eng{The elderly} : terme poli, respectueux, administratif/médical (personnes âgées).
      \item \eng{Old people} : neutre ou familier, peut manquer de respect selon le ton.
      \end{itemize}
\item \textbf{\eng{Addiction} : construction.}
      \begin{itemize}
      \item \eng{Addiction to} + nom/gérondif. \eng{Addiction to smoking}.
      \item Préférez souvent la structure verbale : \eng{to be addicted to} + nom/gérondif. \eng{He is addicted to tramadol.}
      \end{itemize}
\item \textbf{\eng{Ageing} (UK) vs \eng{Aging} (US).} Au Cameroun (anglais britannique) : \eng{ageing} (garde le \eng{e} muet avant \eng{-ing}).
\end{enumerate}
\end{attention}

\courssub{Organigramme de la structure du script}

Le schéma ci-dessous résume l'architecture du reportage radio. Mémorisez cet enchaînement logique (Introduction $\rightarrow$ Constat chiffré $\rightarrow$ Expertise médicale $\rightarrow$ Conseil $\rightarrow$ Conclusion) : c'est la structure type d'un \eng{public health announcement} (message de santé publique) que vous devrez imiter en production orale/écrite.

\begin{popfigure}
\begin{tikzpicture}[
  node distance=1.2cm and 1.5cm,
  box/.style={rectangle, rounded corners, draw=popblue, thick, fill=popblueL, text width=3.2cm, align=center, minimum height=1.8cm, font=\small},
  arrow/.style={-Stealth, thick, popdark}
]
% Nodes
\node[box, align=center] (intro){\textbf{Introduction}\\\eng{Good morning... This is Health Watch on CRTV.}\\Hook \eng{:} \eng{Reporter takes us to Yaoundé Central Hospital.}};
\node[box, below=of intro, align=center] (stats){\textbf{Statistics / Constat}\\\eng{Recent figures: 30\% smoke.}\\ \eng{Alcohol consumption doubled.}\\ \eng{Target:} young adults, teenagers.};
\node[box, below=of stats, align=center] (interview){\textbf{Doctor's Interview}\\\eng{Dr Ngono explains:}\\ \eng{-- Tobacco $\rightarrow$ lungs, cancer.}\\ \eng{-- Alcohol $\rightarrow$ liver.}\\ \eng{-- Hard drugs $\rightarrow$ addiction, madness, death.}};
\node[box, below=of interview, align=center] (advice){\textbf{Advice / Prévention}\\\eng{Elderly need special care.}\\ \eng{Ageing weakens body/memory.}\\ \eng{Advice:} No drugs, sport, eat well, stay informed.};
\node[box, below=of advice, align=center] (concl){\textbf{Conclusion}\\\eng{Stay healthy, stay informed!}\\ Slogan final, signature émission.};

% Arrows
\draw[arrow] (intro.south) -- (stats.north) node[midway, right, font=\tiny, text=poppurple] {Transition};
\draw[arrow] (stats.south) -- (interview.north) node[midway, right, font=\tiny, text=poppurple] {Expertise};
\draw[arrow] (interview.south) -- (advice.north) node[midway, right, font=\tiny, text=poppurple] {Extension};
\draw[arrow] (advice.south) -- (concl.north) node[midway, right, font=\tiny, text=poppurple] {Slogan};
\end{tikzpicture}
\legende{Structure rhétorique du script \eng{Health Watch} — Modèle pour votre production (Section 6).}
\end{popfigure}

\courssub{Entraînement : Reconstitution du script (Gap-fill)}

\begin{exercice}[Reconstitution du script]
Complétez le résumé ci-dessous avec les mots du glossaire (forme correcte). La première écoute et la relecture vous y aideront.
\begin{enumerate}
\item The radio programme is called \_\_\_\_\_\_\_\_\_\_ Watch on CRTV.
\item Dr Ngono \_\_\_\_\_\_\_\_\_\_ young Cameroonians about the \_\_\_\_\_\_\_\_\_\_ of drug, tobacco and alcohol \_\_\_\_\_\_\_\_\_\_.
\item Statistics show that alcohol \_\_\_\_\_\_\_\_\_\_ among \_\_\_\_\_\_\_\_\_\_ has \_\_\_\_\_\_\_\_\_\_ in ten years.
\item Tobacco \_\_\_\_\_\_\_\_\_\_ the \_\_\_\_\_\_\_\_\_\_ and \_\_\_\_\_\_\_\_\_\_ cancer.
\item \_\_\_\_\_\_\_\_\_\_ drugs such as cannabis and tramadol \_\_\_\_\_\_\_\_\_\_ to \_\_\_\_\_\_\_\_\_\_.
\item \_\_\_\_\_\_\_\_\_\_ people need \_\_\_\_\_\_\_\_\_\_ care because \_\_\_\_\_\_\_\_\_\_ weakens the body.
\item Final \_\_\_\_\_\_\_\_\_\_: \eng{Say no to drugs, do sport, eat well, keep informed.}
\end{enumerate}
\end{exercice}

\begin{corrige}[Exercice — Reconstitution du script]
\begin{enumerate}
\item \textbf{Health} (Titre de l'émission : \eng{Health Watch}).
\item \textbf{warns} (3e pers. sg présent) — \textbf{dangers} — \textbf{abuse} (nom indénombrable).
\item \textbf{consumption} — \textbf{teenagers} — \textbf{doubled} (past participle après \eng{has} $\rightarrow$ \eng{has doubled} = Present Perfect).
\item \textbf{damages} (3e pers. sg) — \textbf{lungs} — \textbf{causes} (3e pers. sg).
\item \textbf{Hard} — \textbf{lead} (présent pluriel après \eng{drugs}) — \textbf{addiction}.
\item \textbf{Elderly} — \textbf{special} — \textbf{ageing} (sujet gérondif/nom, orthographe UK).
\item \textbf{advice} (indénombrable, singulier).
\end{enumerate}
\textbf{Notes de grammaire :} \eng{Has doubled} = Present Perfect (lien passé-présent). \eng{Lead} = présent (sujet pluriel \eng{drugs}). \eng{Leads} serait pour \eng{a drug} (singulier). \eng{Ageing} s'écrit avec \eng{e} en anglais britannique.
\end{corrige}

\coursec{Compréhension : questions variées}

Après avoir découvert le script d'écoute et son glossaire dans la section précédente, nous passons maintenant à l'étape cruciale : l'exploitation active du document sonore ou écrit. Au Baccalauréat comme en classe, la compréhension ne se limite pas à « avoir compris vaguement » ; elle exige de \textbf{relever des informations précises}, de \textbf{justifier ses réponses par des citations} et de \textbf{reformuler en anglais correct}. Cette section vous entraîne aux quatre grands types de questions qui reviennent systématiquement sur ce thème : compréhension globale, Vrai/Faux justifié, QCM et réponses courtes.

\courssub{Questions de compréhension globale (Global Understanding)}

Ces questions vérifient que vous avez saisi le \textbf{sujet principal}, les \textbf{intervenants} et le \textbf{type de document}. Elles portent sur le \eng{gist} (l'idée générale). Relisez (ou réécoutez) les premières et dernières phrases du script : elles contiennent souvent ces informations.

\begin{exemplebox}[Questions globales types — Réponses modèles]
\begin{itemize}
    \item \textbf{What is the programme called?} \\
    \eng{The programme is called “Health Watch”.} « L'émission s'appelle “Health Watch”. »
    \item \textbf{Who is the host / presenter?} \\
    \eng{The host is Jane Mbarga.} « L'animatrice est Jane Mbarga. »
    \item \textbf{Who is interviewed? What is his/her profession?} \\
    \eng{Dr. Samuel Ngono, a public health specialist, is interviewed.} « Le Dr Samuel Ngono, spécialiste en santé publique, est interviewé. »
    \item \textbf{What is the main topic of the programme?} \\
    \eng{The main topic is health problems related to drug, tobacco, and alcohol abuse among youth, and aging.} « Le sujet principal est les problèmes de santé liés à l'abus de drogues, tabac, alcool chez les jeunes, et le vieillissement. »
    \item \textbf{Where and when does the scene take place?} \\
    \eng{It takes place in a radio studio in Yaoundé, during a live broadcast.} « Ça se passe dans un studio radio à Yaoundé, lors d'une émission en direct. »
\end{itemize}
\end{exemplebox}

\begin{propriete}[Stratégie : Repérer les indices de globalité]
\begin{itemize}
    \item \textbf{Le titre / le générique :} Souvent donné dans l'énoncé ou dit au tout début (\eng{“Welcome to Health Watch...”}).
    \item \textbf{Les noms propres :} Noms des personnes, lieux, organisations (\eng{Dr. Ngono, Yaoundé, Ministry of Health}).
    \item \textbf{Les mots de liaison de structure :} \eng{Today we focus on... / Our guest is... / Let's talk about... / In conclusion...}
    \item \textbf{Le temps verbal :} Le présent simple pour les faits généraux, le présent continu pour l'action en cours (\eng{“We are discussing...”}).
\end{itemize}
\end{propriete}

\courssub{Exercice Vrai / Faux avec justification (True / False with Justification)}

C'est l'exercice roi du Bac. \textbf{La réponse “True” ou “False” seule ne rapporte aucun point.} Le point est accordé \textbf{uniquement si la justification est exacte et citée}. La justification doit être une \textbf{phrase complète} (ou un segment de phrase) copiée \textbf{mot pour mot} du texte, entre guillemets anglais “ ”, qui prouve votre réponse.

\begin{methode}[Réussir le Vrai / Faux justifié — 4 étapes]
\begin{enumerate}
    \item \textbf{Lisez l'affirmation} et soulignez les \cle{mots-clés} (chiffres, noms, verbes d'action, négations).
    \item \textbf{Scannez le texte} pour localiser le passage correspondant (mots-clés ou synonymes).
    \item \textbf{Comparez} : l'affirmation résume-t-elle exactement l'information du texte ? Attention aux pièges : \eng{some} $\neq$ \eng{all}, \eng{increase} $\neq$ \eng{double}, \eng{young people} $\neq$ \eng{teenagers}.
    \item \textbf{Rédigez :} \eng{True. “Quote from text.”} ou \eng{False. “Quote from text.”} (La citation doit contredire l'affirmation si c'est False).
\end{enumerate}
\end{methode}

\begin{attention}[Pièges fréquents chez les francophones]
\begin{itemize}
    \item \textbf{Traduire la citation :} Interdit. La citation doit être en \textbf{anglais}, identique au texte.
    \item \textbf{Citer trop long :} Ne copiez pas tout le paragraphe. Sélectionnez la \textbf{phrase unique} qui justifie.
    \item \textbf{Inventer la justification :} Écrire \eng{True. “The text says so.”} vaut 0 point.
    \item \textbf{Oublier les guillemets :} Utilisez “ ” (guillemets anglais), pas les guillemets français « ».
    \item \textbf{Confondre “False” et “Not Mentioned” :} Au Bac camerounais, l'option “Not Mentioned” n'existe généralement pas en Vrai/Faux simple. Si l'info n'est pas dans le texte, c'est souvent \textbf{False} car l'affirmation n'est pas vérifiée par le texte. Lisez la consigne exacte du sujet.
\end{itemize}
\end{attention}

\begin{exoresolu}[Vrai / Faux : “Alcohol consumption among teenagers has doubled in ten years.”]
\textbf{Énoncé :} D'après le script, déterminez si l'affirmation est True ou False et justifiez par une citation.

\textbf{Analyse :}
\begin{enumerate}
    \item \textbf{Mots-clés :} \eng{Alcohol consumption}, \eng{teenagers}, \eng{doubled}, \eng{ten years}.
    \item \textbf{Recherche dans le script (paragraphe 2, intervention de Dr Ngono) :} \eng{“Recent surveys show that alcohol consumption among teenagers has doubled over the last decade.”}
    \item \textbf{Comparaison :} \eng{doubled over the last decade} = \eng{doubled in ten years}. \eng{Teenagers} = \eng{teenagers}. L'information est identique.
    \item \textbf{Décision :} \textbf{True}.
\end{enumerate}

\textbf{Réponse attendue au Bac :}
\begin{center}
\eng{True. “Recent surveys show that alcohol consumption among teenagers has doubled over the last decade.”}
\end{center}

\textbf{Variante piège (False) :} \eng{Alcohol consumption among adults has doubled in ten years.}
\textbf{Justification :} \eng{False. “Recent surveys show that alcohol consumption among \textbf{teenagers} has doubled over the last decade.”} (Le texte parle des adolescents, pas des adultes).
\end{exoresolu}

\begin{exemplebox}[Autres exemples Vrai/Faux sur le thème]
\begin{center}
\begin{tabularx}{\linewidth}{|X|X|X|}
\hline
\rowcolor{popblueL} \textbf{Affirmation} & \textbf{Réponse} & \textbf{Justification (Citation)} \\
\hline
Tobacco kills more people than road accidents in Cameroon. & True & “Tobacco remains the leading cause of preventable death, killing more than road accidents.” \\
\hline
The doctor recommends vaping as a safe alternative to smoking. & False & “Vaping is not a safe alternative; it damages the lungs and creates addiction.” \\
\hline
Hard drugs like cocaine and heroin are mentioned. & True & “Hard drugs such as cocaine and heroin destroy the nervous system.” \\
\hline
Aging improves memory and physical strength. & False & “Aging weakens the body and the memory, making seniors more vulnerable.” \\
\hline
\end{tabularx}
\end{center}
\end{exemplebox}

\courssub{QCM de type Baccalauréat (Multiple Choice Questions)}

Les QCM testent la compréhension de détail, l'inférence (déduction) ou le vocabulaire en contexte. Au Bac, il y a souvent 4 options (A, B, C, D). Une seule est correcte.

\begin{methode}[Stratégie d'élimination pour QCM — 3 étapes]
\begin{enumerate}
    \item \textbf{Lisez la question} (pas les options d'abord) et identifiez l'information cherchée.
    \item \textbf{Repérez la zone du texte} qui contient la réponse.
    \item \textbf{Éliminez les distracteurs (réponses fausses) :}
    \begin{itemize}
        \item Réponse \textbf{contraire} au texte.
        \item Réponse \textbf{non mentionnée} (inventée).
        \item Réponse \textbf{trop large} ou \textbf{trop restrictive} (ex: \eng{all} au lieu de \eng{some}).
        \item Réponse qui utilise un \textbf{mot du texte} mais change le sens (piège lexical).
    \end{itemize}
    \item \textbf{Choisissez l'option restante} et vérifiez qu'elle correspond \textbf{exactement} au sens du texte.
\end{enumerate}
\end{methode}

\begin{exoresolu}[QCM : According to Dr. Ngono, tobacco mainly damages...]
\textbf{Question :} According to Dr. Ngono, tobacco mainly damages:
\begin{itemize}
    \item[a)] the liver
    \item[b)] the lungs
    \item[c)] the memory
    \item[d)] the heart
\end{itemize}

\textbf{Analyse :}
\begin{enumerate}
    \item \textbf{Question :} Organe \textbf{principalement} endommagé par le tabac selon le Dr Ngono.
    \item \textbf{Zone texte (Paragraphe 3) :} \eng{“Dr Ngono: ‘Tobacco smoke contains over 7,000 chemicals. It primarily attacks the \textbf{respiratory system}, causing chronic bronchitis and lung cancer. While it affects the heart and blood vessels, the \textbf{lungs} are the first target.’”}
    \item \textbf{Élimination :}
    \begin{itemize}
        \item[a)] \eng{the liver} : Non mentionné comme cible principale (plutôt alcool). $\rightarrow$ \textbf{Éliminé}.
        \item[c)] \eng{the memory} : Associé au vieillissement / drogues dures dans le texte. $\rightarrow$ \textbf{Éliminé}.
        \item[d)] \eng{the heart} : Mentionné comme affecté (\eng{affects the heart}) mais \textbf{pas} la cible \textbf{principale} (\eng{first target}). $\rightarrow$ \textbf{Éliminé}.
        \item[b)] \eng{the lungs} : Correspond à \eng{respiratory system}, \eng{lung cancer}, \eng{first target}. $\rightarrow$ \textbf{Bonne réponse}.
    \end{itemize}
\end{enumerate}

\textbf{Réponse :} \textbf{b) the lungs}
\end{exoresolu}

\begin{exemplebox}[Autres QCM types — Entraînez-vous à éliminer]
\begin{enumerate}
    \item \textbf{What does Dr. Ngono say about aging?}
    \begin{itemize}
        \item[a)] It is a disease that can be cured.
        \item[b)] It weakens the body but wisdom increases. (Correct : \eng{“Aging weakens the body and the memory, yet it brings experience.”})
        \item[c)] It has no impact on health if you exercise.
        \item[d)] It only affects people over 80.
    \end{itemize}
    \item \textbf{The expression “say no to drugs” in the last paragraph is an example of:}
    \begin{itemize}
        \item[a)] A medical prescription.
        \item[b)] A prevention slogan / advice. (Correct : \eng{“My advice to youth: say no to drugs.”})
        \item[c)] A legal obligation.
        \item[d)] A description of addiction symptoms.
    \end{itemize}
\end{enumerate}
\end{exemplebox}

\courssub{Réponses courtes et reformulation (Short Answers \& Reformulation)}

Les questions \eng{Wh-} (Who, What, Where, When, Why, How, How many/much) exigent une \textbf{phrase complète} (Sujet + Verbe + Complément). Répondre par un seul mot (\eng{Drugs.}) ou un groupe nominal (\eng{To say no.}) est \textbf{pénalisé} (souvent 0,5 pt sur 1 ou 0 pt selon la rigueur).

\begin{attention}[Règle d'or : Phrase complète obligatoire]
\begin{itemize}
    \item \textbf{Interdit :} \eng{Name two hard drugs.} $\rightarrow$ \eng{Cocaine and heroin.} (Réponse télégraphique).
    \item \textbf{Exigé :} \eng{The two hard drugs mentioned are cocaine and heroin.} OU \eng{Cocaine and heroin are the two hard drugs mentioned.}
    \item \textbf{Interdit :} \eng{What advice does the doctor give?} $\rightarrow$ \eng{Say no to drugs.}
    \item \textbf{Exigé :} \eng{The doctor advises young people to say no to drugs.} OU \eng{He advises them to avoid drugs.}
\end{itemize}
\end{attention}

\begin{propriete}[Structure de la réponse courte idéale]
\begin{center}
\begin{tabular}{|l|l|l|}
\hline
\rowcolor{popblueL} \textbf{Type de question} & \textbf{Début de réponse type (Subject + Verb)} & \textbf{Exemple} \\
\hline
\eng{Who...?} & \eng{The person who... is...} / \eng{It is...} & \eng{The expert interviewed is Dr. Ngono.} \\
\eng{What...?} & \eng{The text says that...} / \eng{It is...} & \eng{The main topic is drug abuse.} \\
\eng{Where...?} & \eng{It takes place in...} / \eng{The scene is set in...} & \eng{It takes place in a radio studio.} \\
\eng{When...?} & \eng{It happened...} / \eng{The event takes place...} & \eng{The broadcast is live today.} \\
\eng{Why...?} & \eng{Because...} / \eng{The reason is that...} & \eng{Because tobacco damages the lungs.} \\
\eng{How many / much...?} & \eng{There are...} / \eng{The number is...} & \eng{30 percent of young adults smoke.} \\
\hline
\end{tabular}
\end{center}
\end{propriete}

\begin{exemplebox}[Réponses courtes modèles — Phrases complètes]
\begin{itemize}
    \item \textbf{Name two hard drugs mentioned in the text.} \\
    \eng{The two hard drugs mentioned in the text are cocaine and heroin.}
    \item \textbf{What advice does the doctor give to young people?} \\
    \eng{The doctor advises young people to say no to drugs and to avoid peer pressure.}
    \item \textbf{How many young adults smoke regularly according to the survey?} \\
    \eng{Nearly 30 percent of young adults smoke regularly according to the survey.}
    \item \textbf{Why is aging a concern for the health system?} \\
    \eng{Aging is a concern because it weakens the body and the memory, increasing healthcare needs.}
\end{itemize}
\end{exemplebox}

\courssub{Relevé des chiffres et données chiffrées (Figures and Statistics)}

Les chiffres (pourcentages, années, âges, quantités) sont des \textbf{informations à haut coefficient} : ils sont faciles à repérer (scan visuel ou auditif) et font l'objet de questions précises (Vrai/Faux, QCM, Complétion). Ne les négligez jamais.

\begin{exemplebox}[Chiffres clés du script — Que représentent-ils ?]
\begin{center}
\begin{tabularx}{\linewidth}{|X|X|X|}
\hline
\rowcolor{popblueL} \textbf{Chiffre / Donnée} & \textbf{Localisation dans le script} & \textbf{Sens / Représentation} \\
\hline
\eng{30 per cent (30\%)} & \eng{“Nearly \textbf{30 per cent} of young adults smoke regularly.”} & Proportion de jeunes adultes (18-25 ans) fumeurs réguliers. \\
\hline
\eng{doubled in ten years / over the last decade} & \eng{“Alcohol consumption among teenagers has \textbf{doubled} over the \textbf{last decade}.”} & Évolution : multiplication par 2 de la consommation d'alcool chez les ados sur 10 ans. \\
\hline
\eng{7,000 chemicals} & \eng{“Tobacco smoke contains over \textbf{7,000 chemicals}.”} & Nombre de substances chimiques dans la fumée (argument toxicologique). \\
\hline
\eng{leading cause} / \eng{number one} & \eng{“Tobacco remains the \textbf{leading cause} of preventable death.”} & Rang : Cause \#1 de mortalité évitable (devant accidents route, alcool). \\
\hline
\eng{15 to 35 years old} & \eng{“The most affected age group is \textbf{15 to 35 years old}.”} & Tranche d'âge cible pour la prévention (jeunes actifs, étudiants). \\
\hline
\end{tabularx}
\end{center}
\end{exemplebox}

\begin{methode}[Bien traiter les chiffres à l'écoute (Listening)]
\begin{enumerate}
    \item \textbf{Anticipez :} Avant d'écouter, soulignez dans les questions les mots \eng{How many}, \eng{percentage}, \eng{year}, \eng{age}, \eng{number}.
    \item \textbf{Écoutez les “signaux sonores” :} Ralentissement du débit, intonation montante, mots \eng{nearly}, \eng{over}, \eng{about}, \eng{approximately}, \eng{exactly}.
    \item \textbf{Notez en chiffres arabes} (30, 10, 7000) pas en lettres (thirty, ten) : plus vite, moins d'erreurs d'orthographe.
    \item \textbf{Vérifiez l'unité :} \eng{per cent} (\%), \eng{years} (ans), \eng{people} (personnes), \eng{chemicals} (substances).
\end{enumerate}
\end{methode}

\courssub{Correction collective : la preuve par le texte (Textual Evidence)}

En classe comme à l'examen, une réponse sans preuve est une opinion, pas une réponse de compréhension. L'entraînement à la \textbf{citation textuelle} (textual evidence) est la compétence transversale n°1.

\begin{methode}[Comment justifier collectivement (Modèle pour le tableau)]
\begin{enumerate}
    \item L'enseignant (ou l'élève au tableau) écrit la \textbf{question}.
    \item La classe propose la \textbf{réponse} (True/False, Option A/B/C/D, Phrase complète).
    \item \textbf{Exigence :} « \textbf{Où est-ce écrit ? Montrez-moi la phrase.} »
    \item L'élève lit la \textbf{phrase exacte} du script (ou la dicte).
    \item L'enseignant valide : \textbf{Copier-coller de la phrase} au tableau entre guillemets “ ”.
    \item \textbf{Analyse linguistique (bonus) :} Souligner le mot-clé / synonyme / structure qui fait le lien Question $\leftrightarrow$ Texte.
\end{enumerate}
\end{methode}

\begin{exemplebox}[Modèle de correction au tableau — Exemple “False”]
\textbf{Question :} \eng{True or False? “Vaping is presented as a safe alternative to smoking.” Justify.}

\textbf{Échange classe :}
\begin{itemize}
    \item \textbf{Élève A :} False.
    \item \textbf{Prof :} Justification ? La phrase exacte ?
    \item \textbf{Élève B :} \eng{“Vaping is not a safe alternative; it damages the lungs and creates addiction.”}
    \item \textbf{Prof :} Parfait. On note : \eng{False. “Vaping is not a safe alternative; it damages the lungs and creates addiction.”}
    \item \textbf{Analyse :} Le mot \eng{not} nie l'affirmation. \eng{Safe alternative} repris. \eng{Damages} = abîme (contraire de safe).
\end{itemize}
\end{exemplebox}

\courssub{Entraînement à la reformulation (Paraphrasing Practice)}

Au-delà de la citation, le Bac peut demander : \eng{In your own words, explain why...} ou \eng{Rewrite the sentence using...}. Savoir reformuler prouve une compréhension profonde et évite le plagiat pur.

\begin{propriete}[Techniques de reformulation (Paraphrasing)]
\begin{center}
\begin{tabularx}{\linewidth}{|X|X|X|}
\hline
\rowcolor{popblueL} \textbf{Technique} & \textbf{Exemple Original} & \textbf{Reformulation} \\
\hline
\textbf{Synonymes (Vocabulaire)} & \eng{Tobacco \textbf{damages} the lungs.} & \eng{Tobacco \textbf{harms / injures / affects} the lungs.} \\
\hline
\textbf{Changement de structure (Grammaire)} & \eng{\textbf{Dr Ngono advises} youth \textbf{to say no}.} & \eng{\textbf{Youth are advised by Dr Ngono} \textbf{to refuse} drugs.} (Passif) \\
\hline
\textbf{Nominalisation / Verbalisation} & \eng{\textbf{Consumption} has \textbf{doubled}.} & \eng{People \textbf{consume} \textbf{twice as much} alcohol.} \\
\hline
\textbf{Synthèse (Regroupement)} & \eng{It weakens the body. It weakens the memory.} & \eng{It weakens \textbf{both the body and the memory}.} \\
\hline
\end{tabularx}
\end{center}
\end{propriete}

\begin{exoresolu}[Reformulation guidée]
\textbf{Consigne :} Reformulez la phrase suivante sans en changer le sens, en utilisant le mot entre parenthèses.
\begin{enumerate}
    \item \eng{“Aging weakens the body and the memory.”} (\textbf{make}) \\
    \eng{Aging \textbf{makes the body and the memory weaker.}}
    \item \eng{“Teenagers must avoid peer pressure.”} (\textbf{should}) \\
    \eng{Teenagers \textbf{should avoid peer pressure.}}
    \item \eng{“Nearly 30 percent of young adults smoke.”} (\textbf{smokers}) \\
    \eng{\textbf{Nearly 30 percent of young adults are smokers.}}
    \item \eng{“Dr Ngono says tobacco kills more than road accidents.”} (\textbf{leading cause}) \\
    \eng{\textbf{Tobacco is the leading cause of death, ahead of road accidents, according to Dr Ngono.}}
\end{enumerate}
\end{exoresolu}

\begin{popfigure}
\begin{tikzpicture}[
    node distance=0.8cm and 1.2cm,
    box/.style={draw=popblue, thick, rounded corners, fill=popblueL, text width=3.8cm, align=center, minimum height=1.8cm, font=\small},
    arrow/.style={-Stealth, thick, popdark}
]
% Nœuds
\node[box, align=center] (gist){\textbf{Compréhension Globale (Gist)}\\\eng{What is the programme? Who speaks? Main topic?}\\\textit{Teste : Idée générale, contexte}};
\node[box, right=of gist, align=center] (tf){\textbf{Vrai / Faux Justifié}\\\eng{True/False + “Quote”}\\\textit{Teste : Localisation précise, lecture fine, preuve textuelle}};
\node[box, below=of gist, align=center] (mcq){\textbf{QCM (Bac)}\\\eng{4 options, 1 correcte}\\\textit{Teste : Détail, inférence, vocabulaire en contexte, élimination}};
\node[box, right=of mcq, align=center] (sa){\textbf{Réponses Courtes (Wh-)}\\\eng{Full sentence: Subj + Verb}\\\textit{Teste : Reformulation, grammaire, précision, chiffres}};
\node[box, below=of mcq, align=center] (fig){\textbf{Chiffres / Données}\\\eng{30\%, doubled, 7000, 15-35}\\\textit{Teste : Écoute sélective, prise de notes, exactitude}};

% Flèches
\draw[arrow] (gist) -- (tf);
\draw[arrow] (tf) -- (sa);
\draw[arrow] (gist) -- (mcq);
\draw[arrow] (mcq) -- (fig);
\draw[arrow] (fig) -- (sa);
\end{tikzpicture}
\legende{Tableau récapitulatif : Types de questions et compétences testées en compréhension (Listening/Reading).}
\end{popfigure}

\courssub{Exercices d'application (Entraînement Bac)}

\begin{exercice}[Compréhension complète — Simulation Bac]
\textbf{Consignes :} Répondez en anglais. Pour les questions 1 à 3, justifiez par une citation du script (guillemets “ ”). Pour les questions 4 à 7, répondez par une phrase complète.

\textit{Script de référence (rappel section 2) :} \eng{“Welcome to Health Watch... Dr Ngono: ‘Recent surveys show that alcohol consumption among teenagers has doubled over the last decade... Tobacco primarily attacks the respiratory system... Hard drugs such as cocaine and heroin destroy the nervous system... Aging weakens the body and the memory... My advice to youth: say no to drugs.’”} 

\begin{enumerate}
    \item \textbf{True or False?} “Alcohol consumption among teenagers has increased by 50\% in ten years.” Justify.
    \item \textbf{True or False?} “Dr Ngono states that vaping helps smokers quit.” Justify.
    \item \textbf{MCQ:} According to the text, the most affected age group by drug abuse is:
    \begin{itemize}
        \item[a)] Children under 10.
        \item[b)] Teenagers only.
        \item[c)] Young adults aged 15 to 35.
        \item[d)] Seniors over 65.
    \end{itemize}
    \item \textbf{Name two hard drugs mentioned by Dr Ngono.}
    \item \textbf{What is Dr Ngono's final advice to young people?}
    \item \textbf{Figure search:} What does “30 per cent” refer to in the survey?
    \item \textbf{Reformulation:} Rewrite: “Tobacco primarily attacks the respiratory system.” Use \textbf{main target}.
\end{enumerate}
\end{exercice}

\begin{corrige}[Exercice Compréhension complète]
\begin{enumerate}
    \item \textbf{False.} \eng{“Recent surveys show that alcohol consumption among teenagers has \textbf{doubled} over the last decade.”} (Doubled = 100\% increase, pas 50\%).
    \item \textbf{False.} \eng{“Vaping is \textbf{not a safe alternative}; it damages the lungs and creates addiction.”} (Le texte dit le contraire).
    \item \textbf{c) Young adults aged 15 to 35.} Justification : \eng{“The most affected age group is \textbf{15 to 35 years old}.”}
    \item \textbf{The two hard drugs mentioned are cocaine and heroin.} (Ou : \eng{Cocaine and heroin are the two hard drugs mentioned.})
    \item \textbf{Dr Ngono advises young people to say no to drugs.} (Ou : \eng{His advice is to say no to drugs.})
    \item \textbf{It refers to the proportion of young adults who smoke regularly.} (Ou : \eng{Nearly 30 percent of young adults smoke regularly.})
    \item \textbf{The respiratory system is the main target of tobacco.} (Ou : \eng{Tobacco's main target is the respiratory system.})
\end{enumerate}
\end{corrige}

\begin{aretenir}[Check-list avant l'épreuve de compréhension]
\begin{itemize}
    \item [ ] Je lis les questions \textbf{avant} d'écouter/lire.
    \item [ ] Je surligne les \textbf{mots-clés} dans les questions (noms, verbes, chiffres, \eng{Wh-}).
    \item [ ] Au Vrai/Faux : \textbf{True/False + Citation exacte entre “ ”}. Pas de traduction.
    \item [ ] Au QCM : J'\textbf{élimine} 3 réponses fausses avant de choisir.
    \item [ ] Aux réponses courtes : \textbf{Phrase complète} (Sujet + Verbe). Pas de réponse télégraphique.
    \item [ ] Je note les \textbf{chiffres} en chiffres arabes (30, 10, 7000).
    \item [ ] Je relis mes réponses : \textbf{orthographe anglais}, majuscules (I, noms propres), point final.
\end{itemize}
\end{aretenir}

\coursec{Lesson 25 — Listening: Texts on keeping informed about health problems related to drugs, tobacco, alcohol abuse, and aging}

\courssub{Mise en route et anticipation}

\begin{experience}[Activité de démarrage : Images et mots-clés]
Observez les trois affiches de sensibilisation ci-dessous (imaginées). En binôme, associez chaque affiche à son thème principal et listez 5 mots anglais que vous vous attendez à entendre dans un texte audio sur ce thème.

\begin{center}
\begin{tabular}{|p{4cm}|p{4cm}|p{4cm}|}
\hline
\rowcolor{popblueL} \textbf{Affiche A} & \textbf{Affiche B} & \textbf{Affiche C} \\
\hline
Image : Jeune homme tremblant, tenant une seringue / comprimés. Slogan : \eng{“Your future is not for sale.”} & Image : Paquet de cigarettes neutre avec poumon noir. Slogan : \eng{“Smoking kills. Quit now.”} & Image : Personne âgée seule, regard vide, déambulateur. Slogan : \eng{“Respect the elderly. Care for them.”} \\
\hline
\textbf{Thème :} \eng{Drug abuse} & \textbf{Thème :} \eng{Tobacco / Alcohol} & \textbf{Thème :} \eng{Aging / Elderly care} \\
\hline
\textbf{Mots attendus :} & \textbf{Mots attendus :} & \textbf{Mots attendus :} \\
\eng{addict, rehab, overdose, dealer, tramadol} & \eng{lung cancer, nicotine, quit, cirrhosis, drunk} & \eng{dementia, care home, memory loss, retirement, loneliness} \\
\hline
\end{tabular}
\end{center}

\textbf{Consigne :} Avant l'écoute, lisez les questions de compréhension (Section 3). Soulignez les \cle{mots-clés} (noms, verbes, chiffres, négations). Prédisez le type de réponses attendues (un nom, un chiffre, une raison, Vrai/Faux).
\end{experience}

\courssub{Scripts d'écoute et glossaires}

\emph{Le professeur lit les trois textes ci-dessous à vitesse normale (ou les diffuse via l'application OnBuch+). Les élèves ne lisent pas le script pendant la première écoute.}

\begin{textebox}[Text 1 : Interview — “A Second Chance in Buea” (Drugs)]
\textbf{Journalist:} Good morning. Today we talk about the tramadol crisis affecting our youth. With us is Mr. Elias, a recovered addict and now a counsellor at the Buea Rehabilitation Centre. Good morning, Elias.

\textbf{Elias:} Good morning. Thank you for having me.

\textbf{Journalist:} Elias, how did it start for you?

\textbf{Elias:} I was a university student in Douala. Pressure was high. A senior student offered me \textbf{cannabis} to “relax”. It became a daily habit. Within six months, I moved to \textbf{tramadol} — 225mg tablets. They were cheap, easy to get from \textbf{dealers} near campus. I became a full-blown \textbf{addict}. I stole my sister’s school fees. I lost 15 kilos.

\textbf{Journalist:} What was the turning point?

\textbf{Elias:} An \textbf{overdose}. My friends found me unconscious in a \emph{bendskin} seat. I woke up in the emergency ward. A doctor told me: “One more dose and you die \textbf{of} cardiac arrest.” That fear saved me. I entered \textbf{rehab} for nine months. It was hell — \textbf{withdrawal symptoms}, insomnia, depression. But I \textbf{gave up drugs}. Today, I \textbf{warn} others: “Don’t try it. Not even once.”
\end{textebox}

\begin{center}
\begin{tabularx}{\linewidth}{|l|X|}
\hline
\rowcolor{popblueL} \textbf{Mot / Expression} & \textbf{Glossaire (Nature, Sens, Prononciation approx.)} \\
\hline
\textbf{rehabilitation centre / rehab} (n.) & centre de désintoxication, centre de réadaptation (« ri-hè-bi-li-té-chone ») \\
\textbf{senior student} (n.) & étudiant en année supérieure, aîné \\
\textbf{pressure} (n.) & pression (stress, contrainte) \\
\textbf{habit} (n.) & habitude (ici : dépendance) \\
\textbf{full-blown addict} (n.) & toxicomane confirmé, accro total \\
\textbf{unconscious} (adj.) & inconscient (évanoui) (« an-con-chious ») \\
\textbf{emergency ward} (n.) & service des urgences \\
\textbf{cardiac arrest} (n.) & arrêt cardiaque \\
\textbf{withdrawal symptoms} (n. pl.) & symptômes de sevrage (« with-draw-al simp-tomes ») \\
\textbf{insomnia} (n.) & insomnie (« in-som-ni-a ») \\
\hline
\end{tabularx}
\end{center}

\begin{textebox}[Text 2 : Radio Health Report — “Weekend Tragedies on the Ring Road” (Alcohol \& Tobacco)]
\textbf{Presenter:} This is Cameroon Radio Health Watch. I’m Martha Nfor. Last weekend, the Bamenda Ring Road recorded three fatal accidents linked to \textbf{drink-driving}. A driver, three times over the legal alcohol limit, crashed into a taxi. Two passengers died on the spot.

Dr. Atem, a hepatologist at the Regional Hospital, warns: “We see \textbf{cirrhosis} of the \textbf{liver} in patients as young as 30. Ten years ago, it was 50. \textbf{Binge drinking} on cheap local spirits and imported beers is the cause.”

On tobacco, the Ministry of Health reports that \textbf{13\%} of youths aged 13–15 already \textbf{smoke}. \textbf{Graphic warnings} on packets — showing rotting teeth and black lungs — have reduced sales by only 2\%. \textbf{Second-hand smoke} in bars and taxis remains a major risk for non-smokers. \textbf{Nicotine} addiction is powerful. The message is clear: \textbf{quit smoking} now. Call the free helpline 119 for support.
\end{textebox}

\begin{center}
\begin{tabularx}{\linewidth}{|l|X|}
\hline
\rowcolor{popblueL} \textbf{Mot / Expression} & \textbf{Glossaire (Nature, Sens, Prononciation approx.)} \\
\hline
\textbf{fatal} (adj.) & mortel (« fa-til ») \\
\textbf{on the spot} (loc.) & sur le coup, sur place \\
\textbf{hepatologist} (n.) & hépatologue (spécialiste du foie) (« hé-pa-to-lo-giste ») \\
\textbf{binge drinking} (n.) & beuverie expresse, consommation massive d'alcool en peu de temps \\
\textbf{spirits} (n. pl.) & spiritueux (alcools forts : whisky, gin, odontol) \\
\textbf{graphic warnings} (n. pl.) & avertissements illustrés (images chocs) \\
\textbf{rotting} (adj.) & pourri, qui pourrit (« ro-ting ») \\
\textbf{helpline} (n.) & numéro vert, ligne d'assistance téléphonique \\
\hline
\end{tabularx}
\end{center}

\begin{textebox}[Text 3 : Feature Story — “The Forgotten Generation” (Aging)]
\textbf{Reporter:} In the quiet village of Muyuka, 78-year-old Mama Agnes sits on her verandah. She stares at the path, waiting for children who migrated to Douala or abroad. “My \textbf{memory} fails me,” she whispers. “I forget I ate. I forget names. The doctor says it’s early \textbf{dementia}.”

Cameroon’s population is \textbf{ageing}. Life expectancy has risen to 60. But \textbf{elderly people} face new challenges: \textbf{memory loss}, \textbf{arthritis}, isolation. Traditional \textbf{care} — children looking after parents — is \textbf{eroding}. Urban migration, poverty, and the \textbf{breakdown} of the extended family leave many seniors alone.

Few \textbf{care homes} exist; they are private and costly. The Ministry of Social Affairs launched a \textbf{“Senior Watch”} programme: community volunteers visit isolated elders weekly. “We \textbf{keep informed} on their health,” says volunteer Pierre. “We \textbf{raise awareness}: ageing is not a curse. It’s a stage of life demanding dignity.”
\end{textebox}

\begin{center}
\begin{tabularx}{\linewidth}{|l|X|}
\hline
\rowcolor{popblueL} \textbf{Mot / Expression} & \textbf{Glossaire (Nature, Sens, Prononciation approx.)} \\
\hline
\textbf{verandah} (n.) & véranda, galerie couverte (« vé-ran-da ») \\
\textbf{dementia} (n.) & démence (troubles cognitifs graves) (« di-men-cha ») \\
\textbf{life expectancy} (n.) & espérance de vie \\
\textbf{arthritis} (n.) & arthrite, rhumatismes (« ar-thrai-tis ») \\
\textbf{eroding} (part.) & qui s'érode, qui se dégrade progressivement \\
\textbf{extended family} (n.) & famille élargie \\
\textbf{breakdown} (n.) & effondrement, dislocation \\
\textbf{volunteer} (n.) & bénévole (« vol-un-ti-eur ») \\
\textbf{isolated} (adj.) & isolé (« ai-so-lèï-tid ») \\
\textbf{demanding} (adj.) & exigeant, qui demande beaucoup \\
\hline
\end{tabularx}
\end{center}

\courssub{Compréhension : questions variées}

\begin{exercice}[Questions sur les trois textes]
\textbf{A. Global Understanding (Texts 1, 2, 3)} — Tick the correct answer (\eng{True} / \eng{False} / \eng{Not Mentioned}).
\begin{enumerate}
    \item Text 1: Elias started using drugs because he was curious. \hfill \eng{T / F / NM}
    \item Text 2: The legal alcohol limit for drivers was exceeded three times. \hfill \eng{T / F / NM}
    \item Text 3: Mama Agnes lives with her children in Muyuka. \hfill \eng{T / F / NM}
    \item Text 1: Tramadol tablets mentioned are 225mg. \hfill \eng{T / F / NM}
    \item Text 2: Graphic warnings have stopped young people from smoking completely. \hfill \eng{T / F / NM}
    \item Text 3: The “Senior Watch” programme uses paid nurses. \hfill \eng{T / F / NM}
\end{enumerate}

\textbf{B. Detailed Comprehension (Text 1)} — Answer in full sentences.
\begin{enumerate}
    \item What were the two main drugs Elias used?
    \item What specific event made Elias decide to stop?
    \item How long did his rehabilitation last?
    \item What is his message to youth today?
\end{enumerate}

\textbf{C. Detailed Comprehension (Text 2)} — Multiple Choice. Choose a, b, or c.
\begin{enumerate}
    \item According to Dr. Atem, cirrhosis now affects people aged:
    \begin{itemize}
        \item[a)] 50 and above \item[b)] 30 and above \item[c)] 40 and above
    \end{itemize}
    \item What percentage of 13–15 year olds smoke?
    \begin{itemize}
        \item[a)] 2\% \item[b)] 13\% \item[c)] 22\%
    \end{itemize}
    \item What is the number to call for help to quit smoking?
    \begin{itemize}
        \item[a)] 112 \item[b)] 119 \item[c)] 911
    \end{itemize}
\end{enumerate}

\textbf{D. Detailed Comprehension (Text 3)} — Short answers (max 10 words).
\begin{enumerate}
    \item What medical condition does Mama Agnes have?
    \item Why are traditional care systems eroding? (Give two reasons).
    \item What do “Senior Watch” volunteers do?
\end{enumerate}
\end{exercice}

\begin{corrige}[Exercice : Compréhension]
\textbf{A. Global Understanding}
1. \textbf{F} (Started due to pressure, offered by senior student).
2. \textbf{T} (“three times over the legal alcohol limit”).
3. \textbf{F} (“children who migrated to Douala or abroad”).
4. \textbf{T} (“225mg tablets”).
5. \textbf{F} (“reduced sales by only 2\%”).
6. \textbf{F} (“community volunteers”).

\textbf{B. Detailed Comprehension (Text 1)}
1. He used \textbf{cannabis} and \textbf{tramadol} (225mg).
2. An \textbf{overdose}; he was found unconscious and a doctor warned him he would die of cardiac arrest.
3. His rehabilitation lasted \textbf{nine months}.
4. “\textbf{Don’t try it. Not even once.}”

\textbf{C. Detailed Comprehension (Text 2)}
1. \textbf{b} (patients as young as 30).
2. \textbf{b} (13\%).
3. \textbf{b} (119).

\textbf{D. Detailed Comprehension (Text 3)}
1. Early \textbf{dementia} / \textbf{memory loss}.
2. Urban migration, poverty, breakdown of the extended family (any two).
3. They visit isolated elders weekly, keep informed on their health, and raise awareness.
\end{corrige}

\courssub{Lexique thématique : addictions et vieillissement}

\emph{Cette section consolide le vocabulaire rencontré dans les scripts. Elle est organisée en cinq champs thématiques, avec collocations, familles de mots et un texte de réinvestissement.}

\courssub{Champ lexical 1 — Les drogues (Drugs)}

\begin{textebox}[Extrait : Témoignage d'un ancien toxicomane (Text 1)]
"I was a university student. A senior student offered me \cle{cannabis} to relax. Soon, I moved to \cle{tramadol} — 225mg tablets. I became a full-blown \cle{addict}. I entered \cle{rehab} to \cle{give up drugs}. Now I \cle{warn} others."
\end{textebox}

\begin{definition}[Vocabulaire de base — Drugs]
\begin{itemize}
    \item \eng{drug} /drʌɡ/ (n.) : drogue, médicament. \emph{Attention : en anglais, \eng{drug} couvre les deux sens. Le contexte détermine le sens (ex. \eng{prescription drugs} = médicaments sur ordonnance ; \eng{hard drugs} = drogues dures).}
    \item \eng{hard drug} (n.) : drogue dure (héroïne, cocaïne, crack, tramadol).
    \item \eng{cannabis} / \eng{marijuana} / \eng{weed} (argot) (n.) : cannabis.
    \item \eng{tramadol} (n.) : tramadol (antalgique opioïde très détourné au Cameroun).
    \item \eng{dealer} / \eng{drug dealer} (n.) : dealer, trafiquant de drogue.
    \item \eng{addict} /ˈædɪkt/ (n.) : toxicomane, accro.
    \item \eng{addiction} /əˈdɪkʃn/ (n.) : toxicomanie, dépendance.
    \item \eng{overdose} /ˈəʊvədəʊs/ (n.) : surdose.
    \item \eng{rehab} /ˈriːhæb/ (n., abbr. de \eng{rehabilitation}) : centre de désintoxication, cure de désintoxication.
    \item \eng{withdrawal symptoms} (n. pl.) : symptômes de sevrage.
    \item \eng{to be addicted to} + V-ing / nom : être dépendant de, être accro à.
    \item \eng{to give up drugs} / \eng{to kick the habit} (argot) : arrêter la drogue, décrocher.
    \item \eng{to warn} (v.) : avertir, mettre en garde. \eng{Warn sb against / about sth}.
\end{itemize}
\end{definition}

\begin{propriete}[Construction : \eng{to be addicted to}]
\eng{Addicted} est un adjectif (participe passé) suivi de la préposition \cle{to}. \eng{To} est ici une préposition, \textbf{pas} l'infinitif. Elle est donc suivie d'un \textbf{gérondif (V-ing)} ou d'un \textbf{nom}.
\begin{center}
\begin{tabular}{|l|l|l|}
\hline
\rowcolor{popblueL} Structure & Exemple & Traduction \\
\hline
be addicted to + V-ing & He is addicted to \textbf{smoking}. & Il est dépendant de la cigarette. \\
be addicted to + nom & She is addicted to \textbf{tramadol}. & Elle est dépendante du tramadol. \\
\hline
\end{tabular}
\end{center}
\end{propriete}

\begin{exemplebox}[Exemples traduits]
\eng{The police arrested a known drug dealer in Buea.} = La police a arrêté un dealer connu à Buea. \\
\eng{Tramadol addiction is a major public health issue in Cameroon.} = La dépendance au tramadol est un problème majeur de santé publique au Cameroun. \\
\eng{He finally decided to give up drugs after ten years.} = Il a finalement décidé d'arrêter la drogue après dix ans. \\
\eng{She warned her brother against bad company.} = Elle a mis en garde son frère contre les mauvaises fréquentations.
\end{exemplebox}

\courssub{Champ lexical 2 — Le tabac (Tobacco)}

\begin{textebox}[Extrait : Reportage sanitaire (Text 2)]
"13\% of youths aged 13–15 already \cle{smoke}. \cle{Graphic warnings} on packets show \cle{rotting} teeth. \cle{Second-hand smoke} in taxis endangers non-smokers. \cle{Nicotine} addiction makes \cle{quitting smoking} hard."
\end{textebox}

\begin{definition}[Vocabulaire de base — Tobacco]
\begin{itemize}
    \item \eng{to smoke} (v.) : fumer.
    \item \eng{smoker} (n.) : fumeur / fumeuse.
    \item \eng{cigarette} (n. comptable) : une cigarette. \textbf{Pluriel : cigarettes.}
    \item \eng{tobacco} /təˈbækəʊ/ (n. \textbf{indénombrable}) : tabac (la substance). \textbf{Jamais 'tobaccos'.}
    \item \eng{lung cancer} (n.) : cancer du poumon.
    \item \eng{second-hand smoke} / \eng{passive smoking} (n.) : tabagisme passif, fumée secondaire.
    \item \eng{to quit smoking} / \eng{to give up smoking} / \eng{to stop smoking} : arrêter de fumer.
    \item \eng{nicotine} (n.) : nicotine.
    \item \eng{graphic warning} (n.) : avertissement illustré (image choc).
\end{itemize}
\end{definition}

\begin{attention}[Faux amis et pièges — Tobacco vs Cigarette]
\begin{itemize}
    \item \textbf{Une cigarette} = \eng{a cigarette} (comptable). \eng{He smokes 20 cigarettes a day.}
    \item \textbf{Du tabac} = \eng{tobacco} (indénombrable, masse). \eng{He buys rolling tobacco.} (Il achète du tabac à rouler).
    \item \textbf{Erreur fréquente :} \eng{*He smokes tobaccos.} $\rightarrow$ \textbf{Correct :} \eng{He smokes tobacco.} ou \eng{He smokes cigarettes.}
    \item \eng{Smoking} (n.) = l'acte de fumer / le tabagisme. \eng{Smoking is forbidden.} (Il est interdit de fumer / Le tabagisme est interdit). Ne pas confondre avec \eng{a smoking} (FR) = \eng{a dinner jacket} / \eng{a tuxedo} (UK/US).
\end{itemize}
\end{attention}

\courssub{Champ lexical 3 — L'alcool (Alcohol)}

\begin{textebox}[Extrait : Reportage routier (Text 2)]
"Weekend \cle{drink-driving} accidents are rising. A driver was \cle{drunk} — three times the legal limit. Long-term abuse leads to \cle{cirrhosis} of the \cle{liver}. \cle{Binge drinking} on cheap spirits is the cause."
\end{textebox}

\begin{definition}[Vocabulaire de base — Alcohol]
\begin{itemize}
    \item \eng{alcohol} /ˈælkəhɒl/ (n. indénombrable) : alcool.
    \item \eng{drunk} /drʌŋk/ (adj.) : ivre, saoul. \emph{Participe passé de \eng{drink}.}
    \item \eng{drunkenness} (n.) : ivresse, beuverie.
    \item \eng{to get drunk} : s'enivrer, devenir ivre.
    \item \eng{liver} /ˈlɪvə(r)/ (n.) : foie. \emph{Prononciation : « live-eur », pas « lai-veur ».}
    \item \eng{cirrhosis} /sɪˈrəʊsɪs/ (n.) : cirrhose.
    \item \eng{drink-driving} (UK) / \eng{drunk driving} / \eng{DUI} (US — Driving Under the Influence) (n.) : conduite en état d'ivresse.
    \item \eng{binge drinking} (n.) : beuverie expresse (consommer beaucoup d'alcool en peu de temps).
    \item \eng{to be on the wagon} (idiom) : ne plus boire d'alcool (être sobre, en sevrage). \emph{Antonyme : \eng{to fall off the wagon} = replonger.}
    \item \eng{alcoholic} (n./adj.) : alcoolique.
\end{itemize}
\end{definition}

\begin{propriete}[Grammaire : \eng{Die of} vs \eng{Die from}]
Avec les maladies, causes internes ou naturelles, on utilise \textbf{\eng{die of}} (mourir \textbf{de}).
\begin{center}
\begin{tabular}{|l|l|}
\hline
\rowcolor{popblueL} Anglais & Français \\
\hline
He died \textbf{of} lung cancer. & Il est mort \textbf{d'un} cancer du poumon. \\
She died \textbf{of} cirrhosis. & Elle est morte \textbf{d'une} cirrhose. \\
He died \textbf{of} a heart attack. & Il est mort \textbf{d'une} crise cardiaque. \\
\hline
\end{tabular}
\end{center}
\eng{Die from} s'utilise pour les causes externes, violentes, accidentelles : \eng{He died from his injuries.} (Il est mort de ses blessures). \textbf{Au Bac, privilégiez \eng{die of} pour les maladies liées aux addictions.}
\end{propriete}

\courssub{Champ lexical 4 — Le vieillissement (Ageing)}

\begin{textebox}[Extrait : Reportage société (Text 3)]
"Cameroon's population is \cle{ageing}. \cle{Elderly people} face \cle{memory loss} and isolation. Traditional \cle{care} is \cle{eroding}. Few \cle{care homes} exist. Volunteers \cle{keep informed} on their health and \cle{raise awareness}."
\end{textebox}

\begin{definition}[Vocabulaire de base — Ageing]
\begin{itemize}
    \item \eng{to age} (v.) : vieillir. \eng{The population is ageing.}
    \item \eng{ageing} / \eng{aging} (n./adj.) : vieillissement / vieillissant. \emph{Orthographe UK : \textbf{ageing} ; US : \textbf{aging}. Au Cameroun : \textbf{ageing}.}
    \item \eng{elderly people} / \eng{the elderly} (n. pl.) : les personnes âgées, les aînés. \emph{Plus respectueux que \eng{old people}.}
    \item \eng{old age} (n.) : la vieillesse, le grand âge.
    \item \eng{memory loss} (n.) : perte de mémoire, troubles de la mémoire.
    \item \eng{dementia} (n.) : démence (maladie d'Alzheimer, etc.).
    \item \eng{arthritis} (n.) : arthrite, rhumatismes.
    \item \eng{to weaken} (v.) : affaiblir, s'affaiblir. \eng{Age weakens the bones.}
    \item \eng{retirement} (n.) : retraite (professionnelle).
    \item \eng{to take care of} / \eng{to look after} : prendre soin de, s'occuper de. \eng{Who takes care of your grandmother?}
    \item \eng{care home} (UK) / \eng{nursing home} (US) (n.) : maison de retraite, EHPAD.
    \item \eng{isolated} (adj.) : isolé.
    \item \eng{volunteer} (n.) : bénévole.
\end{itemize}
\end{definition}

\courssub{Champ lexical 5 — Prévention et santé (Prevention \& Health)}

\begin{definition}[Vocabulaire de base — Prevention \& Health]
\begin{itemize}
    \item \eng{health} (n.) : santé. \eng{Public health} = santé publique.
    \item \eng{disease} /dɪˈziːz/ (n.) : maladie (pathologie précise). \emph{Ne pas confondre avec \eng{illness} (état de malade) ou \eng{sickness} (maladie, nausée).}
    \item \eng{treatment} (n.) : traitement, prise en charge.
    \item \eng{awareness campaign} (n.) : campagne de sensibilisation. \textbf{Collocation forte.}
    \item \eng{to raise awareness} (loc.) : sensibiliser, faire prendre conscience.
    \item \eng{to warn} (v.) : avertir, mettre en garde. \eng{Warn sb against / about sth}.
    \item \eng{to advise} (v.) : conseiller. \eng{Advise sb to do sth.} (structure : \eng{advise + objet + to + infinitif}).
    \item \eng{to keep informed} : se tenir informé. \eng{Keep informed about health risks.}
    \item \eng{healthy lifestyle} (n.) : mode de vie sain, hygiène de vie.
    \item \eng{helpline} (n.) : numéro vert, ligne d'écoute.
\end{itemize}
\end{definition}

\begin{definition}[Awareness]
\eng{Awareness} /əˈweənəs/ (n.) : conscience, sensibilisation, prise de conscience. Vient de \eng{aware} (conscient) + \eng{-ness}.
\eng{An awareness campaign} = une campagne de sensibilisation. C'est une collocation \textbf{très fréquente au Baccalauréat} (sujets d'expression écrite, compréhension).
\end{definition}

\courssub{Collocations fréquentes (Mots « collés » ensemble)}

Apprenez ces blocs inséparables pour gagner en fluidité et éviter les traductions mot à mot.

\begin{center}
\begin{tabularx}{\linewidth}{|X|X|X|}
\hline
\rowcolor{popblueL} Collocation (Anglais) & Traduction & Exemple d'emploi \\
\hline
\eng{to suffer from a disease} & souffrir d'une maladie & Many elderly people \textbf{suffer from} Alzheimer's disease. \\
\eng{to die of cancer} & mourir d'un cancer & He \textbf{died of} lung cancer last year. \\
\eng{to raise awareness} & sensibiliser & NGOs \textbf{raise awareness} on drug abuse. \\
\eng{to lead a healthy life} & mener une vie saine & You should \textbf{lead a healthy life} to age well. \\
\eng{to launch a campaign} & lancer une campagne & The Ministry \textbf{launched an awareness campaign}. \\
\eng{to quit / give up smoking} & arrêter de fumer & It's never too late to \textbf{quit smoking}. \\
\eng{to be addicted to} & être dépendant de & He is \textbf{addicted to} tramadol. \\
\eng{to warn against} & mettre en garde contre & Doctors \textbf{warn against} self-medication. \\
\hline
\end{tabularx}
\end{center}

\courssub{Familles de mots (Word Formation)}

Travailler la dérivation est crucial pour l'expression écrite et le \textit{Use of English}.

\begin{center}
\begin{tabular}{|c|c|c|c|}
\hline
\rowcolor{popblueL} Base / Verbe & Nom (Personne) & Nom (Concept/Action) & Adjectif \\
\hline
\eng{addict} & \eng{addict} & \eng{addiction} & \eng{addicted} / \eng{addictive} \\
\hline
\eng{smoke} & \eng{smoker} & \eng{smoking} & \eng{smoky} (enfumé) \\
\hline
\eng{drink} & \eng{drinker} & \eng{drinking} & \eng{drunk} (ivre) / \eng{drunken} (attrib. : \eng{a drunken driver}) \\
\hline
\eng{age} & — & \eng{ageing} & \eng{aged} (âgé) / \eng{ageing} \\
\hline
\eng{depend} & \eng{dependent} & \eng{dependence} / \eng{dependency} & \eng{dependent} (on) / \eng{independent} \\
\hline
\eng{prevent} & — & \eng{prevention} & \eng{preventive} / \eng{preventable} \\
\hline
\eng{treat} & \eng{patient} / \eng{practitioner} & \eng{treatment} & \eng{treatable} \\
\hline
\eng{aware} (adj.) & — & \eng{awareness} & \eng{aware} (of) \\
\hline
\eng{rehabilitate} & \eng{rehabilitee} (rare) & \eng{rehabilitation} & \eng{rehabilitative} \\
\hline
\end{tabular}
\end{center}

\begin{exemplebox}[Dérivation en contexte]
\eng{The \textbf{addictive} nature of tramadol leads to severe \textbf{addiction}.} (adj. $\rightarrow$ n.) \\
\eng{He is a heavy \textbf{smoker}; \textbf{smoking} has damaged his lungs.} (n. personne $\rightarrow$ n. action) \\
\eng{\textbf{Prevention} is better than cure. \textbf{Preventive} measures save lives.} (n. $\rightarrow$ adj.) \\
\eng{The government must \textbf{warn} young people. A \textbf{warning} on packets is mandatory.} (v. $\rightarrow$ n.) \\
\eng{\textbf{Rehabilitation} takes time. A \textbf{rehabilitated} addict can help others.} (n. $\rightarrow$ adj.)
\end{exemplebox}

\begin{popfigure}
\begin{tikzpicture}[
    node distance=0.6cm and 1.0cm,
    every node/.style={font=\small},
    champ/.style={rectangle, rounded corners, draw=popblue, thick, fill=popblueL, text width=2.5cm, align=center, minimum height=1.1cm},
    mot/.style={rectangle, draw=poporange, fill=poporangeL, text width=4.0cm, align=left, minimum height=0.9cm},
    trad/.style={rectangle, draw=popgreen, fill=popgreenL, text width=4.0cm, align=left, minimum height=0.9cm},
]
% En-têtes
\node[champ] (h1) {\textbf{Champ lexical}};
\node[mot, right=of h1] (h2) {\textbf{Mots clés (English)}};
\node[trad, right=of h2] (h3) {\textbf{Traduction française}};

% Ligne 1 : Drugs
\node[champ, below=of h1] (c1) {1. Drugs};
\node[mot, right=of c1] (m1) {drug, hard drug, cannabis, tramadol, dealer, addict, addiction, overdose, rehab, withdrawal symptoms, to be addicted to, to give up drugs, to warn};
\node[trad, right=of m1] (t1) {drogue, drogue dure, cannabis, tramadol, dealer, toxicomane, toxicomanie, surdose, centre de désintoxication, symptômes de sevrage, être dépendant de, arrêter la drogue, avertir};

% Ligne 2 : Tobacco
\node[champ, below=of c1] (c2) {2. Tobacco};
\node[mot, right=of c2] (m2) {to smoke, smoker, cigarette, tobacco, lung cancer, second-hand smoke, to quit smoking, nicotine, graphic warning};
\node[trad, right=of m2] (t2) {fumer, fumeur, cigarette, tabac, cancer du poumon, tabagisme passif, arrêter de fumer, nicotine, avertissement choc};

% Ligne 3 : Alcohol
\node[champ, below=of c2] (c3) {3. Alcohol};
\node[mot, right=of c3] (m3) {alcohol, drunk, drunkenness, to get drunk, liver, cirrhosis, drink-driving, binge drinking, to be on the wagon, alcoholic};
\node[trad, right=of m3] (t3) {alcool, ivre, ivresse, s'enivrer, foie, cirrhose, conduite en état d'ivresse, beuverie expresse, être sobre (sevrage), alcoolique};

% Ligne 4 : Ageing
\node[champ, below=of c3] (c4) {4. Ageing};
\node[mot, right=of c4] (m4) {to age, ageing, elderly people, old age, memory loss, dementia, arthritis, to weaken, retirement, to take care of, care home, isolated, volunteer};
\node[trad, right=of m4] (t4) {vieillir, vieillissement, personnes âgées, vieillesse, perte de mémoire, démence, arthrite, affaiblir, retraite, prendre soin de, maison de retraite, isolé, bénévole};

% Ligne 5 : Prevention
\node[champ, below=of c4] (c5) {5. Prevention \& Health};
\node[mot, right=of c5] (m5) {health, disease, treatment, awareness campaign, to raise awareness, to warn, to advise, to keep informed, healthy lifestyle, helpline};
\node[trad, right=of m5] (t5) {santé, maladie, traitement, campagne de sensibilisation, sensibiliser, avertir, conseiller, se tenir informé, mode de vie sain, numéro vert};

% Lignes de séparation
\draw[popline, thick] (h1.south west) -- (h3.south east);
\draw[popline] (c1.south west) -- (t1.south east);
\draw[popline] (c2.south west) -- (t2.south east);
\draw[popline] (c3.south west) -- (t3.south east);
\draw[popline] (c4.south west) -- (t4.south east);
\draw[popline] (c5.south west) -- (t5.south east);
\end{tikzpicture}
\legende{Tableau récapitulatif des 5 champs lexicaux : Mots anglais et traduction française.}
\end{popfigure}

\courssub{Texte de réinvestissement : « Youth at the Crossroads »}

Lisez ce texte original qui reprend l'essentiel du vocabulaire vu ci-dessus. Le glossaire suit le texte.

\begin{textebox}[Youth at the Crossroads: A Call for Action in Bamenda]
In the bustling streets of Bamenda, the sight of young people smoking \textbf{cannabis} or sniffing glue has become alarmingly common. Many start with \textbf{cannabis} in secondary school, unaware that it often acts as a \textbf{gateway drug} leading to \textbf{hard drugs} like \textbf{tramadol} or cocaine. A \textbf{dealer} operates near every corner, targeting vulnerable students who want to escape poverty or academic pressure.

The consequences are devastating. \textbf{Addiction} sets in quickly. Users become \textbf{addicted to} the substance, spending their school fees on their \textbf{habit}. \textbf{Overdose} deaths are rising silently in our communities. Meanwhile, \textbf{alcohol} abuse fuels weekend chaos. \textbf{Drink-driving} accidents claim lives on the Bamenda-Bali road. Long-term \textbf{drinkers} face \textbf{cirrhosis} of the \textbf{liver} in their thirties.

Tobacco remains a silent killer. Despite \textbf{graphic warnings} on packets — "\textbf{Smoking kills}" — many youths light their first \textbf{cigarette} at 14. \textbf{Second-hand smoke} endangers non-smokers in taxis and bars. \textbf{Nicotine} addiction makes \textbf{quitting smoking} a battle few win without help.

Our \textbf{ageing} parents suffer too. \textbf{Elderly people} with \textbf{memory loss} or \textbf{dementia} are often abandoned when families cannot \textbf{cope}. \textbf{Care homes} are rare and expensive. The tradition of \textbf{taking care of} grandparents is \textbf{fading}.

The Ministry of Public Health has \textbf{launched} a national \textbf{awareness campaign} to \textbf{raise awareness} on these scourges. Radio spots in Pidgin English \textbf{warn} youths: "\textbf{Don't die of} ignorance." Schools \textbf{advise} students to \textbf{lead a healthy lifestyle}. To \textbf{keep informed} is the first step to survival. Say no to drugs, limit alcohol, quit tobacco — choose life.
\end{textebox}

\begin{center}
\begin{tabularx}{\linewidth}{|l|X|}
\hline
\rowcolor{popblueL} \textbf{Mot / Expression} & \textbf{Glossaire (Nature, Sens)} \\
\hline
\textbf{bustling} (adj.) & animé, bouillonnant d'activité \\
\textbf{alarmingly} (adv.) & de façon alarmante \\
\textbf{gateway drug} (n.) & drogue d'entrée (qui mène à des drogues plus dures) \\
\textbf{vulnerable} (adj.) & vulnérable, fragile \\
\textbf{habit} (n.) & habitude (ici : dépendance, « la came ») \\
\textbf{devastating} (adj.) & dévastateur \\
\textbf{claim lives} (loc.) & faire des victimes, coûter la vie à \\
\textbf{graphic warnings} (n.) & avertissements chocs (images choquantes) \\
\textbf{endangers} (v.) & met en danger \\
\textbf{cope} (v.) & faire face, s'en sortir \\
\textbf{fading} (part.) & qui s'efface, qui disparaît \\
\textbf{scourges} (n.) & fléaux \\
\textbf{spots} (n.) & spots publicitaires, messages courts radio/TV \\
\textbf{launched} (v.) & lancé (prétérit de \eng{launch}) \\
\hline
\end{tabularx}
\end{center}

\begin{exoresolu}[Entraînement : Complétion et Reformulation]
\textbf{Énoncé :} Complétez les phrases avec le mot juste du texte ou du vocabulaire étudié (forme correcte). Puis traduisez la phrase 3.

1. Despite the \_\_\_\_\_\_\_\_\_\_ on cigarette packets, many teens start \_\_\_\_\_\_\_\_\_\_. (graphic warnings / smoking)
2. Tramadol \_\_\_\_\_\_\_\_\_\_ is a public health emergency in the North West Region. (addiction)
3. The NGO aims to \_\_\_\_\_\_\_\_\_\_ \_\_\_\_\_\_\_\_\_\_ on the dangers of \_\_\_\_\_\_\_\_\_\_ driving. (raise / awareness / drink-driving)
4. My grandfather suffers from severe \_\_\_\_\_\_\_\_\_\_ \_\_\_\_\_\_\_\_\_\_; he doesn't recognize me anymore. (memory loss)
5. She has been \_\_\_\_\_\_\_\_\_\_ \_\_\_\_\_\_\_\_\_\_ the wagon for six months now. (on / the) — \emph{Elle est sobre depuis six mois.}

\textbf{Solution détaillée :}
1. \textbf{warnings} (nom pluriel après \eng{graphic}) / \textbf{smoking} (gérondif après \eng{start} + V-ing ou nom \eng{smoking}). \textit{Traduction : Malgré les avertissements chocs sur les paquets de cigarettes, beaucoup d'adolescents commencent à fumer.}
2. \textbf{addiction} (nom, sujet de \eng{is}). \textit{La dépendance au tramadol est une urgence de santé publique...}
3. \textbf{raise awareness} (collocation figée) / \textbf{drink-driving} (adjectif composé invariable). \textit{L'ONG vise à sensibiliser sur les dangers de la conduite en état d'ivresse.}
4. \textbf{memory loss} (nom composé). \textit{Mon grand-père souffre d'une perte de mémoire sévère ; il ne me reconnaît plus.}
5. \textbf{on the wagon} (idiome figé). \textit{Elle ne boit plus (est en sevrage) depuis six mois.}
\end{exoresolu}

\begin{savaistu}[Culture : Sensibilisation en Afrique anglophone]
Au \textbf{Royaume-Uni} et aux \textbf{États-Unis}, le « \textit{plain packaging} » (paquet neutre) est obligatoire : plus de logos, juste la marque en police standard et une photo choc (poumon noir, dent pourrie) avec l'inscription \eng{"Smoking kills"} ou \eng{"Smoking causes fatal lung cancer"}.
Au \textbf{Nigeria} et au \textbf{Ghana}, des ONG et les ministères de la Santé diffusent des spots radio en \textbf{Pidgin English} (ex. : \eng{"Say no to tramadol, e go spoil your life!"}) pour toucher la jeunesse des rues. Au \textbf{Cameroun anglophone} (Buea, Bamenda, Limbe), des clubs de santé scolaires (\eng{Health Clubs}) organisent des \eng{awareness walks} (marches de sensibilisation) le 26 juin (Journée mondiale contre la drogue) et le 31 mai (Journée mondiale sans tabac). Participer à ces activités, c'est aussi \eng{keep informed} et \eng{raise awareness} !
\end{savaistu}

\courssub{Méthode d'écoute en 3 étapes et exemple traité}

\begin{methode}[Réussir une épreuve de compréhension orale (Listening)]
Suivez ces trois étapes à chaque écoute (examen ou entraînement) :
\begin{enumerate}
    \item \textbf{ANTICIPER (Before listening — 1ère lecture des questions)} : Lisez le titre, les consignes, les questions. Soulignez les \cle{mots-clés} (noms propres, chiffres, verbes d'action, négations). Prédisez le thème, le type de document (interview, reportage, dialogue), le registre (formel/informel). \textit{Astuce : Transformez les questions en "mots-attendus" (ex. "Why?" $\rightarrow$ écoutez "because, since, as").}
    \item \textbf{ÉCOUTER ACTIVEMENT (During listening — 2 écoutes)} :
        \begin{itemize}
            \item \textbf{1ère écoute (Globale) :} Ne notez rien. Identifiez : Qui parle ? De quoi ? Ton général ? Répondez aux questions "Global Understanding" (Vrai/Faux).
            \item \textbf{2ème écoute (Détaillée) :} Notez des \cle{mots-clés} (chiffres, noms, mots de liaison : \eng{however, therefore, but}). Complétez les réponses courtes, QCM. Vérifiez les pièges (négation, modalité : \eng{might, could} $\neq$ certitude).
        \end{itemize}
    \item \textbf{VÉRIFIER ET TRANSFÉRER (After listening)} : Relisez vos réponses. Vérifiez l'orthographe des noms propres et chiffres. Assurez-vous que la réponse correspond \textbf{exactement} à la question (unité, temps verbal). Transférez proprement sur la copie d'examen.
\end{enumerate}
\end{methode}

\begin{exemplebox}[Exemple traité : Question "Global Understanding" sur Text 1]
\textbf{Question :} \textit{The speaker says he started taking drugs because of peer pressure. True or False?}

\textbf{Analyse (Étape 1 - Anticipation) :} Mots-clés : "started", "drugs", "peer pressure" (pression des pairs), "True/False".
\textbf{Écoute (Étape 2) :} Extrait entendu : \eng{"Pressure was high. A senior student offered me cannabis to relax."}
\textbf{Raisonnement (Étape 3) :} Le texte mentionne "Pressure was high" (pression académique) et "A senior student offered me..." (offre d'un aîné). Ce n'est pas strictement de la "peer pressure" (pression d'amis du même âge), mais une offre d'un aîné dans un contexte de stress. Souvent au Bac, la nuance compte. Si la question dit "peer pressure", c'est souvent \textbf{False} ou \textbf{Not Mentioned} si le texte dit "academic pressure".
\textbf{Réponse :} \textbf{False} (ou \textbf{Not Mentioned} selon la rigueur attendue). \textit{Justification : He started due to academic pressure and an offer from a senior student, not peer pressure per se.}
\end{exemplebox}

\courssub{Production guidée : message de sensibilisation}

\begin{experience}[Production écrite/orale : Créez votre campagne "Choose Life"]
\textbf{Contexte :} Le Ministère de la Santé Publique lance un concours national. Vous devez créer un \textbf{spot radio de 30 secondes} (environ 60-75 mots) OU une \textbf{affiche de sensibilisation} (slogan + 4 arguments visuels/textuels) sur l'un des thèmes : \eng{Drugs}, \eng{Alcohol}, \eng{Tobacco}, \eng{Elderly Care}.

\textbf{Structure type d'un spot radio (Modèle) :}
\begin{enumerate}
    \item \textbf{Hook (Accroche) :} Son, choc, question, statistique. \eng{"Did you know? 13\% of Cameroonian teens smoke!"}
    \item \textbf{Facts (Faits) :} 2 dangers concrets (vocabulaire précis). \eng{"Tobacco causes lung cancer. Nicotine traps you."}
    \item \textbf{Consequences (Conséquences) :} Vie gâchée, famille, argent. \eng{"You lose your health, your money, your future."}
    \item \textbf{Call to Action (Appel à l'action) :} Verbe d'action + solution. \eng{"Quit smoking today. Call 119. Choose life."}
    \item \textbf{Slogan (Mémorable) :} Court, percutant. \eng{"Your health. Your choice. Say NO to tobacco."}
\end{enumerate}

\textbf{Boîte à outils linguistiques (Useful Phrases) :}
\begin{center}
\begin{tabularx}{\linewidth}{|X|X|}
\hline
\rowcolor{popblueL} \textbf{Fonction} & \textbf{Expressions utiles} \\
\hline
Alerter / Chocher & \eng{Warning! / Beware of... / Don't risk your life... / Stop!} \\
Donner des faits & \eng{... causes ... / ... leads to ... / ... increases the risk of ... / Statistics show that ...} \\
Conseiller / Ordonner & \eng{You should / must / have to ... / It's time to ... / Why not ...? / Choose to ...} \\
Sensibiliser & \eng{Raise awareness / Spread the word / Talk about it / Break the silence.} \\
\hline
\end{tabularx}
\end{center}

\textbf{Consignes :}
\begin{enumerate}
    \item Choisissez votre thème et votre format (Radio ou Affiche).
    \item Rédigez un brouillon en utilisant \textbf{au moins 10 mots} du lexique de la leçon (surlignez-les).
    \item Vérifiez : \eng{die of} (maladie), \eng{addicted to} + V-ing, \eng{advise/warn sb to/against}, collocations (\eng{raise awareness, lead a healthy life}).
    \item \textbf{Oral :} Enregistrez-vous (application OnBuch+ ou téléphone). Évaluez : prononciation, débit, intonation expressive.
    \item \textbf{Écrit :} Présentez l'affiche en classe. Votez pour la plus percutante.
\end{enumerate}
\end{experience}

\begin{exercice}[Entraînement personnel : Vocabulaire en action — Version Bac]
\textbf{1. Traduisez en anglais (soignez les structures) :}
a) Mon oncle est mort d'un cancer du poumon l'année dernière.
b) Le gouvernement a lancé une campagne de sensibilisation contre la drogue.
c) Elle est dépendante de l'alcool depuis son divorce.
d) Il a arrêté de fumer il y a deux ans.
e) Les personnes âgées ont besoin qu'on prenne soin d'elles.

\textbf{2. Choisissez la bonne préposition :}
a) He died \_\_\_ (of / from) a heart attack.
b) She is addicted \_\_\_ (to / on) tramadol.
c) They warned us \_\_\_ (against / about) the dangers of alcohol.
d) The doctor advised him \_\_\_ (to / for) quit smoking.

\textbf{3. Dérivation : Complétez le tableau.}
\begin{tabular}{|l|l|l|l|}
\hline
\rowcolor{popblueL} Verbe / Base & Nom (Personne) & Nom (Concept) & Adjectif \\
\hline
\eng{prevent} & & \eng{prevention} & \\
\hline
& \eng{smoker} & & \\
\hline
\eng{drink} & & & \eng{drunk} \\
\hline
& \eng{addict} & \eng{addiction} & \eng{addictive} \\
\hline
\eng{age} & & & \\
\hline
\end{tabular}

\textbf{4. Expression écrite (Type Bac — 150 mots) :}
\textit{"You are the president of your school Health Club. Write an article for the school magazine entitled 'Youth and Addictions: Breaking the Chains'. Warn your peers about the dangers of drugs, alcohol and tobacco. Advise them on how to stay healthy and where to seek help."}
\end{exercice}

\begin{corrige}[Exercice : Vocabulaire en action]
\textbf{1. Traduction :}
a) My uncle \textbf{died of lung cancer} last year. (\textit{Rappel : \eng{die of} pour maladie}).
b) The government \textbf{launched an awareness campaign} against drugs. (\textit{Collocations : \eng{launch a campaign} + \eng{awareness campaign}}).
c) She \textbf{has been addicted to alcohol} since her divorce. (\textit{Present perfect + \eng{addicted to} + nom}).
d) He \textbf{quit smoking} / \textbf{gave up smoking} two years ago. (\textit{Prétérit + gérondif après \eng{quit/give up}}).
e) \textbf{Elderly people need to be taken care of} / \textbf{We need to take care of the elderly}. (\textit{Passif ou actif ; \eng{elderly people} préféré à \eng{old people}}).

\textbf{2. Prépositions :}
a) \textbf{of} (maladie/cause interne).
b) \textbf{to} (\eng{addicted to}).
c) \textbf{against} (\eng{warn sb against sth} = mettre en garde contre un danger).
d) \textbf{to} (\eng{advise sb to do sth}).

\textbf{3. Tableau de dérivation :}
\begin{tabular}{|l|l|l|l|}
\hline
\rowcolor{popblueL} Verbe / Base & Nom (Personne) & Nom (Concept) & Adjectif \\
\hline
\eng{prevent} & — & \eng{prevention} & \eng{preventive} / \eng{preventable} \\
\hline
\eng{smoke} & \eng{smoker} & \eng{smoking} & \eng{smoky} \\
\hline
\eng{drink} & \eng{drinker} & \eng{drinking} & \eng{drunk} / \eng{drunken} \\
\hline
\eng{addict} & \eng{addict} & \eng{addiction} & \eng{addictive} / \eng{addicted} \\
\hline
\eng{age} & — & \eng{ageing} & \eng{aged} / \eng{ageing} \\
\hline
\end{tabular}

\textbf{4. Proposition de rédaction (Article) :}
\textbf{Youth and Addictions: Breaking the Chains}

By [Your Name], Health Club President

Dear fellow students,

Did you know that tramadol addiction is destroying the future of many young Cameroonians? Today, I want to \textbf{raise awareness} on a silent epidemic: substance abuse.

Many of us start with \textbf{cannabis} or \textbf{alcohol} just to "fit in" or escape exam pressure. But the \textbf{gateway drug} effect is real. Soon, \textbf{hard drugs} like tramadol take control. You become an \textbf{addict}, spending your school fees on your \textbf{habit}. The risk of \textbf{overdose} is daily. \textbf{Drink-driving} kills on our roads. \textbf{Smoking} causes \textbf{lung cancer} and hurts your friends through \textbf{second-hand smoke}.

\eng{Prevention is better than cure.} \textbf{Lead a healthy lifestyle}: do sport, eat well, \textbf{keep informed} on health risks. If you are trapped, \textbf{seek help}. Talk to a teacher, a counsellor, or call the \textbf{helpline 119}. You can \textbf{quit smoking}, \textbf{give up drugs}, \textbf{get on the wagon}. \textbf{Rehabilitation} works — Elias from Buea is living proof.

Don't \textbf{die of} ignorance. Break the chains. Choose life.

\textit{Word count: ~160 words. Vocabulaire ciblé surligné. Structures : \eng{die of}, \eng{addicted to}, \eng{lead a healthy lifestyle}, \eng{raise awareness}, collocations.}
\end{corrige}

\coursec{Méthode d'écoute en 3 étapes et exemple traité}

Dans la section précédente, nous avons construit le vocabulaire des addictions et du vieillissement (\eng{addiction, tobacco, alcohol abuse, ageing, to keep informed about}...). Ce lexique ne sert pleinement que si vous savez le \emph{reconnaître à l'oreille} et l'exploiter rapidement. C'est exactement l'objet de cette section : transformer votre vocabulaire en \cle{stratégie d'écoute} efficace, celle que les correcteurs du Baccalauréat attendent de vous.

\courssub{Observer d'abord : un extrait de script d'écoute}

Avant toute règle, observons un court extrait de script, du type de celui que vous entendrez en classe ou à l'examen. Lisez-le une première fois, puis demandez-vous : \emph{quelles informations un élève attentif doit-il absolument relever ?}

\begin{textebox}[Listening script — Health awareness on CRTV]
“Good morning, dear listeners. This is CRTV Yaoundé, and today our guest is Dr. Ngono, a specialist in public health at Yaoundé Central Hospital. Doctor, our topic is keeping informed about health problems related to drugs, tobacco and alcohol abuse.”

“Thank you. Well, the figures are worrying. According to a recent survey in our region, about \textbf{30 per cent} of young adults between eighteen and twenty-five smoke regularly. That is almost one young person in three. Alcohol consumption has doubled in ten years, especially among teenagers aged \textbf{thirteen} to seventeen.”

“What can families do, Doctor?”

“First, keep informed about your health and about the dangers of these products. Second, talk openly with your children. And finally, remember that ageing well also means avoiding tobacco and alcohol early in life. Prevention is always better than cure.”

“Thank you, Dr. Ngono. That was our health programme for this morning.”
\end{textebox}

\textbf{Glossaire :} \eng{survey} (n.) : enquête ; \eng{worrying} (adj.) : inquiétant ; \eng{aged} (adj.) : âgé de ; \eng{ageing} (n.) : vieillissement ; \eng{prevention} (n.) : prévention ; \eng{cure} (n.) : remède, traitement.

\textbf{Observation.} Un élève qui écoute ce script sans méthode retient « un docteur parle de tabac ». Un élève méthodique relève : un \cle{nom propre} (Dr. Ngono, Yaoundé Central Hospital, CRTV), des \cle{chiffres} (30 per cent ; thirteen to seventeen), et un \cle{conseil-clé} (\eng{keep informed about your health}). Toute la méthode qui suit vise à faire de vous ce second élève.

\courssub{La règle des 3 P : Predict, Pick out, Prove}

\begin{methode}[La règle des 3 P — Predict, Pick out, Prove]
Toute écoute réussie suit trois étapes, résumées par trois verbes anglais commençant par P :
\begin{enumerate}
\item \textbf{Predict} (prédire) — \emph{avant} l'écoute : lisez toutes les questions, identifiez le thème annoncé, et activez le vocabulaire que vous connaissez déjà sur ce thème. Demandez-vous : \eng{What kind of words will I probably hear?}
\item \textbf{Pick out} (repérer, extraire) — \emph{pendant} l'écoute : notez uniquement les mots-clés, les chiffres et les noms propres, sous forme abrégée. Ne vous arrêtez jamais sur un mot inconnu : le sens global suffit souvent pour répondre.
\item \textbf{Prove} (prouver, justifier) — \emph{après} l'écoute : rédigez vos réponses en phrases complètes, vérifiez leur cohérence avec le sens général, et justifiez chaque réponse par un élément précis du script (\eng{The speaker says that...}).
\end{enumerate}
\end{methode}

\begin{popfigure}
\begin{tikzpicture}[
box/.style={rectangle, rounded corners=4pt, draw=popblue, thick, fill=popblueL, text width=5.2cm, align=center, minimum height=1.5cm, font=\small},
arr/.style={-{Stealth[length=3mm]}, thick, popdark}]
\node[box, align=center] (before) at (0,0){\textbf{BEFORE — Predict}\\ lire les questions, prédire le thème, activer le vocabulaire};
\node[box, fill=poporangeL, draw=poporange, align=center] (during) at (0,-2.4){\textbf{DURING — Pick out}\\ noter mots-clés, chiffres (en chiffres arabes), noms propres};
\node[box, fill=popgreenL, draw=popgreen, align=center] (after) at (0,-4.8){\textbf{AFTER — Prove}\\ répondre en phrases complètes, vérifier, justifier par le texte};
\draw[arr] (before) -- (during);
\draw[arr] (during) -- (after);
\end{tikzpicture}
\legende{Organigramme de la méthode d'écoute en 3 étapes (les 3 P)}
\end{popfigure}

\courssub{Étape 1 — Avant l'écoute : Predict}

Lisez les questions \emph{avant} que l'enregistrement ne commence. Chaque question est un indice sur le contenu du texte.

\begin{exemplebox}[Prédire à partir des questions]
Supposons que vous lisiez ces deux questions avant l'écoute :
\begin{itemize}
\item \eng{What percentage of young adults smoke regularly?} — le mot \eng{percentage} vous annonce qu'il faudra écouter un \cle{chiffre} suivi de \eng{per cent}.
\item \eng{Where does Dr. Ngono work?} — le mot \eng{where} vous annonce qu'il faudra écouter un \cle{lieu}, probablement un hôpital ou une ville.
\end{itemize}
Vous pouvez aussi prédire le vocabulaire : le thème « health problems » évoque \eng{smoke, tobacco, alcohol, drugs, hospital, doctor, survey, per cent}. Ainsi préparé, votre oreille reconnaîtra ces mots instantanément.
\end{exemplebox}

\begin{propriete}[Règle — l'activation du vocabulaire]
Avant l'écoute, listez mentalement (ou au brouillon) cinq à dix mots anglais liés au thème annoncé. Cette \cle{activation} réduit fortement le temps de reconnaissance pendant l'écoute. En français, on fait la même chose naturellement : si l'on vous dit « on va parler d'hôpital », vous vous attendez à entendre « médecin, patient, traitement ». La stratégie est identique en anglais — il faut simplement la faire \emph{consciemment}.
\end{propriete}

\courssub{Étape 2 — Pendant l'écoute : Pick out}

Pendant l'écoute, votre brouillon doit ressembler à un télégramme, pas à une rédaction. Notez :
\begin{itemize}
\item les \cle{chiffres} — \textbf{toujours en chiffres arabes} (30, 13, 18--25), jamais en lettres : écrire « thirty » prend du temps et vous fait perdre la suite ;
\item les \cle{noms propres} — avec la première lettre en majuscule, ce qui les repère visuellement : \eng{Dr. Ngono, CRTV, Yaoundé} ;
\item les \cle{mots-clés} du thème — un ou deux mots par idée : \eng{smoke, alcohol doubled, prevention}.
\end{itemize}

\begin{attention}[Piège majeur : thirteen (13) ou thirty (30) ?]
C'est l'erreur d'écoute la plus fréquente chez les francophones. La différence ne tient pas seulement à la fin du mot, mais à l'\cle{accent tonique} :
\begin{itemize}
\item nombres en \textbf{-teen} : l'accent tombe sur la \textbf{deuxième} syllabe — \eng{thirTEEN} (13), \eng{fourTEEN} (14), \eng{fifTEEN} (15) ;
\item nombres en \textbf{-ty} : l'accent tombe sur la \textbf{première} syllabe — \eng{THIRty} (30), \eng{FORty} (40), \eng{FIFty} (50).
\end{itemize}
Prononciation approchée : \eng{thirteen} « seur-TINE », \eng{thirty} « SEUR-ti ». Dans notre script, \eng{30 per cent of young adults} concerne les jeunes adultes, tandis que \eng{teenagers aged thirteen to seventeen} concerne les adolescents : confondre les deux chiffres fausse deux réponses d'un coup. Entraînez-vous à répéter les paires à voix haute en marquant l'accent.
\end{attention}

\begin{aretenir}[Les deux réflexes du brouillon]
\begin{enumerate}
\item Chiffre entendu $\rightarrow$ écrire immédiatement en chiffres arabes : \eng{thirteen} $\rightarrow$ 13 ; \eng{thirty per cent} $\rightarrow$ 30\%.
\item Nom propre entendu $\rightarrow$ noter avec majuscule : \eng{Dr. Ngono}, \eng{CRTV}, \eng{Yaoundé Central Hospital}. La majuscule vous permet de retrouver l'information d'un coup d'œil au moment de rédiger.
\end{enumerate}
\end{aretenir}

\courssub{Étape 3 — Après l'écoute : Prove}

L'écoute terminée, il reste à transformer vos notes en réponses. Trois exigences :
\begin{enumerate}
\item \textbf{Répondre en phrases complètes.} À la question \eng{What percentage of young adults smoke regularly?}, n'écrivez pas seulement « 30\% », mais : \eng{According to the survey, 30 per cent of young adults smoke regularly.} (« Selon l'enquête, 30 \% des jeunes adultes fument régulièrement. »)
\item \textbf{Vérifier la cohérence.} Si votre réponse contredit le sens général du texte (par exemple « 13\% of young adults smoke » alors que le chiffre 13 concernait l'\emph{âge} des adolescents), c'est le signal d'une erreur d'écoute.
\item \textbf{Justifier par le texte.} Pour les Vrai/Faux, appuyez-vous toujours sur une formulation du script : \eng{The speaker says that...} / \eng{According to Dr. Ngono,...}
\end{enumerate}

\begin{center}
\begin{tabularx}{\linewidth}{|X|X|X|}
\hline
\rowcolor{popblueL} Étape & En anglais & Ce que fait l'élève \\
\hline
Avant l'écoute & Predict & lit les questions, prédit le thème, active le vocabulaire \\
\hline
Pendant l'écoute & Pick out & note chiffres (en chiffres arabes), noms propres (majuscules), mots-clés \\
\hline
Après l'écoute & Prove & rédige en phrases complètes, vérifie la cohérence, justifie par le script \\
\hline
\end{tabularx}
\end{center}

\courssub{Exemples traités pas à pas}

\begin{exoresolu}[Exemple 1 — Question à réponse courte]
Énoncé : \eng{What percentage of young adults smoke regularly?} Répondez en une phrase complète.
\tcblower
\textbf{Solution pas à pas.}
\begin{enumerate}
\item \textbf{Predict} : le mot \eng{percentage} annonce un chiffre suivi de \eng{per cent} ; je tends l'oreille vers tous les nombres.
\item \textbf{Pick out} : j'entends \eng{about 30 per cent of young adults between eighteen and twenty-five smoke regularly} $\rightarrow$ je note au brouillon : « 30\% — young adults — smoke ».
\item \textbf{Prove} : je rédige : \eng{According to the survey mentioned by Dr. Ngono, 30 per cent of young adults smoke regularly.} (« Selon l'enquête mentionnée par le Dr Ngono, 30 \% des jeunes adultes fument régulièrement. ») Je vérifie : le chiffre 30 concerne bien les jeunes adultes, pas les adolescents — cohérent.
\end{enumerate}
\end{exoresolu}

\begin{exoresolu}[Exemple 2 — Vrai ou Faux, avec justification]
Énoncé : \eng{Dr. Ngono works at Douala General Hospital.} — True or False? Justify your answer.
\tcblower
\textbf{Solution pas à pas.}
\begin{enumerate}
\item \textbf{Predict} : la question porte sur un \cle{lieu de travail} ; j'écoute les noms de villes et d'hôpitaux.
\item \textbf{Pick out} : j'entends \eng{a specialist in public health at Yaoundé Central Hospital} $\rightarrow$ je note : « Dr. Ngono — Yaoundé Central Hospital ».
\item \textbf{Prove} : je rédige : \eng{False. Dr. Ngono does not work at Douala General Hospital; the speaker says he works at Yaoundé Central Hospital.} (« Faux. Le Dr Ngono ne travaille pas à l'Hôpital Général de Douala ; le présentateur dit qu'il travaille à l'Hôpital Central de Yaoundé. »)
\end{enumerate}
Remarque : au Bac, un « False » sans justification ne vaut souvent que la moitié des points. Justifiez toujours.
\end{exoresolu}

\courssub{Réemployer le lexique : exemples traduits}

\begin{exemplebox}[Phrases-modèles à réutiliser dans vos réponses]
\eng{Alcohol consumption has doubled in ten years.} — « La consommation d'alcool a doublé en dix ans. » (Notez le present perfect : bilan d'une évolution jusqu'à aujourd'hui.)

\eng{Keep informed about your health.} — « Restez informé(e) sur votre santé. » (Impératif ; attention à la préposition : \eng{informed \textbf{about}}, jamais « informed of your health » dans ce sens.)

\eng{Prevention is always better than cure.} — « Mieux vaut prévenir que guérir. » (Proverbe anglais : l'ordre des mots est inversé par rapport au français.)

\eng{That is almost one young person in three.} — « Cela représente presque un jeune sur trois. » (Structure des proportions : \eng{one in three}, « un sur trois ».)
\end{exemplebox}

\begin{attention}[Erreurs fréquentes des francophones dans ce thème]
\begin{itemize}
\item \eng{to keep informed} ne se traduit pas par « garder informé » au sens de retenir : c'est « se tenir informé, rester informé ».
\item Ne dites pas « the \emph{actual} figures » pour « les chiffres actuels » : \eng{actual} signifie « réel, véritable » (faux ami). Dites \eng{the current figures}.
\item \eng{a survey} (« une enquête », prononcé « SEUR-véï ») ne se confond pas avec \eng{to survey} (« enquêter sur », accent sur la 2e syllabe).
\item « Fumer » se dit \eng{to smoke} ; « un fumeur » est \eng{a smoker}, jamais « a smoking man ».
\end{itemize}
\end{attention}

\begin{savaistu}[Le double passage à l'écoute]
Au Baccalauréat, le texte de compréhension orale est le plus souvent lu \textbf{deux fois}. Ce n'est pas un luxe, c'est une stratégie à exploiter méthodiquement : utilisez la \textbf{première écoute} pour saisir le sens global (qui parle ? de quoi ? quel ton ?) et noter les grandes informations ; réservez la \textbf{seconde écoute} aux détails (chiffres exacts, noms propres) et à la vérification de vos réponses. Entre les deux passages, relisez vos notes et repérez les questions encore sans réponse : vous saurez ainsi exactement quoi écouter la seconde fois. Les candidats qui paniquent à la première écoute parce qu'ils n'ont pas tout compris oublient que le second passage existe précisément pour cela.
\end{savaistu}

\begin{exercice}[À vous de jouer]
En vous aidant du script de la section « Observer d'abord », répondez en phrases complètes :
\begin{enumerate}
\item \eng{How many pieces of advice does Dr. Ngono give to families? Name one.}
\item \eng{Alcohol consumption has decreased in ten years.} — True or False? Justify.
\item \eng{What age group of teenagers is mentioned in the survey?}
\end{enumerate}
\end{exercice}

\begin{corrige}[Exercice — À vous de jouer]
\begin{enumerate}
\item \eng{Dr. Ngono gives three pieces of advice to families. For example, he says that parents should talk openly with their children.} (Les trois conseils : \eng{keep informed}, \eng{talk openly}, \eng{avoid tobacco and alcohol early in life}.)
\item \eng{False. The speaker says that alcohol consumption has doubled in ten years, not decreased.}
\item \eng{The survey mentions teenagers aged thirteen to seventeen.} (Attention : ici, il s'agit bien de \eng{thirteen} — \eng{thirTEEN} — et non de \eng{thirty}.)
\end{enumerate}
\end{corrige}

\begin{aretenir}[L'essentiel de la méthode]
Une écoute réussie repose sur les \cle{3 P} : \textbf{Predict} avant (questions lues, thème prédit, vocabulaire activé), \textbf{Pick out} pendant (chiffres en chiffres arabes, noms propres en majuscules, mots-clés seulement), \textbf{Prove} après (phrases complètes, cohérence vérifiée, justification par le script). Et n'oubliez jamais : \eng{thirTEEN} (13) s'oppose à \eng{THIRty} (30) par l'accent tonique.
\end{aretenir}

\coursec{Production guidée : message de sensibilisation}

\begin{experience}[Mise en route : Anticipation et rappel]
En classe entière, l'enseignant projette ou écrit au tableau trois images choc : un poumon noirci par le tabac, un foie cirrhotique, une personne âgée isolée. 
\begin{enumerate}
    \item \textbf{Brainstorming éclair (2 min) :} Chaque élève note 3 mots-clés en anglais qui lui viennent à l'esprit. Mise en commun au tableau : construction d'un nuage de mots (\textit{word cloud}) : \cle{addiction}, \cle{lungs}, \cle{liver}, \cle{elderly}, \cle{abuse}, \cle{prevention}...
    \item \textbf{Rappel des stratégies d'écoute :} « Avant d'écouter ou de produire, on identifie le \textbf{thème}, on anticipe le \textbf{lexique} et on repère les \textbf{chiffres/noms propres}. »
\end{enumerate}
\end{experience}

\courssub{Analyse d'un modèle de message radio}

Avant de rédiger le vôtre, étudions un script modèle respectant exactement le plan imposé (accroche $\rightarrow$ dangers $\rightarrow$ conseils $\rightarrow$ slogan) et intégrant le lexique cible.

\begin{textebox}{Modèle de script radio : « Tobacco Free Youth » (Durée : env. 45 secondes)}
\textbf{Host:} Good morning, listeners! Did you know that \textbf{tobacco} kills more than 8 million people every year worldwide according to the WHO? That is one death every four seconds.

\textbf{Host:} In Cameroon, many young people start smoking \emph{shisha} or cigarettes thinking it is trendy. But the \textbf{dangers} are real. First, the toxic smoke destroys your \textbf{lungs}, causing chronic bronchitis and cancer. Second, nicotine creates a powerful \textbf{addiction} that controls your brain and your wallet.

\textbf{Host:} So, what can you do? You \textbf{must} say no to the first cigarette. You \textbf{should} avoid places where people smoke to protect yourself from passive smoking. You \textbf{ought to} call a friend or do sport if you feel the urge.

\textbf{Host (energetic):} Remember: \textbf{Your health is your wealth!} Say no to tobacco, say yes to life! This is a message from the Ministry of Public Health.
\end{textebox}

\begin{definition}[Glossaire du modèle]
\begin{itemize}
    \item \textbf{listeners} (n. pl.) : auditeurs / auditrices.
    \item \textbf{trendy} (adj.) : à la mode, branché.
    \item \textbf{shisha} (n.) : chicha, narguilé.
    \item \textbf{chronic bronchitis} (n.) : bronchite chronique.
    \item \textbf{passive smoking} (n.) : tabagisme passif (syn. \cle{secondhand smoke}).
    \item \textbf{urge} (n.) : envie pressante, pulsion.
    \item \textbf{wealth} (n.) : richesse, bien-être.
    \item \textbf{wallet} (n.) : portefeuille.
\end{itemize}
\end{definition}

\begin{exercice}[Questions de repérage sur le modèle]
\begin{enumerate}
    \item Identifiez l'\textbf{accroche} (chiffre + source).
    \item Relevez les \textbf{deux dangers} cités (organe + conséquence).
    \item Surlignez les \textbf{trois structures de conseil} utilisées (\textit{must, should, ought to}).
    \item Quel est le \textbf{slogan final} ? Est-il court, rythmé, mémorisable ?
\end{enumerate}
\end{exercice}

\begin{corrige}[Exercice : Questions de repérage]
\begin{enumerate}
    \item \eng{« Did you know that tobacco kills more than 8 million people every year worldwide according to the WHO? »} (Chiffre : 8 millions ; Source : OMS).
    \item 1. Destruction des \textbf{poumons (lungs)} $\rightarrow$ bronchite chronique, cancer. 2. \textbf{Addiction (addiction)} à la nicotine $\rightarrow$ contrôle du cerveau et du portefeuille.
    \item \textbf{Must} say no (obligation forte) ; \textbf{Should} avoid places (conseil) ; \textbf{Ought to} call a friend (conseil moral / alternative).
    \item \eng{« Your health is your wealth! Say no to tobacco, say yes to life! »} Court, rime (health/wealth), parallélisme (no/yes, tobacco/life).
\end{enumerate}
\end{corrige}

\courssub{Structure imposée du message}

Votre message doit suivre \textbf{impérativement} ce schéma en 4 étapes (6 à 8 lignes au total). Chaque étape correspond à une fonction communicationnelle précise.

\begin{popfigure}
\begin{tikzpicture}[
    node distance=1.2cm and 1.5cm,
    box/.style={rectangle, rounded corners, draw=popblue, thick, fill=popblueL, text width=3.2cm, align=center, font=\small},
    arrow/.style={-Stealth, thick, popdark}
]
\node[box, align=center] (hook){\textbf{1. ACCROCHE (Hook)}\\\textit{Question ou Chiffre choc}\\\eng{Did you know...? / 1 in 3...}};
\node[box, right=of hook, align=center] (dangers){\textbf{2. DANGERS (Faits)}\\\textit{2 faits précis : organe, statistique}\\\eng{Destroys lungs / Causes addiction}};
\node[box, right=of dangers, align=center] (advice){\textbf{3. CONSEILS (Advice)}\\\textit{2 impératifs / modaux}\\\eng{Must / Should / Ought to + BV}};
\node[box, right=of advice, align=center] (slogan){\textbf{4. SLOGAN FINAL}\\\textit{Court, positif, mémorable}\\\eng{Health is wealth!}};

\draw[arrow] (hook.east) -- (dangers.west);
\draw[arrow] (dangers.east) -- (advice.west);
\draw[arrow] (advice.east) -- (slogan.west);

% Annotations sous les boîtes
\node[below=0.3cm of hook, font=\tiny, text=popmuted] {Lignes 1--2};
\node[below=0.3cm of dangers, font=\tiny, text=popmuted] {Lignes 3--4};
\node[below=0.3cm of advice, font=\tiny, text=popmuted] {Lignes 5--6};
\node[below=0.3cm of slogan, font=\tiny, text=popmuted] {Lignes 7--8};
\end{tikzpicture}
\legende{Schéma de la structure obligatoire du message radio de sensibilisation}
\end{popfigure}

\begin{aretenir}[Plan type pour la rédaction (6--8 lignes)]
\begin{enumerate}
    \item \textbf{Lignes 1--2 : ACCROCHE.} Une question rhétorique (\eng{Did you know...?}) ou un chiffre choc avec source (\eng{According to the Ministry of Health...}).
    \item \textbf{Lignes 3--4 : DANGERS.} Deux faits concrets : nommer l'organe touché (\cle{lungs}, \cle{liver}, \cle{brain}, \cle{heart}) et la conséquence (cancer, cirrhose, \cle{addiction}, perte de mémoire). Utilisez le \textbf{Present Simple}.
    \item \textbf{Lignes 5--6 : CONSEILS.} Deux injonctions distinctes. Une avec \textbf{must/should/ought to} + base verbale. Une à l'\textbf{impératif} (\eng{Say no to...}, \eng{Avoid...}, \eng{Call...}).
    \item \textbf{Lignes 7--8 : SLOGAN.} Une phrase courte, percutante, si possible avec rime, allitération ou parallélisme.
\end{enumerate}
\end{aretenir}

\courssub{Outils linguistiques : Modaux de conseil et impératifs}

Pour exprimer la recommandation, l'obligation morale ou le conseil avisé, l'anglais utilise trois modaux principaux + l'impératif. Maîtrisez leurs nuances.

\begin{propriete}[Modaux de conseil : \texorpdfstring{\eng{should / must / ought to}}{should / must / ought to}]
\begin{center}
\begin{tabularx}{\linewidth}{|X|X|X|X|}
\hline
\rowcolor{popblueL}
\textbf{Modal} & \textbf{Force / Nuance} & \textbf{Forme} & \textbf{Exemple (Thème santé)} \\
\hline
\textbf{MUST} & Obligation forte, nécessité vitale, ordre médical. & \eng{Must + BV} (pas de 'to') & \eng{You \textbf{must} stop smoking immediately.} (Tu dois / Il faut absolument...) \\
\hline
\textbf{SHOULD} & Conseil avisé, recommandation standard, « il serait bon que ». & \eng{Should + BV} & \eng{You \textbf{should} see a doctor regularly.} (Tu devrais / Il faut que tu voies...) \\
\hline
\textbf{OUGHT TO} & Devoir moral, attente sociale, très proche de \eng{should} mais plus formel. & \eng{Ought to + BV} & \eng{We \textbf{ought to} protect the \textbf{elderly}} from neglect. (Nous devons / Il convient de...) \\
\hline
\textbf{IMPÉRATIF} & Ordre direct, slogan, instruction claire, urgence. & \eng{BV} (sujet implicite \eng{you}) & \eng{\textbf{Say no} to drugs! \textbf{Protect} your \textbf{liver}!} \\
\hline
\end{tabularx}
\end{center}

\textbf{Forme interrogative (inversion pour le conseil) :} \eng{\textbf{Should} I see a doctor?} (Fréquent). \eng{\textbf{Ought} I to go?} (Très formel, rare). 

\textbf{Négation :} \eng{must not / mustn't} (interdiction formelle), \eng{should not / shouldn't} (déconseillé), \eng{ought not to / oughtn't to} (rare à l'oral).
\eng{Don't} + BV pour l'impératif négatif : \eng{\textbf{Don't} start smoking.}
\end{propriete}

\begin{attention}[Piège majeur : Noms indénombrables \texorpdfstring{\eng{alcohol, tobacco, drugs}}{alcohol, tobacco, drugs}]
\eng{Alcohol}, \eng{tobacco} et \eng{drugs} (au sens général de « substances illicites ») sont des \textbf{noms indénombrables (uncountable)} en anglais.
\begin{itemize}
    \item \textbf{FAUX} : \eng{The alcohol is bad. / A tobacco kills. / The drugs are dangerous.} (si sens général).
    \item \textbf{VRAI} : \eng{\textbf{Alcohol} damages the \textbf{liver}.} (Zéro article). \eng{\textbf{Tobacco} causes cancer.} \eng{\textbf{Drug} abuse destroys lives.}
    \item \textbf{Exception} : \eng{A drug} (un médicament / une substance précise), \eng{drugs} (pluriel = plusieurs types de substances).
\end{itemize}
\emph{Astuce :} Si vous pouvez dire « de l'alcool / du tabac » (partitif) en français $\rightarrow$ pas de \eng{the} / \eng{a} en anglais.
\end{attention}

\begin{exemplebox}[Comparaison Français / Anglais : Article zéro]
\begin{center}
\begin{tabularx}{\linewidth}{|X|X|X|}
\hline
\rowcolor{popblueL}
\textbf{Français (Partitif)} & \textbf{Anglais (Zéro article)} & \textbf{Remarque} \\
\hline
L'alcool tue. / \textit{De l'}alcool tue. & \eng{Alcohol kills.} & Règle générale : vérité générale = zéro article. \\
\hline
Le tabac abîme les poumons. & \eng{Tobacco damages the lungs.} & \eng{The lungs} = parties du corps définies (possesseur implicite). \\
\hline
La drogue rend dépendant. & \eng{Drug addiction creates dependence.} & Préférer \eng{Drug abuse / addiction} (noms composés) à \eng{Drugs} seul. \\
\hline
Il boit \textit{de l'}alcool. & \eng{He drinks alcohol.} & Après verbe de consommation (\eng{drink, smoke, take}). \\
\hline
\end{tabularx}
\end{center}
\end{exemplebox}

\courssub{Lexique obligatoire à réemployer (Minimum 5 mots)}

Vous devez intégrer \textbf{au moins 5} de ces mots/expressions dans votre message. Cochez-les lors de la rédaction.

\begin{definition}[Banque lexicale thématique — Addictions \& Vieillissement]
\begin{center}
\begin{tabularx}{\linewidth}{|l|l|X|}
\hline
\rowcolor{popblueL}
\textbf{Mot / Expression} & \textbf{Nature} & \textbf{Traduction / Contexte} \\
\hline
\cle{abuse} & n. & abus (usage nocif) — \eng{drug abuse, alcohol abuse} \\
\cle{addiction} & n. & dépendance, accoutumance — \eng{nicotine addiction} \\
\cle{awareness} & n. & sensibilisation, prise de conscience — \eng{raise awareness} \\
\cle{elderly} & adj./n. & personnes âgées (politesse) — \eng{care for the elderly} \\
\cle{liver} & n. & foie — \eng{liver cirrhosis / damage} \\
\cle{lungs} & n. pl. & poumons — \eng{damage the lungs} \\
\cle{prevention} & n. & prévention — \eng{prevention campaign} \\
\cle{recovery} & n. & guérison, sevrage — \eng{long recovery} \\
\cle{secondhand smoke} & n. & fumée secondaire / tabagisme passif \\
\cle{withdrawal} & n. & sevrage, syndrome de manque — \eng{withdrawal symptoms} \\
\hline
\cle{isolation} & n. & isolement — \eng{social isolation} \\
\cle{loneliness} & n. & solitude (souffrance) — \eng{feel loneliness} \\
\cle{neglect} & n. & négligence, abandon — \eng{elderly neglect} \\
\cle{shisha} & n. & chicha, narguilé \\
\cle{trendy} & adj. & à la mode, branché \\
\hline
\end{tabularx}
\end{center}
\end{definition}

\courssub{Méthode de présentation orale (Simulation radio)}

\begin{methode}[Parler comme un pro à la radio : 4 étapes]
\begin{enumerate}
    \item \textbf{Préparez votre « fiche de lecture » :} Écrivez le texte en gros caractères, double interligne. \textbf{Surlignez} les mots-clés (chiffres, modaux, slogans) et \textbf{marquez les pauses} ( // ) après chaque idée.
    \item \textbf{Articulez les chiffres et noms propres :} \eng{Eight million} (pas « ète »), \eng{W-H-O} (épelé), \eng{Cameroon} (/ˌkæməˈruːn/ « ka-me-roun »).
    \item \textbf{Modulez votre voix :} Ton grave et lent pour les dangers (gravité) $\rightarrow$ Ton dynamique, plus rapide et souriant pour les conseils et le slogan (espoir, action).
    \item \textbf{Regardez le « micro » (caméra/classe) :} Ne lisez pas le nez sur la feuille. Levez les yeux à la fin de chaque phrase. Souriez pour le slogan final (ça s'entend !).
\end{enumerate}
\end{methode}

\courssub{Exemple traité : Rédaction collaborative (Thème : Prise en charge des personnes âgées)}

Voici comment un groupe de 4 élèves a construit son message en suivant la méthode.

\begin{exoresolu}[Rédaction pas à pas — Thème : \eng{Elderly Care}]
\textbf{Étape 1 : Choix du thème et brainstorming lexique.}
Mots choisis : \cle{elderly}, \cle{awareness}, \cle{neglect}, \cle{loneliness}, \cle{prevention}, \cle{isolation}.

\textbf{Étape 2 : Rédaction de l'accroche (Hook).}
\eng{Did you know that 1 in 5 elderly people in Cameroon suffer from loneliness?} (Chiffre inventé mais réaliste, source implicite. \textbf{Note :} \eng{1 in 5 people} $\rightarrow$ pluriel \eng{suffer}).

\textbf{Étape 3 : Rédaction des dangers (2 faits).}
\eng{First, social isolation leads to depression and memory loss. Second, neglect can cause malnutrition and untreated chronic diseases.}
(Vérification : \eng{isolation, neglect} = indénombrables, pas de \eng{the}. \eng{The elderly} = groupe défini, donc \eng{the}).

\textbf{Étape 4 : Rédaction des conseils (Modaux + Impératif).}
\eng{You \textbf{should} visit your grandparents every week.} (Conseil).
\eng{Communities \textbf{must} build day-care centres for the \textbf{elderly}.} (Obligation forte).
\eng{\textbf{Don't ignore} their silence.} (Impératif négatif).

\textbf{Étape 5 : Slogan final.}
\eng{Respect their past, secure their future!} (Parallélisme : respect/secure, past/future).

\textbf{Étape 6 : Assemblage et relecture (7 lignes).}
\tcblower
\begin{textebox}{Message final du groupe (Lu par l'élève « Speaker »)}
\textbf{Speaker:} Good morning! Did you know that 1 in 5 \textbf{elderly} people \textbf{suffer} from \textbf{loneliness}? First, social \textbf{isolation} leads to depression. Second, \textbf{neglect} causes malnutrition. You \textbf{should} visit your grandparents weekly. Communities \textbf{must} support \textbf{prevention} centres. \textbf{Don't ignore} their silence. Respect their past, secure their future!
\end{textebox}

\textbf{Analyse du correcteur :} 7 lignes. 7 mots du lexique (\cle{elderly} x2, \cle{loneliness}, \cle{isolation}, \cle{neglect}, \cle{prevention}, \cle{awareness} implicite). 3 structures de conseil (\eng{should, must, Don't}). Slogan percutant. \textbf{Note : 8/8}.
\end{exoresolu}

\courssub{À votre tour : Production en groupe}

\begin{exercice}[Tâche finale : « On Air : Health Alert »]
\textbf{Consignes :}
\begin{enumerate}
    \item \textbf{Groupes de 3--4.} Choisissez \textbf{UN} thème : \eng{Drug abuse} / \eng{Tobacco} / \eng{Alcohol abuse} / \eng{Elderly care}.
    \item \textbf{Rédigez} votre script (6--8 lignes) sur une feuille A4 (format radio). Respectez le plan : Hook $\rightarrow$ 2 Dangers $\rightarrow$ 2 Conseils (modaux + impératif) $\rightarrow$ Slogan.
    \item \textbf{Cochez} au moins \textbf{5 mots} de la banque lexicale (section 6.4).
    \item \textbf{Répétez} à voix basse 3 fois en appliquant la \textbf{Méthode} (section 6.5).
    \item \textbf{Passage « On Air » :} Un élève par groupe lit le message devant la classe (simulation micro). Les autres notent sur la grille d'auto-évaluation.
\end{enumerate}
\end{exercice}

\courssub{Grille d'auto-évaluation et évaluation par les pairs}

Après chaque passage, l'orateur et deux camarades remplissent cette grille (sur 8 points). L'enseignant valide ou ajuste.

\begin{aretenir}[Grille d'évaluation du message radio — Total : 8 points]
\begin{center}
\begin{tabularx}{\linewidth}{|X|c|c|c|X|}
\hline
\rowcolor{popblueL}
\textbf{Critère} & \textbf{2 pts (Maîtrisé)} & \textbf{1 pt (En cours)} & \textbf{0 pt (Non acquis)} & \textbf{Indicateurs observables} \\
\hline
\textbf{Lexique (2 pts)} & 5+ mots de la leçon utilisés correctement. & 3--4 mots utilisés ou 1--2 erreurs d'emploi. & $<$ 3 mots ou erreurs graves (faux amis). & \eng{abuse, addiction, lungs, liver, elderly, awareness...} \\
\hline
\textbf{Grammaire (2 pts)} & Modaux (\eng{must/should/ought to}) + Impératifs corrects. Noms indénombrables respectés. & 1 erreur sur modaux ou article (\eng{the alcohol}). & Erreurs systématiques sur modaux / articles. & \eng{You must stop} / \eng{Alcohol kills} (pas \eng{The alcohol}). \\
\hline
\textbf{Structure (2 pts)} & Plan Hook/2 Dangers/2 Conseils/Slogan respecté. 6--8 lignes. Cohérence logique. & Plan respecté mais déséquilibré (ex: 3 dangers, 1 conseil) ou hors lignes. & Plan non respecté, idées déconnectées. & Transitions fluides : \eng{First... Second... So... Remember...} \\
\hline
\textbf{Prononciation (2 pts)} & Débit fluide, articulation claire des chiffres/mots-clés, intonation expressive (grave $\rightarrow$ joyeux). & Débit saccadé ou monotone, 2--3 mots mal prononcés. & Lecture mot à mot, inaudible, intonation plate. & \eng{WHO} / \eng{eight} / \eng{lungs} / \eng{elderly} bien dits. Pause // marquée. \\
\hline
\end{tabularx}
\end{center}
\end{aretenir}

\courssub{Exemples de slogans percutants (Banque d'idées)}

\begin{exemplebox}[Slogans modèles — Structure : Court + Rime/Parallélisme + Positif]
\begin{center}
\begin{tabularx}{\linewidth}{|X|X|}
\hline
\rowcolor{popblueL}
\textbf{English Slogan} & \textbf{Traduction française} \\
\hline
\eng{Say no to drugs, say yes to life!} & Dis non à la drogue, dis oui à la vie ! \\
\hline
\eng{Your health is your wealth.} & Ta santé est ta richesse. \\
\hline
\eng{Don't let tobacco steal your breath.} & Ne laisse pas le tabac voler ton souffle. \\
\hline
\eng{Love your liver, leave alcohol.} & Aime ton foie, laisse l'alcool. \\
\hline
\eng{Respect the elderly, they built our future.} & Respecte les aînés, ils ont bâti notre avenir. \\
\hline
\eng{Be smart, don't start.} & Sois malin, ne commence pas. \\
\hline
\eng{Addiction is a prison, prevention is the key.} & L'addiction est une prison, la prévention est la clé. \\
\hline
\end{tabularx}
\end{center}
\emph{Conseil :} Pour créer le vôtre, utilisez l'impératif (\eng{Choose, Protect, Love, Respect, Say}) + un nom fort (\eng{life, health, future, family, brain}).
\end{exemplebox}

\begin{savaistu}[Culture : La Journée Mondiale Sans Tabac au Cameroun]
Chaque \textbf{31 mai}, le Cameroun célèbre la \eng{World No Tobacco Day} (Journée mondiale sans tabac). Le Ministère de la Santé Publique (\eng{MINSANTÉ}) organise des marches à Yaoundé et Douala, des spots radio en langues nationales (Bassa, Ewondo, Fulfulde, Pidgin) et des consultations gratuites de sevrage tabagique dans les hôpitaux de district. Le slogan officiel 2023 était \eng{« We need food, not tobacco »} (Nous avons besoin de nourriture, pas de tabac), liant culture du tabac et sécurité alimentaire. Votre message radio pourrait être diffusé lors de cet événement !
\end{savaistu}

\newpage

\coursec{Activité d'intégration}

\coursec{Lesson 25 — Listening \& Speaking: Health Problems, Addictions and Aging}

\courssub{Objectifs de la leçon}
\begin{aretenir}[Compétences visées (Programme MINESEC Terminale A -- Chapter 3)]
À l'issue de cette séquence de 2 h, l'élève doit être capable de :
\begin{itemize}
    \item \textbf{Listening :} Comprendre un texte authentique (script radio) sur les problèmes de santé liés aux addictions (drogues, tabac, alcool) et au vieillissement ; repérer le thème, les mots-clés, les chiffres et les noms propres.
    \item \textbf{Vocabulary :} Mobiliser le lexique thématique (substance abuse, withdrawal, cirrhosis, cognitive decline, prescription drug misuse, life expectancy, etc.) en compréhension et en production.
    \item \textbf{Grammar :} Réinvestir le discours rapporté (backshifting), la voix passive, les modaux de conseil/obligation (\eng{should, must, had better}) et les questions indirectes.
    \item \textbf{Speaking :} Produire une mini-émission radio structurée (3 min) en équipe : accroche, données chiffrées, expertise médicale, interaction auditeur, conclusion.
    \item \textbf{Critical Thinking :} Évaluer la production d'autrui via une grille d'écoute active (thème, chiffres, conseils).
\end{itemize}
\end{aretenir}

\courssub{Étape 1 -- Mise en route et anticipation (10 min)}
\begin{experience}[Anticipation : « Health Watch » -- Prédire le contenu]
\begin{enumerate}
    \item Observez le titre de l'émission : \eng{“Health Watch”} et le nom de la radio : \eng{“Cameroon Health FM 94.5”}. Quel type d'informations attendez-vous ? (Santé publique, prévention, actualité médicale, numéros d'urgence).
    \item L'animatrice s'appelle \eng{Dr. Grace Nfor}. Quel registre de langue prévoit-on ? (Standard, formel mais accessible, ton engageant).
    \item \textbf{Brainstorming (2 min) :} En binôme, listez 10 mots-clés anglais que vous pourriez entendre sur le thème \eng{“substance abuse and healthy aging”}. Partagez au tableau.
    \item \textbf{Prédiction chiffrée :} L'animatrice annonce des « \eng{shocking figures} ». Quels types de données chiffrées ? (Pourcentages, nombres de cas, âges, coûts, années).
\end{enumerate}
\end{experience}

\courssub{Étape 2 -- Script d'écoute et glossaire (15 min)}
\begin{textebox}[Script audio / Texte support -- Introduction de l'émission « Health Watch »]
\eng{“Good morning, Yaoundé! Good morning, Douala, Buea, Bamenda and all our listeners across the ten regions! You are tuned to Cameroon Health FM, 94.5 FM, and this is \textbf{Health Watch}, your weekly rendezvous for reliable medical information and practical wellness tips. I’m your host, Dr. Grace Nfor.”}

\eng{“Today’s edition focuses on a silent epidemic that threatens both our youth and our elders: \textbf{substance abuse and the challenges of healthy aging}. According to a recent report by the Ministry of Public Health, non-communicable diseases linked to tobacco and alcohol are rising sharply in urban centres like Douala and Yaoundé. Meanwhile, our ageing population in the North-West and West regions faces specific risks: isolation, misuse of prescription drugs, and lack of adapted care.”}

\eng{“In the next twenty minutes, we will hear \textbf{the facts} from our data journalist, \textbf{the medical perspective} from a specialist in addiction medicine, and \textbf{your questions} from the streets of Mfoundi market. Stay with us — your health is your wealth.”}
\end{textebox}

\begin{definition}[Glossaire -- Mots et expressions clés du script]
\begin{center}
\begin{tabularx}{\linewidth}{|l|X|l|}
\hline
\rowcolor{popblueL} \textbf{English} & \textbf{Context / Definition} & \textbf{Français} \\ \hline
\eng{tuned to} & listening to a specific station & branché sur / à l'écoute de \\ \hline
\eng{rendezvous} & regular meeting/appointment & rendez-vous (émission régulière) \\ \hline
\eng{wellness tips} & practical health advice & conseils bien-être / santé \\ \hline
\eng{silent epidemic} & widespread health issue, unnoticed & épidémie silencieuse \\ \hline
\eng{substance abuse} & harmful use of drugs/alcool & abus de substances (drogues/alcool) \\ \hline
\eng{non-communicable diseases (NCDs)} & chronic diseases not passed person-to-person & maladies non transmissibles (MNT) \\ \hline
\eng{urban centres} & big cities & centres urbains / grandes villes \\ \hline
\eng{ageing population} & growing proportion of old people & population vieillissante \\ \hline
\eng{misuse of prescription drugs} & taking meds not as prescribed & détournement / usage abusif de médicaments sur ordonnance \\ \hline
\eng{adapted care} & healthcare suited to specific needs & prise en charge adaptée \\ \hline
\eng{data journalist} & reporter specializing in statistics & journaliste data / d'investigation chiffrée \\ \hline
\eng{addiction medicine} & medical specialty for substance use disorders & addictologie / médecine des addictions \\ \hline
\eng{your health is your wealth} & proverbial slogan & la santé est notre plus grande richesse \\ \hline
\end{tabularx}
\end{center}
\end{definition}

\courssub{Étape 3 -- Compréhension de l'oral / de l'écrit (20 min)}
\begin{exercice}[Questions sur le script d'introduction (Dr. Grace Nfor)]
\textbf{Consigne :} Lisez le texte (ou écoutez l'enregistrement) et répondez en anglais. Pour Vrai/Faux, justifiez par une citation du texte.

\begin{enumerate}
    \item \textbf{True / False / Not Mentioned -- Justify with a quote.}
    \begin{enumerate}
        \item The radio station broadcasts only in Yaoundé and Douala.
        \item Dr. Grace Nfor is a medical doctor.
        \item The show \eng{Health Watch} is broadcast daily.
        \item Non-communicable diseases are decreasing in Cameroonian cities.
        \item The ageing population in the North-West and West regions suffers from isolation.
    \end{enumerate}

    \item \textbf{Multiple Choice. Choose the correct answer.}
    \begin{enumerate}
        \item What is the main theme of today's edition?
        \begin{itemize}
            \item[A] Road safety in Cameroon.
            \item[B] Substance abuse and healthy aging.
            \item[C] Malaria prevention.
            \item[D] Nutrition for children.
        \end{itemize}
        \item Who will speak during the show?
        \begin{itemize}
            \item[A] Only the host.
            \item[B] A data journalist, a medical expert, and callers.
            \item[C] The Minister of Health.
            \item[D] Students from a local school.
        \end{itemize}
    \end{enumerate}

    \item \textbf{Short Answers (Full sentences).}
    \begin{enumerate}
        \item What is the frequency of the radio station?
        \item Which ministry published the recent report mentioned?
        \item Name two specific risks faced by the ageing population according to the text.
        \item Where are the callers speaking from?
    \end{enumerate}
\end{enumerate}
\end{exercice}

\begin{corrige}[Exercice -- Compréhension]
\begin{enumerate}
    \item \textbf{True / False}
    \begin{enumerate}
        \item \textbf{False.} \eng{“...and all our listeners across the ten regions!”}
        \item \textbf{True.} \eng{“I’m your host, Dr. Grace Nfor.”} (Title \eng{Dr.} implies medical doctor).
        \item \textbf{False.} \eng{“...your weekly rendezvous...”} (Weekly = hebdomadaire).
        \item \textbf{False.} \eng{“...are rising sharply...”}
        \item \textbf{True.} \eng{“...faces specific risks: isolation...”}
    \end{enumerate}
    \item \textbf{MCQ}
    \begin{enumerate}
        \item B (\eng{“substance abuse and the challenges of healthy aging”}).
        \item B (\eng{“we will hear the facts from our data journalist, the medical perspective from a specialist..., and your questions...”}).
    \end{enumerate}
    \item \textbf{Short Answers}
    \begin{enumerate}
        \item The frequency is 94.5 FM.
        \item The Ministry of Public Health published the report.
        \item Isolation and misuse of prescription drugs (or lack of adapted care).
        \item The callers are speaking from Mfoundi market.
    \end{enumerate}
\end{enumerate}
\end{corrige}

\courssub{Étape 4 -- Lexique thématique : Addictions et vieillissement (15 min)}
\begin{propriete}[Lexique clé -- Addictions \& Vieillissement (Consolidation)]
\begin{center}
\begin{tabularx}{\linewidth}{|l|X|l|}
\hline
\rowcolor{popblueL} \textbf{English Term} & \textbf{Definition / Context} & \textbf{Français} \\ \hline
\eng{Substance abuse} & Misuse of drugs/alcohol & Abus de substances \\ \hline
\eng{Binge drinking} & Heavy alcohol consumption in short time & Beuverie express / biture \\ \hline
\eng{Second-hand smoke} & Involuntary inhalation & Tabagisme passif \\ \hline
\eng{Withdrawal symptoms} & Physical effects when stopping & Symptômes de sevrage \\ \hline
\eng{Cirrhosis (of the liver)} & Irreversible liver scarring & Cirrhose (du foie) \\ \hline
\eng{Cognitive decline / Dementia} & Loss of memory, reasoning & Déclin cognitif / Démence \\ \hline
\eng{Prescription drug misuse} & Using meds without doctor's order & Détournement de médicaments \\ \hline
\eng{Peer pressure} & Influence from friends & Pression des pairs \\ \hline
\eng{Rehabilitation centre} & Facility for recovery & Centre de désintoxication / de réadaptation \\ \hline
\eng{Life expectancy} & Average years to live & Espérance de vie \\ \hline
\eng{Opioid / Tramadol} & Strong painkiller, risk of addiction & Opioïde / Tramadol \\ \hline
\eng{Overdose} & Excessive dangerous dose & Surdose \\ \hline
\eng{Harm reduction} & Policies to reduce negative consequences & Réduction des risques \\ \hline
\eng{Relapse} & Return to addiction after stopping & Rechute \\ \hline
\eng{Caregiver} & Person looking after elderly/sick & Aidant / soignant familial \\ \hline
\end{tabularx}
\end{center}
\end{propriete}

\begin{exercice}[Entraînement lexical -- Complétion]
Complétez les phrases avec le terme approprié (forme correcte).
\begin{enumerate}
    \item Many seniors suffer from \_\_\_\_\_\_\_\_\_\_ because they take too many different pills prescribed by different doctors. (\eng{prescription drug misuse / polypharmacy})
    \item \_\_\_\_\_\_\_\_\_\_ is a major cause of lung cancer for non-smokers living with smokers. (\eng{Second-hand smoke})
    \item The \_\_\_\_\_\_\_\_\_\_ in Cameroon has improved but remains lower than the global average. (\eng{life expectancy})
    \item Sudden cessation of alcohol by a dependent person can trigger severe \_\_\_\_\_\_\_\_\_\_. (\eng{withdrawal symptoms})
    \item \eng{Tramadol} is an \_\_\_\_\_\_\_\_\_\_ often misused for its euphoric effects. (\eng{opioid})
\end{enumerate}
\end{exercice}

\begin{corrige}[Exercice -- Lexique]
\begin{enumerate}
    \item prescription drug misuse (ou polypharmacy)
    \item Second-hand smoke
    \item life expectancy
    \item withdrawal symptoms
    \item opioid
\end{enumerate}
\end{corrige}

\courssub{Étape 5 -- Méthode d'écoute en 3 étapes et structures grammaticales (15 min)}
\begin{methode}[Méthode d'écoute active en 3 étapes (Pour l'auditeur)]
\begin{enumerate}
    \item \textbf{ANTICIPER (Avant) :} Lire la grille/consignes. Prédire le vocabulaire : Chiffres $\to$ \eng{percent, rate, number, cases, doubled} ; Conseils $\to$ \eng{should, must, avoid, stop, start, consult, reduce} ; Noms propres $\to$ \eng{Ministry, WHO, Committee, Hospital, Dr.}
    \item \textbf{ÉCOUTER SÉLECTIVEMENT (Pendant) :} Ne pas viser la compréhension totale. Cibler : \textbf{Mots-clés} (noms, verbes d'action), \textbf{Chiffres} (écouter \eng{teen} vs \eng{ty}, \eng{hundred, thousand}), \textbf{Connecteurs} (\eng{However, Moreover, Consequently}).
    \item \textbf{VÉRIFIER / RECONSTITUER (Après) :} Comparer notes avec un voisin. Déduire l'info manquante par le contexte. Rédiger la réponse finale en phrase complète et correcte.
\end{enumerate}
\end{methode}

\begin{propriete}[Structures grammaticales utiles pour la production orale]
\begin{itemize}
    \item \textbf{Discours rapporté (Journaliste / Faits) :} Backshifting si verbe introducteur au passé.
    \eng{The Ministry \textbf{reported} that 15\% of youths \textbf{had tried} tramadol.}
    \textit{Exception :} Si le fait est encore vrai $\to$ présent conservé. \eng{A survey \textbf{reveals} that 18\% \textbf{use} tramadol daily.}
    \item \textbf{Voix passive (Faits médicaux / Statistiques) :} Focus sur l'objet/patient.
    \eng{The liver \textbf{is damaged} by toxins.} / \eng{400 cases \textbf{were recorded} last year.}
    \item \textbf{Modaux de conseil, obligation, recommandation (Médecin / Animateur) :}
    \begin{center}
    \begin{tabular}{|l|l|l|}
    \hline
    \rowcolor{popblueL} \textbf{Modal} & \textbf{Nuance} & \textbf{Exemple} \\ \hline
    \eng{should / ought to} & Conseil standard & \eng{You \textbf{should} consult a doctor.} \\ \hline
    \eng{must / have to} & Obligation forte / Règle & \eng{Patients \textbf{must not} stop treatment abruptly.} \\ \hline
    \eng{had better} & Avertissement / Urgence (conséquence négative si non fait) & \eng{He \textbf{had better} reduce his intake.} \\ \hline
    \end{tabular}
    \end{center}
    \item \textbf{Connecteurs logiques (Animateur -- Structure) :}
    \eng{Firstly, / To begin with, / Moreover, / Furthermore, / However, / On the other hand, / Consequently, / As a result, / In conclusion, / To sum up.}
    \item \textbf{Questions indirectes polies (Animateur $\to$ Expert / Auditeur $\to$ Médecin) :}
    \eng{Could you tell us \textbf{what the main risks are}?} (Ordre S+V, pas d'auxiliaire \eng{do/does/did}).
    \eng{Would you mind explaining \textbf{how the treatment works}?}
\end{itemize}
\end{propriete}

\courssub{Étape 6 -- Production guidée : Mini-émission « Health Watch » (50 min)}

\courssub{Situation}
\begin{textebox}[Contexte : Vous êtes l'équipe de production de « Health Watch »]
\eng{“Good morning, Yaoundé! Good morning, Douala, Buea, Bamenda and all our listeners across the ten regions! You are tuned to Cameroon Health FM, 94.5 FM, and this is \textbf{Health Watch}, your weekly rendezvous for reliable medical information and practical wellness tips. I’m your host, Dr. Grace Nfor.”}

\eng{“Today’s edition focuses on a silent epidemic that threatens both our youth and our elders: \textbf{substance abuse and the challenges of healthy aging}. According to a recent report by the Ministry of Public Health, non-communicable diseases linked to tobacco and alcohol are rising sharply in urban centres like Douala and Yaoundé. Meanwhile, our ageing population in the North-West and West regions faces specific risks: isolation, misuse of prescription drugs, and lack of adapted care.”}

\eng{“In the next twenty minutes, we will hear \textbf{the facts} from our data journalist, \textbf{the medical perspective} from a specialist in addiction medicine, and \textbf{your questions} from the streets of Mfoundi market. Stay with us — your health is your wealth.”}
\end{textebox}

\textbf{Votre mission :} Par groupes de 4, vous êtes l'équipe de production de l'émission \eng{Health Watch}. Vous devez préparer et présenter une \textbf{mini-émission de 3 minutes} condensant la structure ci-dessus. Les autres groupes joueront le rôle des auditeurs et rempliront une grille d'écoute.

\courssub{Consignes détaillées}
\begin{enumerate}
    \item \textbf{Étape 1 -- Répartition des rôles (5 min).} Dans votre groupe, désignez :
    \begin{itemize}
        \item \textbf{L'Animateur/Animatrice (Host) :} Lance l'émission, annonce le thème, assure les transitions, conclut.
        \item \textbf{Le/La Journaliste Data (Data Journalist) :} Présente 2 à 3 chiffres-clés (inventés mais réalistes : prévalence, coûts, âges) sur une addiction précise (tabac, alcool, tramadol, drogue synthétique type \eng{“kush”}).
        \item \textbf{Le/La Médecin Expert(e) (Medical Expert) :} Explique \emph{un} mécanisme de danger (ex. : cirrhose, cancer du poumon, troubles cognitifs chez le senior) et donne \textbf{deux conseils concrets} de prévention ou de prise en charge.
        \item \textbf{L'Auditeur/L'Auditrice (Caller) :} Pose une question courte et réaliste (ex. : \eng{“My father drinks whisky every night, how can I help?”} ou \eng{“Is tramadol dangerous for my back pain?”}).
    \end{itemize}

    \item \textbf{Étape 2 -- Rédaction du canevas (15 min).} Rédigez \textbf{en anglais} les grandes lignes (mots-clés, phrases-clés, pas de texte lu mot à mot) sur une feuille A4 partagée. Utilisez les \textbf{Ressources} (Lexique + Grammaire) ci-dessus. Respectez la chronologie :
    \begin{enumerate}
        \item \textit{Jingle / Accroche (Host) :} 20 sec. (\eng{“Good morning Cameroon! This is Health Watch... Today: ... Stay tuned.”})
        \item \textit{Flash Info / Chiffres (Journalist) :} 40 sec. (\eng{“The numbers are alarming. A 2023 survey reveals that... cases were recorded...”})
        \item \textit{Interview Expert (Host + Doctor) :} 1 min 30 sec. (\eng{“Dr X, why is... risky?”} $\to$ \eng{“As we age... risks include... My advice: 1. ... 2. ...”})
        \item \textit{Appel Auditeur + Réponse (Caller + Doctor) :} 30 sec. (\eng{“Hello Doctor. My husband... Should I stop it?”} $\to$ \eng{“Do not stop suddenly. Book an appointment...”})
        \item \textit{Conclusion / Numéros utiles (Host) :} 20 sec. (\eng{“Remember: Health is a daily choice. Call Green Line 1510... Stay healthy!”})
    \end{enumerate}

    \item \textbf{Étape 3 -- Répétition et ajustement (5 min).} Chronométrez-vous. Vérifiez la prononciation des termes médicaux :
    \begin{center}
    \begin{tabular}{|l|l|}
    \hline
    \rowcolor{popblueL} \textbf{Terme} & \textbf{Prononciation (approx. FR)} \\ \hline
    \eng{cirrhosis} & « si-ro-sis » /sɪˈrəʊsɪs/ \\ \hline
    \eng{dementia} & « di-men-sia » /dɪˈmenʃə/ \\ \hline
    \eng{withdrawal} & « wi-th-draw-al » /wɪðˈdrɔːəl/ \\ \hline
    \eng{second-hand smoke} & « se-kend hand smouk » \\ \hline
    \eng{opioid} & « o-pi-oïd » /ˈəʊpiɔɪd/ \\ \hline
    \eng{geriatric} & « dje-ri-a-trik » /ˌdʒeriˈætrɪk/ \\ \hline
    \eng{seizures} & « si-jures » /ˈsiːʒəz/ \\ \hline
    \eng{physiotherapy} & « fi-zio-the-ra-pi » /ˌfɪziəʊˈθerəpi/ \\ \hline
    \end{tabular}
    \end{center}
    L'animateur doit parler \emph{à} l'auditeur (ton direct, \eng{you}, questions rhétoriques : \eng{“Did you know...?”}).

    \item \textbf{Étape 4 -- Enregistrement / Direct \& Écoute active (20 min).} Chaque groupe passe devant la classe (ou s'enregistre sur téléphone).
    \begin{itemize}
        \item \textbf{Groupes « Auditeurs » :} Remplissez la \textbf{Grille d'écoute} (fournie ci-après) pendant la prestation.
        \item \textbf{Groupe « Producteurs » :} Parlez distinctement, sans lire, en regardant le public.
    \end{itemize}

    \item \textbf{Étape 5 -- Retour et correction (10 min).} Après chaque passage :
    \begin{itemize}
        \item Les auditeurs donnent leurs réponses (thème, 2 chiffres, 2 conseils).
        \item Le groupe producteur confirme ou corrige.
        \item L'enseignant note les points forts (vocabulaire, structure, prononciation) et les points à revoir (grammaire, fluidité).
    \end{itemize}
\end{enumerate}

\begin{center}
\begin{tabularx}{\linewidth}{|X|X|X|}
\hline
\rowcolor{popblueL} \textbf{Grille d'écoute -- Health Watch} & \textbf{Groupe n° \_\_\_} & \textbf{Thème principal} \\ \hline
\textbf{Chiffre 1} (sujet + valeur) & \textbf{Chiffre 2} (sujet + valeur) & \textbf{Conseil 1} (verbe + action) \\ \hline
& & \\ \hline
\textbf{Conseil 2} (verbe + action) & \textbf{Question de l'auditeur} (résumée) & \textbf{Réponse du médecin} (idée clé) \\ \hline
& & \\ \hline
\end{tabularx}
\end{center}

\courssub{Proposition de production et corrigé commenté}
\begin{exoresolu}[Modèle de mini-émission « Health Watch » (3 min -- Niveau Terminale A)]
\eng{\textbf{[JINGLE: Upbeat music, fade under]}} \\
\eng{\textbf{Host:} Good morning Cameroon! This is \textbf{Health Watch} on Cameroon Health FM. I'm Grace Nfor. Today: \textbf{“Tramadol and the Elderly: A Hidden Danger”}. Did you know that in the Centre Region alone, emergency admissions for opioid misuse among over-60s have doubled in five years? Stay tuned.}
\\
\eng{\textbf{[TRANSITION: Short beep]}} \\
\eng{\textbf{Journalist:} Thanks, Grace. The numbers are alarming. A 2023 survey by the \textbf{National Anti-Drug Committee} reveals that \textbf{18\% of seniors} in Yaoundé and Douala use tramadol daily for chronic pain — often without a prescription. Even worse: \textbf{400 emergency cases} were recorded last year just in Yaoundé Central Hospital linked to tramadol overdose in patients over 60. That’s more than one per day.}
\\
\eng{\textbf{Host:} Shocking figures. Dr. Atem, you are a geriatric specialist at the Yaoundé General Hospital. Why is tramadol so risky for older people?} \\
\eng{\textbf{Doctor:} Thank you. As we age, our \textbf{kidney function slows down}. The drug stays in the body longer, leading to \textbf{toxicity}. Risks include \textbf{confusion, falls, respiratory depression} and \textbf{seizures}. Many seniors take it for arthritis, but they mix it with other painkillers or alcohol — a deadly cocktail.} \\
\eng{\textbf{Host:} What should families do?} \\
\eng{\textbf{Doctor:} \textbf{First, never share medication.} A dose for a 30-year-old can kill a 75-year-old. \textbf{Second, ask the doctor for a “medication review”} every six months. Non-drug options like physiotherapy are safer for chronic pain.} \\
\eng{\textbf{Host:} We have a call from Madame Ngo in Mfoundi. Madame Ngo, you're on air.} \\
\eng{\textbf{Caller:} Hello Doctor. My husband, 72, takes tramadol for his knee. He says he “needs it to walk”. But he sleeps all day and forgets names. Should I stop it?} \\
\eng{\textbf{Doctor:} Madam, \textbf{do not stop suddenly} — withdrawal can cause seizures. \textbf{Book an appointment this week} for a dose reduction plan. The memory loss and drowsiness are classic side effects. There are alternatives.} \\
\eng{\textbf{Host:} Wise advice. Remember: \textbf{Health is a daily choice.} For help, call the \textbf{Green Line 1510} — free, anonymous, 24/7. This was Health Watch. I'm Grace Nfor. Stay healthy, Cameroon!} \\
\eng{\textbf{[JINGLE: Music up, fade out]}}
\tcblower
\textbf{Commentaire des choix de langue (niveau Terminale A) :}
\begin{itemize}
    \item \textbf{Registre :} Anglais standard radiophonique, clair, accessible, sans argot excessif. Ton \textbf{engageant} (\eng{Stay tuned, Shocking figures, Wise advice}).
    \item \textbf{Vocabulaire ciblé :} \eng{opioid misuse, chronic pain, prescription, kidney function, toxicity, respiratory depression, seizures, medication review, withdrawal, side effects, alternatives, anonymous helpline}. Tous issus du programme.
    \item \textbf{Grammaire :}
    \begin{itemize}
        \item \textit{Passif} : \eng{400 emergency cases \textbf{were recorded}} (focus sur le fait, pas l'auteur).
        \item \textit{Discours rapporté} : \eng{A survey \textbf{reveals that} 18\% \textbf{use}...} (Present reporting verb $\to$ present tense kept for current fact).
        \item \textit{Modaux} : \eng{\textbf{should} ask, \textbf{must not} stop, \textbf{had better} reduce} (nuances : conseil fort, interdiction médicale, urgence).
        \item \textit{Questions indirectes} : \eng{Could you tell us \textbf{why it is} so risky?} (Structure complexe attendue au Bac).
    \end{itemize}
    \item \textbf{Prononciation (points de vigilance) :} \eng{Tramadol /ˈtræmədɒl/}, \eng{opioid /ˈəʊpiɔɪd/}, \eng{geriatric /ˌdʒeriˈætrɪk/}, \eng{seizures /ˈsiːʒəz/}, \eng{physiotherapy /ˌfɪziəʊˈθerəpi/}. Liaisons : \eng{over-60s /əʊvə sɪkstiːz/}, \eng{one per day /wʌn pɜː deɪ/}.
    \item \textbf{Contexte camerounais :} Références géographiques (Yaoundé, Douala, Mfoundi, Centre Region), institutionnelles (Ministry of Public Health, National Anti-Drug Committee, Green Line 1510), sociologiques (arthritis, family care, sharing meds).
\end{itemize}
\end{exoresolu}

\courssub{Bilan de l'activité}
\begin{aretenir}[Auto-évaluation -- Grille de compétences]
Remplis ce tableau honnêtement après l'activité. Coche la case correspondante.
\begin{center}
\begin{tabularx}{\linewidth}{|X|c|c|c|}
\hline
\rowcolor{popblueL} \textbf{Compétence évaluée} & \textbf{Non acquis} & \textbf{En cours} & \textbf{Acquis} \\ \hline
Je sais identifier le thème, 2 chiffres et 2 conseils en écoutant un texte authentique. & & & \\ \hline
J'utilise correctement le lexique des addictions et du vieillissement à l'oral. & & & \\ \hline
Je structure une prise de parole continue (rôle animateur/expert) avec connecteurs. & & & \\ \hline
J'interagis spontanément (rôle auditeur/médecin) en utilisant les modaux de conseil. & & & \\ \hline
Je maîtrise la prononciation des termes médicaux clés (stress, sons voyelles). & & & \\ \hline
Je collabore efficacement en groupe (respect temps, écoute, aide aux pairs). & & & \\ \hline
\end{tabularx}
\end{center}
\end{aretenir}

\begin{attention}[Pièges fréquents relevés pendant l'activité]
\begin{itemize}
    \item \textbf{Confusion \eng{addiction} / \eng{dependence} :} \eng{Addiction} = comportement compulsif (psychologique) ; \eng{Dependence} = adaptation physiologique (sevrage). Ne pas les utiliser comme synonymes parfaits.
    \item \textbf{Faux ami \eng{“Une prescription”} :} En anglais, \eng{a prescription} = l'ordonnance (le papier) ; \eng{prescription drugs} = médicaments sur ordonnance. Ne pas dire \eng{“He takes a prescription”} pour “il prend un médicament” $\to$ \eng{He takes prescription medication}.
    \item \textbf{Ordre des mots question indirecte :} \eng{*Could you tell us what are the risks?} $\to$ \textbf{FAUX}. Correct : \eng{Could you tell us \textbf{what the risks are}?}
    \item \textbf{Prononciation \eng{“senior”} :} Pas « sénior » à la française \eng{/senjɔʁ/} mais \eng{/ˈsiːniər/} (US) ou \eng{/ˈsiːnɪə/} (UK).
    \item \textbf{Oubli du \eng{the} avec \eng{elderly} :} \eng{The elderly are vulnerable} (nominalisé, pluriel), pas \eng{*Elderly are...}.
    \item \textbf{Accord \eng{“The number of” vs “A number of”} :} \eng{The number of cases \textbf{is} rising} (singulier) vs \eng{A number of cases \textbf{were} recorded} (pluriel).
\end{itemize}
\end{attention}

\begin{savaistu}[Le saviez-vous ? -- Le numéro vert 1510 au Cameroun]
Le \textbf{1510} (ou \textbf{119} pour les violences) est un numéro d'urgence \textbf{gratuit} et \textbf{anonyme} accessible 24h/24 au Cameroun. Il est géré par le \textbf{Ministère de la Santé Publique} en partenariat avec des ONG. Il ne sert pas qu'aux urgences vitales : on peut y appeler pour une \textbf{orientation vers un centre de prise en charge des addictions} (ex. Centre Jamot à Yaoundé, Centre de santé mentale de Douala), pour un conseil sur l'arrêt du tabac ou de l'alcool, ou pour signaler une vente illicite de tramadol près d'une école. Dans votre émission, citer ce numéro donne une \textbf{valeur de service public} à votre production orale -- critère valorisé à l'oral du Baccalauréat.
\end{savaistu}

\newpage

\coursec{Méthodes et exercices résolus}

\courssub{Méthodes et exercices résolus}

\begin{methode}[Stratégie d'écoute globale : identifier le thème, les mots-clés et les chiffres]
\textbf{Objectif :} Comprendre le message principal d'un document audio (ou lu) sans comprendre chaque mot.
\begin{enumerate}
    \item \textbf{Avant l'écoute (Anticipation) :} Lis le titre, les images ou les questions. Prédis le vocabulaire probable (\eng{addiction, withdrawal, rehabilitation, life expectancy, chronic disease}). Note les mots-clés attendus.
    \item \textbf{Première écoute (Globalité) :} Ne t'arrête pas sur les mots inconnus. Repère : \textbf{Qui parle ?} (expert, témoin, journaliste), \textbf{De quoi ?} (drogues, tabac, alcool, vieillissement), \textbf{Ton ?} (alarmiste, informatif, espoir).
    \item \textbf{Repérage des \cle{signaux sonores} :} Insiste sur les \textbf{chiffres} (dates, pourcentages, âges), les \textbf{noms propres} (OMS, Yaoundé, Ministère de la Santé), les \textbf{connecteurs logiques} (\eng{however, therefore, consequently, in contrast}).
    \item \textbf{Deuxième écoute (Validation) :} Coche les hypothèses, complète les chiffres, note les arguments principaux (causes $\rightarrow$ conséquences $\rightarrow$ solutions).
\end{enumerate}
\textbf{Reflexe Bac :} Les questions \eng{True/False} ou \eng{Multiple Choice} portent souvent sur l'idée générale ou un chiffre clé. Élimine les réponses trop précises non entendues.
\end{methode}

\begin{exoresolu}[Compréhension globale : « Youth and Substance Abuse in Cameroon »]
\textbf{Écoute le script suivant (lu par le professeur ou audio) et réponds aux questions.}
\begin{textebox}[Script d'écoute original — Niveau Terminale A]
Good morning, everyone. Today, we focus on a growing concern in our major cities like Douala and Yaoundé: substance abuse among young people. A recent survey by the Ministry of Public Health reveals that fifteen percent of students aged fifteen to twenty have already experimented with hard drugs such as tramadol or cocaine. Alcohol remains the most accessible substance; forty percent of secondary school students admit drinking regularly, often starting as early as thirteen years old. Tobacco use is also rising, especially shisha in trendy cafés. The consequences are dramatic: school dropout, mental health disorders like anxiety and depression, and a burden on families. However, solutions exist. The government has launched a toll-free number, 1510, for counselling. NGOs organise peer-education campaigns in schools. Remember: early intervention saves lives. Stay informed, stay safe.
\end{textebox}

\end{exoresolu}

\begin{aretenir}[Glossaire du script — Mots clés]
\begin{itemize}
    \item \textbf{substance abuse} (n) : abus de substances / toxicomanie
    \item \textbf{survey} (n) : enquête
    \item \textbf{reveal} (v) : révéler
    \item \textbf{experiment with} (v) : expérimenter, essayer
    \item \textbf{tramadol} (n) : tramadol (opioïde antidouleur détourné)
    \item \textbf{accessible} (adj) : accessible
    \item \textbf{admit} (v) : admettre, avouer
    \item \textbf{shisha} (n) : chicha / narguilé
    \item \textbf{trendy} (adj) : branché, à la mode
    \item \textbf{dropout} (n) : décrochage (scolaire)
    \item \textbf{mental health disorders} (n) : troubles de la santé mentale
    \item \textbf{anxiety} (n) : anxiété
    \item \textbf{depression} (n) : dépression
    \item \textbf{burden} (n) : fardeau
    \item \textbf{toll-free number} (n) : numéro vert / gratuit
    \item \textbf{counselling} (n) : conseil, écoute psychologique
    \item \textbf{NGO} (n) : ONG (Organisation Non Gouvernementale)
    \item \textbf{peer-education} (n) : éducation par les pairs
    \item \textbf{early intervention} (n) : intervention précoce
\end{itemize}
\end{aretenir}

\begin{exoresolu}[Questions et corrigé détaillé]
\textbf{Questions :}
\begin{enumerate}
    \item \textbf{True or False?} Justify with a phrase from the text.
    \begin{itemize}
        \item[a)] The survey concerns university students only.
        \item[b)] Alcohol is the most consumed substance among youth.
        \item[c)] The toll-free number for help is 1501.
    \end{itemize}
    \item \textbf{Multiple Choice:} The main purpose of this text is to:
    \begin{itemize}
        \item A) Advertise a new rehabilitation centre.
        \item B) Inform about a health issue and existing solutions.
        \item C) Criticize the government's inaction.
    \end{itemize}
    \item \textbf{Pick out two consequences of substance abuse mentioned in the text.}
\end{enumerate}

\tcblower
\textbf{Solution détaillée :}

\textbf{1. True / False + Justification}
\begin{itemize}
    \item[a)] \textbf{False.} Justification : « \eng{A recent survey by the Ministry of Public Health reveals that fifteen percent of students aged fifteen to twenty...} » $\rightarrow$ Le texte précise \eng{students aged fifteen to twenty} (lycéens/collégiens), pas \eng{university students}.
    \item[b)] \textbf{True.} Justification : « \eng{Alcohol remains the most accessible substance; forty percent of secondary school students admit drinking regularly...} » $\rightarrow$ \eng{most accessible} + chiffre \eng{forty percent} confirment.
    \item[c)] \textbf{False.} Justification : « \eng{The government has launched a toll-free number, 1510, for counselling.} » $\rightarrow$ Le chiffre exact est \eng{1510}, pas 1501. \textbf{Piège :} inversion de chiffres.
\end{itemize}

\textbf{2. Multiple Choice}
\textbf{Bonne réponse : B.}
\textbf{Raisonnement :} Le texte donne des \textbf{statistiques} (15\%, 40\%, âge 13), cite des \textbf{conséquences} (dropout, mental health) et présente des \textbf{solutions} (numéro 1510, campagnes ONG). L'option A n'est pas mentionnée (pas de centre spécifique). L'option C est fausse car le gouvernement agit (numéro vert).

\textbf{3. Conséquences (deux attendues parmi) :}
\begin{itemize}
    \item \eng{School dropout} (décrochage scolaire)
    \item \eng{Mental health disorders like anxiety and depression} (troubles santé mentale)
    \item \eng{A burden on families} (fardeau pour les familles)
\end{itemize}

\textbf{Conclusion :} La méthode « Anticipation $\rightarrow$ Écoute globale $\rightarrow$ Repérage chiffres/noms $\rightarrow$ Validation » permet de répondre sans comprendre \eng{tramadol} ou \eng{peer-education} parfaitement.


\end{exoresolu}

\begin{methode}[Écoute détaillée : relever informations spécifiques, causes et conséquences]
\textbf{Objectif :} Répondre à des questions précises (QCM, réponses courtes, tableau à compléter) exigeant une écoute fine.
\begin{enumerate}
    \item \textbf{Lis les questions AVANT la 2ème écoute.} Souligne les \cle{mots-questions} (\eng{Who, What, Why, How many, Which consequence, According to...}).
    \item \textbf{Prédis la nature de la réponse :} Un chiffre ? Un nom propre ? Un verbe d'action ? Une opinion ?
    \item \textbf{Écoute activement :} Trace au crayon les réponses potentielles. Méfie-toi des \textbf{distracteurs} : le locuteur cite une fausse piste puis la corrige (\eng{Some say 20\%, but the actual figure is 15\%}).
    \item \textbf{Gère le vocabulaire inconnu :} Déduis le sens par le contexte (ex. \eng{peer pressure} = pression des pairs/amis grâce à \eng{friends influence you}).
    \item \textbf{Vérifie la cohérence :} La réponse respecte-t-elle la grammaire de la question ? (Ex. \eng{How many...?} $\rightarrow$ chiffre ; \eng{Why...?} $\rightarrow$ \eng{because / due to}).
\end{enumerate}
\textbf{Structure utile :} \eng{The main cause is... / One consequence is... / According to the speaker, ... / Statistics show that...}
\end{methode}

\begin{exoresolu}[Compréhension détaillée : « The Hidden Costs of Tobacco and Alcohol »]
\begin{textebox}[Script d'écoute original — Niveau Terminale A]
Dr. Amadou Bello, a pulmonologist at the Yaoundé Central Hospital, warns against the trivialisation of smoking and drinking. ``Many patients think shisha is harmless because it smells like fruit,'' he explains. ``But one session of shisha equals smoking fifty cigarettes in terms of carbon monoxide intake.'' Regarding alcohol, he notes a shift: ``We used to see cirrhosis in patients over sixty. Now, we diagnose fatty liver disease in thirty-year-olds who consume local spirits like \eng{odontol} or cheap beer daily.'' The economic cost is staggering. The state spends billions of CFA francs annually treating tobacco-related cancers and alcohol-induced liver failures. Dr. Bello advocates for higher taxes on these products and plain packaging laws, already effective in countries like Senegal and Mauritius.
\end{textebox}

\end{exoresolu}

\begin{aretenir}[Glossaire du script — Mots clés]
\begin{itemize}
    \item \textbf{pulmonologist} (n) : pneumologue
    \item \textbf{warn against} (v) : mettre en garde contre
    \item \textbf{trivialisation} (n) : banalisation
    \item \textbf{harmless} (adj) : inoffensif
    \item \textbf{carbon monoxide} (n) : monoxyde de carbone
    \item \textbf{intake} (n) : absorption, inhalation
    \item \textbf{shift} (n) : changement, évolution
    \item \textbf{cirrhosis} (n) : cirrhose
    \item \textbf{fatty liver disease} (n) : stéatose hépatique (foie gras)
    \item \textbf{local spirits} (n) : alcools locaux / forts
    \item \textbf{odontol} (n) : odontol (alcool local artisanal très fort)
    \item \textbf{staggering} (adj) : ahurissant, énorme
    \item \textbf{annually} (adv) : annuellement
    \item \textbf{advocate for} (v) : plaider pour, recommander
    \item \textbf{plain packaging} (n) : paquet neutre (sans logo/marque)
\end{itemize}
\end{aretenir}

\begin{exoresolu}[Exercices et corrigé détaillé]
\textbf{Exercices :}
\begin{enumerate}
    \item \textbf{Complete the table with EXACT information from the text.}
    \begin{center}
    \begin{tabular}{|l|l|l|}\hline
    \rowcolor{popblueL} \textbf{Substance} & \textbf{Specific Danger Mentioned} & \textbf{Target Group / Age} \\ \hline
    Shisha & 1. \dots & Young people (trendy cafés) \\ \hline
    Local spirits (\eng{odontol}) / Cheap beer & 2. \dots & 3. \dots \\ \hline
    \end{tabular}
    \end{center}
    \item \textbf{Answer in English (full sentences).}
    \begin{itemize}
        \item[a)] What comparison does Dr. Bello make about shisha?
        \item[b)] What two policy measures does he recommend?
    \end{itemize}
    \item \textbf{Vocabulary in context:} Match the words (a-c) to their definitions (1-3).
    \begin{itemize}
        \item[a)] Trivialisation \quad [b)] Staggering \quad [c)] Advocates
        \item[1)] Shocking, surprisingly large \quad 2)] Making something seem less serious \quad 3)] Publicly supports/recommends
    \end{itemize}
\end{enumerate}

\tcblower
\textbf{Solution détaillée :}

\textbf{1. Tableau de réponses courtes (Information extraction)}
\begin{center}
\begin{tabular}{|l|l|l|}\hline
\rowcolor{popblueL} \textbf{Substance} & \textbf{Specific Danger Mentioned} & \textbf{Target Group / Age} \\ \hline
Shisha & \eng{One session equals smoking fifty cigarettes (carbon monoxide)} & Young people / trendy cafés \\ \hline
Local spirits / Cheap beer & \eng{Fatty liver disease} & \eng{Thirty-year-olds} \\ \hline
\end{tabular}
\end{center}
\textit{Note :} Pour \eng{Shisha}, le danger spécifique est l'équivalence \eng{50 cigarettes}. Pour l'alcool, la pathologie précise est \eng{fatty liver disease} (stéatose hépatique) et l'âge \eng{thirty-year-olds} (trentenaires), plus \eng{over sixty} pour l'ancienne cible (cirrhose).

\textbf{2. Réponses rédigées (Full sentences)}
\begin{itemize}
    \item[a)] \eng{Dr. Bello compares one shisha session to smoking fifty cigarettes in terms of carbon monoxide intake.} (Utilise \eng{compares... to} + structure de la phrase originale).
    \item[b)] \eng{He recommends higher taxes on tobacco and alcohol and plain packaging laws.} (Deux mesures : \eng{higher taxes} + \eng{plain packaging laws}). \textbf{Attention :} \eng{advocates for} $\rightarrow$ \eng{recommends/supports}.
\end{itemize}

\textbf{3. Vocabulaire (Contextualisation)}
\begin{itemize}
    \item[a)] Trivialisation $\rightarrow$ \textbf{2} (Making something seem less serious). \textit{Indice :} ``think shisha is harmless''.
    \item[b)] Staggering $\rightarrow$ \textbf{1} (Shocking, surprisingly large). \textit{Indice :} ``billions of CFA francs''.
    \item[c)] Advocates $\rightarrow$ \textbf{3} (Publicly supports/recommends). \textit{Indice :} ``advocates for higher taxes...''.
\end{itemize}

\textbf{Conclusion :} La lecture préventive des questions guide l'oreille vers \eng{fifty cigarettes}, \eng{fatty liver}, \eng{thirty-year-olds}, \eng{higher taxes}, \eng{plain packaging}.


\end{exoresolu}

\begin{methode}[Lexique thématique : addictions, santé, vieillissement — familles de mots et collocations]
\textbf{Objectif :} Maîtriser le vocabulaire précis pour comprendre et produire (writing/speaking).
\begin{enumerate}
    \item \textbf{Classe les mots par catégorie :} \cle{Substances} (\eng{opioids, cannabis, nicotine, spirits, sedatives}), \cle{Verbes d'action} (\eng{abuse, misuse, overdose, inject, inhale, quit, relapse, recover}), \cle{Noms abstraits} (\eng{addiction, dependence, withdrawal, tolerance, craving, relapse, rehabilitation, stigma}), \cle{Vieillissement} (\eng{life expectancy, chronic illness, dementia, arthritis, mobility, loneliness, palliative care, caregiver}).
    \item \textbf{Apprends les \cle{collocations} (associations figées) :}
    \begin{itemize}
        \item \eng{Substance abuse / Drug abuse} (pas \eng{maltraitance})
        \item \eng{Heavy smoker / drinker} (pas \eng{big})
        \item \eng{Withdrawal symptoms} (symptômes de sevrage)
        \item \eng{Peer pressure} (pression des pairs)
        \item \eng{Health hazard / risk factor}
        \item \eng{Ageing population / Elderly people} (pas \eng{old people} seul)
        \item \eng{Life expectancy has increased / dropped}
    \end{itemize}
    \item \textbf{Attention aux \cle{faux-amis} et calques :}
    \begin{center}
    \begin{tabularx}{\linewidth}{|X|X|X|}\hline
    \rowcolor{popblueL} \textbf{Français} & \textbf{Anglais correct} & \textbf{Piège / Erreur fréquente} \\ \hline
    Une drogue & \eng{A drug / A substance} & \eng{A dope} (argot), \eng{Medicine} = médicament \\ \hline
    Se soigner / Guérir & \eng{To recover / To heal / To quit} & \eng{To cure} (transitif : \eng{The doctor cures the patient}) \\ \hline
    Un toxicomane & \eng{A drug addict / A user} & \eng{A toxicoman} (n'existe pas) \\ \hline
    Le sevrage & \eng{Withdrawal / Detox} & \eng{Weaning} (sevrage bébé/animal) \\ \hline
    Une overdose & \eng{An overdose} & \eng{An overdosage} (rare) \\ \hline
    Vieillir & \eng{To age / To grow old} & \eng{To old} (verbe inexistant) \\ \hline
    Une personne âgée & \eng{An elderly person / A senior} & \eng{An old} (nom), \eng{An ancient} (antique) \\ \hline
    \end{tabularx}
    \end{center}
    \item \textbf{Entraîne-toi :} Transforme \eng{addict} (n/v) $\rightarrow$ \eng{addiction} (n) $\rightarrow$ \eng{addictive} (adj) $\rightarrow$ \eng{addictively} (adv). \eng{Depend} $\rightarrow$ \eng{dependence} $\rightarrow$ \eng{dependent} $\rightarrow$ \eng{independent}.
\end{enumerate}
\end{methode}

\begin{exoresolu}[Lexique en contexte : transformation et emploi]
\textbf{Exercice : Complète les phrases avec la forme correcte du mot entre parenthèses. Traduisez chaque phrase.}
\begin{enumerate}
    \item Many young people start smoking due to peer \_\_\_\_\_\_\_\_. (PRESS)
    \item The \_\_\_\_\_\_\_\_ symptoms were so severe he had to be hospitalised. (WITHDRAW)
    \item The government launched a campaign to fight drug \_\_\_\_\_\_\_\_ among students. (ABUSE)
    \item She works as a \_\_\_\_\_\_\_\_ in a centre for \_\_\_\_\_\_\_\_ people. (CARE / AGE)
    \item \_\_\_\_\_\_\_\_ packaging laws aim to make cigarettes less attractive. (PLAIN)
    \item Life \_\_\_\_\_\_\_\_ in Cameroon has risen to about 60 years. (EXPECT)
    \item He is \_\_\_\_\_\_\_\_ on painkillers after his accident. (DEPEND)
\end{enumerate}

\tcblower
\textbf{Solution détaillée :}

\begin{enumerate}
    \item \textbf{pressure} (Nom). \eng{Peer pressure} = collocation figée. \eng{« La pression des pairs pousse beaucoup de jeunes à fumer. »}
    \item \textbf{withdrawal} (Nom). \eng{Withdrawal symptoms} = collocation médicale. \textit{Verbe :} \eng{withdraw} $\rightarrow$ \eng{withdrawn} (participe). \eng{« Les symptômes de sevrage étaient si graves... »}
    \item \textbf{abuse} (Nom). \eng{Drug abuse / Substance abuse} = terme technique. \eng{« Le gouvernement lutte contre l'abus de drogues... »}
    \item \textbf{caregiver} (Nom agent) / \textbf{elderly} (Adj) ou \textbf{aged} (Adj). \eng{Caregiver} = personne qui soigne/aide (proche aidant). \eng{Elderly people} = personnes âgées (politesse). \eng{« Elle travaille comme aide-soignante dans un centre pour personnes âgées. »}
    \item \textbf{Plain} (Adj invariable). \eng{Plain packaging} = paquet neutre (loi). \eng{« Les lois sur le paquet neutre visent... »}
    \item \textbf{expectancy} (Nom). \eng{Life expectancy} = espérance de vie (terme statique). \eng{« L'espérance de vie au Cameroun a atteint... »}
    \item \textbf{dependent} (Adj). \eng{To be dependent ON} (préposition obligatoire). \eng{Addicted TO} aussi possible. \eng{« Il est dépendant aux antidouleurs... »}
\end{enumerate}

\textbf{Point de grammaire intégré :} \eng{Dependent ON / Addicted TO / Keen ON / Aware OF}. Prépositions figées après adjectifs.
\end{exoresolu}

\begin{methode}[Prise de notes structurée et restitution (Summary / Reporting)]
\textbf{Objectif :} Synthétiser un document audio long (2-3 min) sous forme de schéma, tableau ou résumé écrit (\eng{Reported Speech}).
\begin{enumerate}
    \item \textbf{Structure tes notes :} Utilise un organigramme : \textbf{Thème central} $\rightarrow$ \textbf{Causes} (flèches entrantes) $\rightarrow$ \textbf{Effets} (flèches sortantes) $\rightarrow$ \textbf{Solutions/Recommandations}.
    \item \textbf{Abrège :} Symboles (\&, +, -, $\uparrow$ augment, $\downarrow$ diminut, $\rightarrow$ entraîne, \$ coût, \# nombre). Mots-clés seulement (pas de phrases complètes).
    \item \textbf{Repère les marqueurs de structure :} \eng{Firstly / The main reason is / Another factor / As a result / Consequently / To tackle this / Authorities urge...}
    \item \textbf{Rédige le résumé (Reported Speech) :} Sujet + verbe introducteur (\eng{The speaker highlights / warns / states / explains / reveals that...}) + contenu. \textbf{Règle de concordance des temps :} Présent $\rightarrow$ Prétérit ; \eng{will} $\rightarrow$ \eng{would} ; \eng{can} $\rightarrow$ \eng{could}.
    \item \textbf{Vérifie :} As-tu les chiffres clés ? Les noms propres ? La conclusion ?
\end{enumerate}
\textbf{Modèle de phrase :} \eng{The report highlights that drug abuse is rising among teenagers. It states that 15\% are affected. It warns that consequences include dropout. It recommends calling 1510.}
\end{methode}

\begin{exoresolu}[Prise de notes et rapport : « Healthy Ageing in Buea »]
\begin{textebox}[Script d'écoute original — Niveau Terminale A]
In Buea, the Regional Delegate of Public Health, Dr. Eunice Njie, presented a plan for healthy ageing. She stated that the population over sixty has doubled in twenty years, reaching eight percent today. She identified three pillars: first, prevention through free screening for hypertension and diabetes at district hospitals every Wednesday. Second, social inclusion: the creation of ``Senior Clubs'' in every subdivision to fight isolation. Third, training caregivers: a new module at the University of Buea for nurses specialising in geriatrics. She concluded: ``Ageing is not a disease. It is a triumph of development. But it demands investment.''
\end{textebox}

\end{exoresolu}

\begin{aretenir}[Glossaire du script — Mots clés]
\begin{itemize}
    \item \textbf{Regional Delegate} (n) : Délégué régional
    \item \textbf{double} (v) : doubler
    \item \textbf{pillar} (n) : pilier
    \item \textbf{prevention} (n) : prévention
    \item \textbf{screening} (n) : dépistage
    \item \textbf{hypertension} (n) : hypertension
    \item \textbf{diabetes} (n) : diabète
    \item \textbf{district hospital} (n) : hôpital de district
    \item \textbf{social inclusion} (n) : inclusion sociale
    \item \textbf{subdivision} (n) : arrondissement / sous-division
    \item \textbf{isolation} (n) : isolement
    \item \textbf{caregiver} (n) : aidant / soignant (proche aidant)
    \item \textbf{geriatrics} (n) : gériatrie
    \item \textbf{triumph} (n) : triomphe
    \item \textbf{demand} (v) : exiger, nécessiter
    \item \textbf{investment} (n) : investissement
\end{itemize}
\end{aretenir}

\begin{exoresolu}[Exercices et corrigé détaillé]
\textbf{Exercices :}
\begin{enumerate}
    \item \textbf{Complete the mind-map / diagram based on the text.}
    \begin{center}
    \begin{tikzpicture}[mindmap, grow cyclic, every node/.style=concept, concept color=popblue!60,
        level 1/.append style={level distance=4.5cm, sibling angle=120},
        level 2/.append style={level distance=3cm, sibling angle=60}]
    \node {Healthy Ageing Plan (Buea)}
        child[concept color=popgreen!60] { node {1. Prevention}
            child { node {Free screening} }
            child { node {Hypertension / Diabetes} }
            child { node {Wednesdays} }
        }
        child[concept color=poporange!60] { node {2. Social Inclusion}
            child { node {Senior Clubs} }
            child { node {Fight isolation} }
        }
        child[concept color=poppurple!60] { node {3. Training}
            child { node {University of Buea} }
            child { node {Geriatrics nurses} }
        };
    \end{tikzpicture}
    \end{center}
    \item \textbf{Write a summary (Reported Speech) in 60-80 words. Start with:} \eng{Dr. Njie presented a plan for healthy ageing in Buea. She stated that...}
\end{enumerate}

\tcblower
\textbf{Solution détaillée :}

\textbf{1. Mind-map complété (voir schéma ci-dessus).}
\textbf{Infos clés extraites :}
\begin{itemize}
    \item Chiffres : \eng{doubled in 20 years}, \eng{8\%}, \eng{every Wednesday}.
    \item Noms propres : \eng{Dr. Eunice Njie}, \eng{Buea}, \eng{University of Buea}.
    \item 3 Piliers : Prevention (screening), Inclusion (Senior Clubs), Training (Geriatrics).
    \item Citation clé : \eng{Ageing is not a disease... triumph of development}.
\end{itemize}

\textbf{2. Résumé (Reported Speech) — \textit{environ 70 mots} :}
\eng{Dr. Njie presented a plan for healthy ageing in Buea. She stated that the population over sixty had doubled in twenty years, reaching eight percent. She outlined three pillars. Firstly, prevention through free screening for hypertension and diabetes at district hospitals every Wednesday. Secondly, social inclusion via the creation of Senior Clubs in every subdivision to fight isolation. Thirdly, training caregivers with a new geriatrics module for nurses at the University of Buea. She concluded that ageing was a triumph of development demanding investment.}

\textbf{Analyse grammaticale (Concordance des temps) :}
\begin{itemize}
    \item \eng{has doubled} (Present Perfect) $\rightarrow$ \eng{had doubled} (Past Perfect).
    \item \eng{identified} (Past) $\rightarrow$ \eng{outlined} (verbe de rapport varié) + \eng{was/were} (Past) pour les détails.
    \item \eng{is not a disease} (Present) $\rightarrow$ \eng{was not a disease} (Past).
    \item \eng{demands} (Present) $\rightarrow$ \eng{demanded} (Past).
\end{itemize}
\textbf{Note :} Pas de \eng{that} obligatoire après \eng{stated/outlined}, mais recommandé pour la clarté.


\end{exoresolu}

\begin{methode}[Production guidée : rédiger un message de sensibilisation (SMS, Affiche, Discours court)]
\textbf{Objectif :} Réemployer le lexique et les structures pour \textbf{informer, alerter, convaincre} (Compétence Writing/Speaking).
\begin{enumerate}
    \item \textbf{Cible \& Canal :} \eng{SMS} (court, impératif, abréviations acceptées) vs \eng{Poster} (visuel, slogan + 3 faits + numéro) vs \eng{Speech} (formel, \eng{Ladies and gentlemen, I urge you to...}).
    \item \textbf{Structure type (AIDA) :}
    \begin{itemize}
        \item \textbf{Attention :} Accroche choc (\eng{Don't let shisha steal your future! / 15\% of our youth at risk.}).
        \item \textbf{Interest :} Fait clé / Chiffre (\eng{One session = 50 cigarettes.}).
        \item \textbf{Desire/Action :} Solution concrète (\eng{Call 1510 today. / Say NO to peer pressure.}).
        \item \textbf{Action :} Appel à l'action clair (\eng{Visit the health centre. / Talk to a teacher.}).
    \end{itemize}
    \item \textbf{Outils linguistiques :}
    \begin{itemize}
        \item \textbf{Impératif :} \eng{Stop smoking. / Protect your liver. / Know the facts.}
        \item \textbf{Modaux d'obligation/conseil :} \eng{You must / should / have to quit.} \eng{Don't start.}
        \item \textbf{Conditionnel 1er type (Real consequence) :} \eng{If you drink daily, you will damage your liver.}
        \item \textbf{Passif (Focus victime) :} \eng{Lives are destroyed. / Families are burdened.}
    \end{itemize}
    \item \textbf{Relecture :} Vérifie orthographe (\eng{addiction, alcohol, cigarettes, ageing}), majuscules (OMS, Cameroun), numéros utiles.
\end{enumerate}
\end{methode}

\begin{exoresolu}[Production écrite : Message de sensibilisation « Zero Tramadol in Schools »]
\textbf{Consigne :} Tu es délégué(e) de ta classe à Yaoundé. Écris le texte d'une \textbf{affiche A4} (120-150 mots) pour alerter sur le Tramadol. Utilise : titre, 3 dangers, 1 statistique, numéro vert 1510, slogan final. Ton : ferme mais bienveillant.

\tcblower
\textbf{Proposition de rédaction (Modèle corrigé) :}

\begin{textebox}[Modèle d'affiche — Writing Task]
\textbf{STOP TRAMADOL: PROTECT YOUR FUTURE!}

\textbf{Did you know?} In Cameroon, 15\% of students aged 15-20 have already tried tramadol. It is \textbf{NOT} a simple painkiller. It is a dangerous opioid.

\textbf{3 REASONS TO SAY NO:}
\begin{itemize}
    \item \textbf{Addiction traps you fast.} One pill creates dependence. Withdrawal is agony.
    \item \textbf{Your brain pays the price.} Memory loss, seizures, inability to study.
    \item \textbf{Your family suffers.} Violence, debt, broken trust.
\end{itemize}

\textbf{YOU ARE NOT ALONE.}
\textbf{FREE HELP 24/7: CALL 1510} (Toll-free)
Talk to a teacher, a nurse, or a trusted friend.

\textbf{YOUR LIFE MATTERS. CHOOSE HEALTH. CHOOSE SUCCESS.}
\end{textebox}

\textbf{Analyse didactique du modèle :}
\begin{itemize}
    \item \textbf{Title} : Impératif fort + Bénéfice (\eng{Protect your future}).
    \item \textbf{Hook} : Statistique contextualisée (\eng{In Cameroon... 15\%}).
    \item \textbf{Clarification} : \eng{NOT a simple painkiller} (corrige représentation erronée).
    \item \textbf{List 3 dangers} : Parallélisme structurel (\eng{Addiction traps... / Brain pays... / Family suffers}). Vocabulaire précis (\eng{opioid, dependence, withdrawal, seizures}).
    \item \textbf{Call to Action} : Numéro mis en évidence (\eng{CALL 1510}), alternatives sociales (\eng{teacher, nurse, friend}).
    \item \textbf{Slogan final} : Parallélisme \eng{Choose... Choose...}, valorisation (\eng{Your life matters}).
    \item \textbf{Grammaire} : Présent de vérité générale, Impératifs, Modaux (\eng{can}), Passif (\eng{is created}).
\end{itemize}
\end{exoresolu}

\begin{methode}[Stratégie « Gap-fill / Cloze test » : compréhension de la structure langue]
\textbf{Objectif :} Réussir les exercices de « texte à trous » (grammaire + lexique) fréquents au Bac.
\begin{enumerate}
    \item \textbf{Lis le texte entier d'abord} (sans regarder les trous) pour saisir le sens global (thème, temps dominant).
    \item \textbf{Analyse chaque trou :} Quelle \textbf{nature de mot} manque ? (Nom, Verbe, Adj, Adv, Préposition, Déterminant, Conjonction).
    \item \textbf{Indices grammaticaux :}
    \begin{itemize}
        \item Après \eng{a/an/the/this/my} $\rightarrow$ \textbf{Nom} (ou Adj + Nom).
        \item Après \eng{to} $\rightarrow$ \textbf{Base verbale} (Infinitif).
        \item Après \eng{be/have/do/modal} $\rightarrow$ \textbf{Forme verbale} (Participe, Base verbale).
        \item Après verbe/préposition $\rightarrow$ \textbf{Gérondif (-ing)} ou Nom.
        \item Fin de phrase / après \eng{very/so/really} $\rightarrow$ \textbf{Adjectif / Adverbe (-ly)}.
    \end{itemize}
    \item \textbf{Indices lexicaux (Collocations) :} \eng{Heavy \_\_\_\_\_} (smoker/drinker), \eng{At \_\_\_\_\_ risk} (high), \eng{Lead \_\_\_\_\_} (to).
    \item \textbf{Vérifie l'accord (Sujet-Verbe) et le temps} (contexte passé/présent/futur).
    \item \textbf{Relis le texte complet} à la fin : ça sonne juste ?
\end{enumerate}
\end{methode}

\begin{exoresolu}[Exercice de consolidation : Texte à trous (Grammaire + Lexique)]
\textbf{Complète le texte suivant avec UN SEUL MOT par trou.}
\textit{« The fight \_\_\_\_\_ (1) drug abuse in schools \_\_\_\_\_ (2) a national priority. Yesterday, the Minister \_\_\_\_\_ (3) launched a campaign targeting secondary students. He warned that \_\_\_\_\_ (4) shisha and tramadol \_\_\_\_\_ (5) becoming popular among teenagers. ``If we \_\_\_\_\_ (6) not act now,'' he said, ``we \_\_\_\_\_ (7) lose a whole generation.'' The programme includes talks by former addicts who \_\_\_\_\_ (8) \_\_\_\_\_ (9) recovered \_\_\_\_\_ (10) their addiction. Students \_\_\_\_\_ (11) encouraged to call 1510 \_\_\_\_\_ (12) they need help. \_\_\_\_\_ (13) is free and confidential. Let's work together \_\_\_\_\_ (14) a drug-free Cameroon. »}

\tcblower
\textbf{Solution détaillée :}

\begin{enumerate}
    \item \textbf{against} (Préposition). \eng{Fight against sth} (collocation). \textit{Piège :} \eng{Fight drug abuse} (verbe direct) possible mais ici nom \eng{The fight} $\rightarrow$ \eng{against}.
    \item \textbf{is} (Verbe BE, Présent, 3e pers. sing.). Sujet = \eng{The fight} (singulier). Temps = présent (actualité).
    \item \textbf{has} (Auxiliaire HAVE, Présent Perfect). \eng{has launched} (action passée conséquence présente). Sujet = \eng{The Minister} (singulier).
    \item \textbf{both} (Déterminant). \eng{Both shisha and tramadol} = les deux. \textit{Indice :} structure \eng{Both X and Y}.
    \item \textbf{are} (Verbe BE, Présent, 3e pers. plur.). Sujet = \eng{Both shisha and tramadol} (pluriel). \eng{are becoming}.
    \item \textbf{do} (Auxiliaire DO, Présent). Conditionnel 1er type : \eng{If we DO not act...} (négation). Sujet \eng{we}.
    \item \textbf{will} (Modal). Conséquence future : \eng{we WILL lose}.
    \item \textbf{have} (Auxiliaire HAVE, Présent Perfect). Sujet = \eng{former addicts} (pluriel). \eng{have recovered}.
    \item \textbf{successfully} / \textbf{fully} / \textbf{completely} (Adverbe en \textbf{-ly}). Modifie le verbe \eng{recovered}. \textit{Attention :} \eng{recovered} est verbe ici, pas adj $\rightarrow$ adv en -ly.
    \item \textbf{from} (Préposition). \eng{Recover FROM an illness/addiction} (verbe prépositionnel obligatoire).
    \item \textbf{are} (Verbe BE, Passif, Présent). \eng{Students ARE encouraged} (Passif : encourage $\rightarrow$ be encouraged). Sujet pluriel.
    \item \textbf{if} / \textbf{when} / \textbf{whenever} (Conjonction). \eng{Call 1510 IF you need help}. Condition / Temps.
    \item \textbf{It} (Sujet impersonnel / Dummy subject). \eng{IT is free} (coût / service). \eng{Calling is free} possible mais \eng{It} attendu (1 mot).
    \item \textbf{for} / \textbf{towards} (Préposition). \eng{Work FOR a drug-free Cameroon} (but / objectif). \eng{Towards} aussi possible.
\end{enumerate}

\textbf{Récapitulatif des natures grammaticales testées :} Prépositions (1, 10, 12, 14), Verbes auxiliaires/accords (2, 3, 5, 7, 8, 11), Déterminant (4), Modal (7), Adverbe (9), Conjonction (12), Sujet impersonnel (13).
\end{exoresolu}

\begin{attention}[Les 5 erreurs qui coûtent des points au Bac — Spécial « Health \& Addictions »]
\begin{enumerate}
    \item \textbf{« Drug abuse » vs « Drugs abuse » (Orthographe) :} N'écris jamais \eng{drugs abuse} (pluriel sur le premier nom). \textbf{Règle :} Nom composé = premier nom au singulier. \eng{Drug addiction, Alcohol consumption, Cigarette smoke, Health hazard}.
    \item \textbf{Faux-ami : « Une drogue » $\neq$ \eng{A dope}.} \eng{Dope} = argot (souvent cannabis/héroïne) ou adjectif \eng{cool}. Termes corrects : \eng{A drug, A substance, A narcotic, Medication} (médicament). \eng{To dope} = se doper (sport).
    \item \textbf{Calque : « Il est toxicomane » $\neq$ \eng{He is toxicoman}.} \textbf{Correct :} \eng{He is a drug addict / He is addicted to drugs / He uses drugs}. \eng{Addict} = nom OU adjectif. \eng{User} = neutre. \eng{Junkie} = très péjoratif (éviter).
    \item \textbf{Préposition : « Dépendant \textbf{de} » $\rightarrow$ \eng{Dependent \textbf{ON}} / \eng{Addicted \textbf{TO}}.}
    \begin{center}
    \begin{tabular}{|l|l|}\hline
    \rowcolor{popblueL} \textbf{Français} & \textbf{Anglais} \\ \hline
    Dépendre de & \eng{Depend ON} \\ \hline
    Accro à / Dépendant de (drogue) & \eng{Be addicted TO} \\ \hline
    Se sevrer de & \eng{Withdraw FROM / Quit / Give up} \\ \hline
    \end{tabular}
    \end{center}
    \item \textbf{Temps verbaux : « Depuis 20 ans, ça augmente » $\neq$ \eng{Since 20 years, it increases}.}
    \textbf{Correct :} \eng{It \textbf{has been increasing} \textbf{for} 20 years / \textbf{since} 2004.} (Present Perfect Continuous = durée action non terminée). \eng{For} + durée, \eng{Since} + point de départ.
\end{enumerate}
\textbf{Astuce finale :} Relis toujours tes réponses en te demandant : « Est-ce que ça sonne \textbf{anglais} ou \textbf{français traduit} ? » (Ex. \eng{The youth are sensitized} $\rightarrow$ \eng{Young people are made aware / sensitised}).
\end{attention}

\newpage

\coursec{Exercices}

\courssub{Je vérifie mes connaissances}

\begin{exercice}[QCM : Compréhension globale du script d'écoute]
Écoutez à nouveau l'extrait \eng{« Health Alert: Substance Abuse and Aging »} (script page précédente) et choisissez la bonne réponse parmi les trois propositions.
\begin{enumerate}
    \item According to the speaker, what is the \textbf{main} reason teenagers start smoking in secondary schools in Yaoundé?
    \begin{itemize}
        \item[a)] Curiosity about the taste of tobacco.
        \item[b)] Peer pressure and the desire to look mature.
        \item[c)] Influence of cigarette advertising on television.
    \end{itemize}
    \item What specific health consequence of \textbf{alcohol abuse} is mentioned regarding the liver?
    \begin{itemize}
        \item[a)] Liver enlargement (hepatomegaly).
        \item[b)] Cirrhosis and increased risk of liver cancer.
        \item[c)] Temporary liver inflammation that heals quickly.
    \end{itemize}
    \item The speaker states that \textbf{drug addiction} affects the brain by:
    \begin{itemize}
        \item[a)] Destroying neurons permanently after the first use.
        \item[b)] Altering the reward system, making it hard to feel pleasure naturally.
        \item[c)] Increasing intelligence and concentration for short periods.
    \end{itemize}
    \item Concerning the \textbf{elderly} (aging population), what does the report recommend for families in Buea and Bamenda?
    \begin{itemize}
        \item[a)] To send them to specialised hospitals in Douala immediately.
        \item[b)] To ensure regular medical check-ups and monitor medication side effects.
        \item[c)] To rely solely on traditional herbal medicine for chronic pains.
    \end{itemize}
\end{enumerate}
\end{exercice}

\begin{exercice}[Vrai / Faux justifié par une citation du texte]
Lisez le texte ci-dessous, puis dites si les affirmations sont \textbf{Vraies (V)} ou \textbf{Fausses (F)}. Justifiez chaque réponse par une phrase exacte copiée du texte (entre guillemets anglais “…”).
\begin{textebox}[Script : « The Silent Epidemic: Youth and Substance Abuse in Cameroon »]
“Recent studies conducted by the Ministry of Public Health in collaboration with WHO reveal a worrying trend in the Centre and Littoral regions. Approximately 21\% of students aged 13 to 17 have experimented with cannabis, locally known as ‘banga’, while 15\% consume alcohol regularly. Dr. Amadou Bello, a psychiatrist at the Jamot Hospital in Yaoundé, warns: ‘Early exposure rewires the adolescent brain, impairing memory and decision-making capacities. We are not just talking about bad grades; we are talking about a generation at risk of chronic psychosis.’ Furthermore, the rise of shisha bars near university campuses normalises tobacco inhalation among young adults who mistakenly believe it is safer than cigarettes. The nicotine concentration in one shisha session can equal that of a whole pack of cigarettes.”
\end{textebox}
\begin{enumerate}
    \item The study mentioned concerns only the city of Yaoundé.
    \item Cannabis is referred to locally as ‘banga’.
    \item Dr. Amadou Bello works in a hospital located in Douala.
    \item Shisha smoking delivers less nicotine than a pack of cigarettes.
    \item Early drug use can lead to chronic mental health disorders.
\end{enumerate}
\end{exercice}

\begin{exercice}[Vocabulaire thématique : Complétion]
Complétez les phrases suivantes avec le mot approprié choisi dans la liste ci-dessous. Attention, il y a deux mots en trop.
\begin{center}
\textbf{addiction \quad withdrawal \quad cirrhosis \quad passive \quad rehabilitation \quad overdose \quad carcinogen \quad secondhand}
\end{center}
\begin{enumerate}
    \item Long-term alcohol abuse is the leading cause of \_\_\_\_\_\_\_\_\_\_\_, a severe scarring of the liver.
    \item Many smokers ignore that \_\_\_\_\_\_\_\_\_\_\_ smoke is just as dangerous for children sitting nearby.
    \item The \_\_\_\_\_\_\_\_\_\_\_ centre in Buea offers psychological support and medical detoxification for young patients.
    \item Heroin users suffer painful \_\_\_\_\_\_\_\_\_\_\_ symptoms like nausea, anxiety and tremors when they stop taking the drug.
    \item Tar is a powerful \_\_\_\_\_\_\_\_\_\_\_ found in cigarette smoke that directly causes lung cancer.
    \item Sadly, a student from our high school died last year from a heroin \_\_\_\_\_\_\_\_\_\_\_ during the holidays.
\end{enumerate}
\end{exercice}

\begin{exercice}[Grammaire : Modaux de probabilité et conseil (Present \& Past)]
Complétez les phrases en conjuguant le verbe entre parenthèses au modal approprié (\eng{must, can't, could, might, should, ought to, have to}) + forme requise (base verbale ou \eng{have + V-en}). Justifiez brièvement le choix du modal (obligation, interdiction, conseil, déduction forte, possibilité).
\begin{enumerate}
    \item Look at those dark circles under his eyes and the smell of smoke. He \_\_\_\_\_\_\_\_\_\_\_ (smoke) at least a pack a day. \textit{(Déduction forte : certitude quasi absolue)}
    \item The doctor said my grandfather \_\_\_\_\_\_\_\_\_\_\_ (take) his medication for hypertension every morning without fail. \textit{(Obligation médicale stricte)}
    \item You \_\_\_\_\_\_\_\_\_\_\_ (not / mix) alcohol with antidepressants; it is extremely dangerous for the liver. \textit{(Interdiction / Danger)}
    \item Teenagers \_\_\_\_\_\_\_\_\_\_\_ (not / enter) shisha bars legally in Cameroon since the 2023 decree. \textit{(Interdiction légale)}
    \item She looks very tired; she \_\_\_\_\_\_\_\_\_\_\_ (stay) up late studying, or perhaps she is just stressed. \textit{(Possibilité faible)}
    \item The government \_\_\_\_\_\_\_\_\_\_\_ (invest) more in public rehab centres to help the youth. \textit{(Conseil moral / Recommandation forte)}
    \item The patient \_\_\_\_\_\_\_\_\_\_\_ (be) honest with his therapist about his relapse last week, otherwise the treatment won't work. \textit{(Nécessité pour un résultat)}
\end{enumerate}
\end{exercice}

\courssub{Je m'entraîne}

\begin{exercice}[Traduction : Français $\rightarrow$ Anglais (Registre formel / Sensibilisation)]
Traduisez les phrases suivantes en anglais. Soignez le vocabulaire spécifique (santé, addiction) et la structure syntaxique.
\begin{enumerate}
    \item La consommation excessive d'alcool détruit progressivement le foie et augmente le risque de cancers.
    \item Il est impératif que les parents parlent à leurs enfants des dangers de la drogue dès le collège.
    \item Bien que le tabac soit interdit dans les lieux publics, beaucoup de jeunes fument encore la chicha dans les quartiers populaires de Douala.
    \item Si le gouvernement avait investi plus tôt dans la prévention, il y aurait moins de toxicomanes aujourd'hui dans nos rues.
    \item Les personnes âgées souffrent souvent de solitude, ce qui peut les pousser à abuser de médicaments tranquillisants.
\end{enumerate}
\end{exercice}

\begin{exercice}[Transformation : Voix passive et reformulation lexicale]
Transformez les phrases actives en phrases passives (conserver le temps verbal). Puis, pour les phrases 6 à 8, reformulez en remplaçant le mot souligné par un synonyme du champ lexical (nominalisation ou verbe plus précis).
\begin{enumerate}
    \item Tobacco kills more than 8 million people every year worldwide.
    \item Doctors are treating hundreds of young patients for shisha-related lung infections this month.
    \item The Ministry of Health has launched a new awareness campaign against banga consumption.
    \item Peer pressure often pushes teenagers to try their first cigarette.
    \item Alcohol abuse has destroyed many promising careers in our community.
    \item The \textbf{addict} stole money to buy his daily dose. (\textit{Remplacer \textbf{addict} par \eng{substance abuser} ou \eng{person with substance use disorder}})
    \item Smoking \textbf{damages} the lungs irreversibly. (\textit{Remplacer \textbf{damages} par \eng{impairs} ou \eng{compromises}})
    \item We need more \textbf{treatment} centres in rural areas. (\textit{Remplacer \textbf{treatment} par \eng{rehabilitation} ou \eng{recovery}})
\end{enumerate}
\end{exercice}

\begin{exercice}[Texte à trous : Cloze test (Grammaire + Lexique)]
Complétez le texte suivant en conjuguant les verbes entre parenthèses au temps qui convient OU en choisissant le mot manquant parmi les options proposées (a, b, c).
\begin{textebox}[Gap-fill Text: « A Community Response »]
“Last year, the village chief of a small locality near Bamenda \textbf{(1. decide)} to act after three young men \textbf{(2. die)} from a drug overdose in just one month. He \textbf{(3. call)} a meeting where elders \textbf{(4. speak)} openly about the shame and fear \textbf{(5. \_\_\_\_\_\_\_\_)} (a) which (b) who (c) that gripped the families. ‘We \textbf{(6. cannot / sit)} idle while our future \textbf{(7. destroy)},’ he declared. Since then, a local vigilance committee \textbf{(8. set up)}. They \textbf{(9. patrol)} the streets at night and \textbf{(10. report)} suspicious dealers to the gendarmerie. So far, the strategy \textbf{(11. work)}: no new cases \textbf{(12. \_\_\_\_\_\_\_\_)} (a) were reported (b) have been reported (c) are reported. However, the chief insists that repression \textbf{(13. \_\_\_\_\_\_\_\_)} (a) is not enough (b) not enough (c) isn't enough; education and job creation \textbf{(14. \_\_\_\_\_\_\_\_)} (a) is (b) are (c) be essential for a lasting solution.”
\end{textebox}
\end{exercice}

\begin{exercice}[Appariement : Lexique $\leftrightarrow$ Définition + Production]
Associez chaque terme anglais (1--10) à sa définition française (A--J). Ensuite, choisissez \textbf{5 termes} et écrivez une phrase personnelle en anglais pour chacun (contexte camerounais si possible).
\begin{center}
\begin{tabular}{|c|l|c|l|}
\hline
\textbf{Terme} & \textbf{Anglais} & \textbf{Code} & \textbf{Définition française} \\
\hline
1 & Addiction & A & Perte de mémoire et de capacités cognitives liée à l'âge ou à la toxicité. \\
2 & Withdrawal & B & Maladie du foie caractérisée par une fibrose irréversible. \\
3 & Cirrhosis & C & Dépendance physique et psychologique à une substance. \\
4 & Dementia & D & Inhalation de la fumée d'autrui par un non-fumeur. \\
5 & Secondhand smoke & E & Substance provoquant le cancer. \\
6 & Carcinogen & F & Ensemble des symptômes physiques à l'arrêt d'une drogue. \\
7 & Rehabilitation & G & Processus de rétablissement et de réinsertion sociale/médicale. \\
8 & Overdose & H & Prise excessive de drogue entraînant un coma ou la mort. \\
9 & Peer pressure & I & Influence du groupe d'amis pour adopter un comportement risqué. \\
10 & Awareness campaign & J & Action de communication pour informer et prévenir. \\
\hline
\end{tabular}
\end{center}
\textit{Exemple de production attendue :} \eng{The awareness campaign in our school helped students understand the risks of shisha.}
\end{exercice}

\courssub{Je me prépare au Bac}

\begin{exercice}[Compréhension écrite type Baccalauréat (Texte original + Questions)]
Lisez attentivement le texte ci-dessous et répondez aux questions qui suivent en anglais.
\begin{textebox}[Reading Comprehension: « The Hidden Cost of Shisha in Accra's Youth »]
“In the bustling neighbourhoods of Osu and East Legon in Accra, Ghana, a new social ritual has taken root among university students and young professionals: shisha smoking. Every evening, cafes fill with clouds of sweet-smelling vapour — apple, mint, grape — masking a harsh reality. According to a 2023 report by the Ghana Health Service, shisha use among young adults aged 18--25 has surged by 40\% in five years. Many users, like 22-year-old Kwame, a student at the University of Ghana, believe the water pipe filters out toxins. ‘It’s just flavoured steam,’ he argues, ‘much safer than the cigarettes my father smoked.’ Medical experts strongly disagree. Dr. Ama Serwah, a pulmonologist at Korle-Bu Teaching Hospital, explains: ‘A typical one-hour shisha session involves 200 puffs, compared to 20 puffs for a cigarette. The user inhales smoke equivalent to 100 cigarettes. The charcoal used to heat the tobacco produces high levels of carbon monoxide and heavy metals.’ The consequences are already visible. Emergency wards report rising cases of acute carbon monoxide poisoning and severe respiratory infections among young non-cigarette smokers. Beyond health, the economic burden is heavy: treating tobacco-related diseases costs Ghana’s NHS millions of cedis annually. The Food and Drugs Authority has banned shisha importation, but enforcement remains weak as smugglers bring it through porous borders. Civil society groups now demand stricter licensing for cafes and mandatory health warnings on shisha menus, hoping to clear the air for the next generation.”
\end{textebox}
\textbf{Questions :}
\begin{enumerate}
    \item \textbf{True or False?} Justify each answer with a specific quote from the text (max 15 words per quote).
    \begin{enumerate}
        \item Shisha smoke is safer than cigarette smoke because it passes through water.
        \item The Ghanaian authorities have completely stopped shisha from entering the country.
    \end{enumerate}
    \item \textbf{Multiple Choice.} Choose the correct answer (a, b, or c).
    According to Dr. Ama Serwah, why is a one-hour shisha session particularly dangerous?
    \begin{itemize}
        \item[a)] The flavours contain addictive chemicals.
        \item[b)] The volume of smoke inhaled is massively higher than a cigarette.
        \item[c)] The water pipe explodes frequently due to charcoal heat.
    \end{itemize}
    \item \textbf{Short Answers.} Answer in your own words (complete sentences, max 20 words each).
    \begin{enumerate}
        \item What two specific toxic substances are produced by the burning charcoal?
        \item What two measures do civil society groups recommend to reduce shisha use?
    \end{enumerate}
    \item \textbf{Vocabulary in Context.} Find in the text words or phrases that mean the following:
    \begin{enumerate}
        \item Increased rapidly (paragraph 1, verb phrase): \_\_\_\_\_\_\_\_\_\_\_
        \item Breathing difficulties / lung diseases (paragraph 3, noun phrase): \_\_\_\_\_\_\_\_\_\_\_
        \item Not strict / not effective (paragraph 4, adjective): \_\_\_\_\_\_\_\_\_\_\_
    \end{enumerate}
\end{enumerate}
\end{exercice}

\begin{exercice}[Traduction type Baccalauréat : Anglais $\rightarrow$ Français]
Traduisez en français le passage suivant (lignes 1 à 12 du texte ci-dessus, de “In the bustling neighbourhoods…” jusqu'à “…100 cigarettes.”).
\textbf{Consigne :} Respectez le registre de langue (journalistique/informatif), les temps (passé composé / imparfait), et le vocabulaire médical.
\begin{textebox}[Extract to translate]
“In the bustling neighbourhoods of Osu and East Legon in Accra, Ghana, a new social ritual has taken root among university students and young professionals: shisha smoking. Every evening, cafes fill with clouds of sweet-smelling vapour — apple, mint, grape — masking a harsh reality. According to a 2023 report by the Ghana Health Service, shisha use among young adults aged 18--25 has surged by 40\% in five years. Many users, like 22-year-old Kwame, a student at the University of Ghana, believe the water pipe filters out toxins. ‘It’s just flavoured steam,’ he argues, ‘much safer than the cigarettes my father smoked.’ Medical experts strongly disagree. Dr. Ama Serwah, a pulmonologist at Korle-Bu Teaching Hospital, explains: ‘A typical one-hour shisha session involves 200 puffs, compared to 20 puffs for a cigarette. The user inhales smoke equivalent to 100 cigarettes.’”
\end{textebox}
\end{exercice}

\begin{exercice}[Expression écrite type Baccalauréat : Sujet de composition]
\textbf{Choose ONE of the following topics. Write about 120--150 words.}
\begin{enumerate}
    \item \textbf{Article for School Magazine.}
    Your school is organising a \eng{“Health and Well-being Week”}. You have been asked to write an article for the school magazine warning students against the dangers of \textbf{drug and alcohol abuse}.
    \textit{In your article you should:}
    \begin{itemize}
        \item describe the current situation among youth in your city (Yaoundé, Douala, Bafoussam, etc.);
        \item explain the health consequences (brain, liver, lungs, mental health);
        \item propose concrete solutions (school clubs, parental dialogue, sports, laws).
    \end{itemize}
    \item \textbf{Formal Letter.}
    You are the President of the \eng{“Youth Health Club”} in your school. Write a formal letter to the \textbf{Minister of Public Health} suggesting measures to fight \textbf{shisha smoking among students} near school campuses.
    \textit{Your letter must include:}
    \begin{itemize}
        \item the sender's and recipient's addresses, date, subject line;
        \item a clear description of the problem (proximity of bars, misconception of safety);
        \item three specific recommendations (ban within 500m, ID checks, education).
    \end{itemize}
    \item \textbf{Argumentative Essay.}
    \eng{“Old age should be a time of rest and respect, yet many elderly people in Cameroon suffer from neglect and poor healthcare.”} Do you agree?
    \textit{Develop your argument in 3 paragraphs:} Introduction (thesis), Development (2 arguments + examples: pension issues, family breakdown, lack of geriatric centres), Conclusion (personal opinion + call to action).
\end{enumerate}
\end{exercice}

\newpage

\coursec{Corrigés des exercices}

\begin{corrige}[Exercice 1 — QCM : Compréhension globale du script d'écoute]
\begin{enumerate}
    \item \textbf{Réponse : b)} \eng{Peer pressure and the desire to look mature.} Le locuteur insiste sur le fait que les adolescents commencent à fumer pour \eng{fit in with their friends} et paraître plus âgés, et non par curiosité du goût (a) ni à cause de la publicité télévisée (c), qui est interdite.
    \item \textbf{Réponse : b)} \eng{Cirrhosis and increased risk of liver cancer.} Le script mentionne explicitement la \eng{cirrhosis} (cicatrisation irréversible du foie) et l'augmentation du risque de \eng{liver cancer} comme conséquences de l'abus d'alcool.
    \item \textbf{Réponse : b)} \eng{Altering the reward system, making it hard to feel pleasure naturally.} La drogue modifie le circuit de la récompense (\eng{reward system}) : le cerveau ne produit plus de plaisir naturellement, ce qui crée la dépendance. Les propositions (a) et (c) sont des contrevérités.
    \item \textbf{Réponse : b)} \eng{To ensure regular medical check-ups and monitor medication side effects.} Le rapport recommande des bilans de santé réguliers (\eng{regular medical check-ups}) et la surveillance des effets secondaires des médicaments (\eng{medication side effects}) pour les personnes âgées.
\end{enumerate}
\textbf{Point méthode :} au QCM d'écoute, repérez d'abord les mots-clés de la question (\eng{main reason}, \eng{liver}, \eng{brain}, \eng{elderly}), puis éliminez les distracteurs qui reprennent des mots du texte hors de leur contexte.
\end{corrige}

\begin{corrige}[Exercice 2 — Vrai / Faux justifié par une citation]
\begin{enumerate}
    \item \textbf{F} — L'étude porte sur deux régions, pas seulement Yaoundé : \eng{“a worrying trend in the Centre and Littoral regions.”}
    \item \textbf{V} — \eng{“cannabis, locally known as ‘banga’”}.
    \item \textbf{F} — L'hôpital se trouve à Yaoundé, pas à Douala : \eng{“a psychiatrist at the Jamot Hospital in Yaoundé”}.
    \item \textbf{F} — C'est le contraire : \eng{“The nicotine concentration in one shisha session can equal that of a whole pack of cigarettes.”}
    \item \textbf{V} — Le psychiatre parle de \eng{“a generation at risk of chronic psychosis”} ; la psychose chronique est bien un trouble mental durable.
\end{enumerate}
\textbf{Point de vigilance :} la justification doit être une citation \emph{exacte} du texte, entre guillemets anglais. Ne paraphrasez pas dans la citation : c'est la paraphrase qui vous fait perdre le point au Bac.
\end{corrige}

\begin{corrige}[Exercice 3 — Vocabulaire thématique : Complétion]
\begin{enumerate}
    \item \eng{cirrhosis} — \eng{Long-term alcohol abuse is the leading cause of cirrhosis.} (la cirrhose : cicatrisation du foie).
    \item \eng{secondhand} — \eng{secondhand smoke} (fumée passive inhalée par les personnes proches).
    \item \eng{rehabilitation} — \eng{The rehabilitation centre in Buea offers psychological support.} (centre de désintoxication et de réinsertion).
    \item \eng{withdrawal} — \eng{withdrawal symptoms} (symptômes de sevrage : nausées, anxiété, tremblements).
    \item \eng{carcinogen} — \eng{Tar is a powerful carcinogen.} (substance cancérigène).
    \item \eng{overdose} — \eng{died from a heroin overdose} (surdose mortelle).
\end{enumerate}
\textbf{Mots en trop :} \eng{addiction} et \eng{passive}. Attention au faux ami : \eng{passive smoking} existe, mais ici la collocation correcte avec \eng{smoke} est \eng{secondhand smoke}.
\end{corrige}

\begin{corrige}[Exercice 4 — Grammaire : Modaux de probabilité et conseil]
\begin{enumerate}
    \item \eng{He \textbf{must smoke} at least a pack a day.} — \eng{must + base verbale} exprime une déduction forte au présent, fondée sur des indices visibles (cernes, odeur).
    \item \eng{My grandfather \textbf{has to take} his medication every morning.} — \eng{have to} exprime une obligation externe (prescription médicale), plus objective que \eng{must}.
    \item \eng{You \textbf{must not mix} alcohol with antidepressants.} — \eng{mustn't + base verbale} exprime l'interdiction absolue pour cause de danger. Ne pas confondre avec \eng{don't have to} (absence d'obligation).
    \item \eng{Teenagers \textbf{cannot / may not enter} shisha bars legally.} — \eng{can't} exprime ici l'interdiction légale (\eng{it is not allowed}).
    \item \eng{She \textbf{might have stayed} up late studying.} — \eng{might + have + V-en} exprime une possibilité faible concernant le passé.
    \item \eng{The government \textbf{should invest} / \textbf{ought to invest} more in public rehab centres.} — \eng{should / ought to} expriment le conseil ou la recommandation morale.
    \item \eng{The patient \textbf{has to be} / \textbf{must be} honest with his therapist.} — nécessité absolue pour obtenir un résultat.
\end{enumerate}
\textbf{Rappel :} après un modal, toujours la base verbale sans \eng{to} (sauf \eng{ought to}) ; pour une déduction sur le passé : modal + \eng{have} + participe passé.
\end{corrige}

\begin{corrige}[Exercice 5 — Traduction : Français $\rightarrow$ Anglais]
\begin{enumerate}
    \item \eng{Excessive alcohol consumption gradually destroys the liver and increases the risk of cancer.} — « consommation » se dit \eng{consumption}, pas \eng{consommation} ; « augmente le risque de » = \eng{increases the risk of}.
    \item \eng{It is imperative that parents (should) talk to their children about the dangers of drugs from secondary school onwards.} — après \eng{it is imperative that}, le subjonctif : base verbale (ou \eng{should} + base verbale), jamais de \eng{-s} à la 3\textsuperscript{e} personne.
    \item \eng{Although tobacco is banned in public places, many young people still smoke shisha in the working-class neighbourhoods of Douala.} — \eng{although} et \eng{but} ne se combinent jamais dans la même phrase.
    \item \eng{If the government had invested earlier in prevention, there would be fewer drug addicts in our streets today.} — conditionnel mixte : \eng{if + past perfect} (passé), \eng{would + base verbale} (présent).
    \item \eng{Elderly people often suffer from loneliness, which can drive them to abuse tranquillisers.} — « souffrir de » = \eng{suffer from} ; « pousser à » = \eng{drive / push somebody to do}.
\end{enumerate}
\textbf{Points de vigilance :} \eng{drugs} au pluriel pour « la drogue » en général ; \eng{elderly people} (jamais \eng{the old people}, trop direct) ; pas de majuscule à \eng{tranquillisers}.
\end{corrige}

\begin{corrige}[Exercice 6 — Voix passive et reformulation lexicale]
\textbf{Voix passive :}
\begin{enumerate}
    \item \eng{More than 8 million people are killed by tobacco every year worldwide.} (présent simple : \eng{is/are + V-en}).
    \item \eng{Hundreds of young patients are being treated for shisha-related lung infections this month.} (présent continu : \eng{is/are being + V-en}).
    \item \eng{A new awareness campaign against banga consumption has been launched by the Ministry of Health.} (present perfect : \eng{has/have been + V-en}).
    \item \eng{Teenagers are often pushed by peer pressure to try their first cigarette.} (présent simple ; \eng{push} est irrégulier en orthographe : \eng{pushed}).
    \item \eng{Many promising careers in our community have been destroyed by alcohol abuse.} (present perfect).
\end{enumerate}
\textbf{Reformulation lexicale :}
\begin{enumerate}
    \setcounter{enumi}{5}
    \item \eng{The substance abuser stole money to buy his daily dose.} / \eng{The person with a substance use disorder stole money to buy his daily dose.}
    \item \eng{Smoking impairs the lungs irreversibly.} / \eng{Smoking compromises lung function irreversibly.}
    \item \eng{We need more rehabilitation centres in rural areas.} / \eng{We need more recovery centres in rural areas.}
\end{enumerate}
\textbf{Rappel :} à la voix passive, le complément d'agent est introduit par \eng{by} ; il peut être omis s'il est évident ou sans importance.
\end{corrige}

\begin{corrige}[Exercice 7 — Texte à trous : Cloze test]
\begin{enumerate}
    \item \eng{decided} — prétérit simple (\eng{last year}).
    \item \eng{died} — prétérit simple, action passée achevée.
    \item \eng{called} — prétérit simple, enchaînement d'actions passées.
    \item \eng{spoke} — prétérit de l'irrégulier \eng{speak -- spoke -- spoken}.
    \item \textbf{(c) that} — pronom relatif pour une chose (\eng{the shame and fear that gripped the families}) ; \eng{who} est réservé aux personnes.
    \item \eng{cannot sit} — modal + base verbale : \eng{We cannot sit idle.}
    \item \eng{is being destroyed} — présent continu passif : notre avenir « est en train d'être détruit ».
    \item \eng{has been set up} — present perfect passif (\eng{since then}) ; \eng{set up} est irrégulier : \eng{set -- set -- set}.
    \item \eng{patrol} — présent simple (habitude actuelle) ; attention : \eng{patrol} double le \eng{l} au prétérit (\eng{patrolled}).
    \item \eng{report} — présent simple, coordonné à \eng{patrol}.
    \item \eng{has worked} — present perfect (\eng{so far}).
    \item \textbf{(b) have been reported} — present perfect passif, cohérent avec \eng{so far}.
    \item \textbf{(c) isn't enough} — contraction correcte ; (b) est agrammatical sans verbe.
    \item \textbf{(b) are} — sujet pluriel coordonné : \eng{education and job creation are essential}.
\end{enumerate}
\end{corrige}

\begin{corrige}[Exercice 8 — Appariement : Lexique $\leftrightarrow$ Définition + Production]
\textbf{Correspondances :} 1--C ; 2--F ; 3--B ; 4--A ; 5--D ; 6--E ; 7--G ; 8--H ; 9--I ; 10--J.

\textbf{Exemples de production personnelle (5 phrases) :}
\begin{enumerate}
    \item \eng{Addiction to tramadol has become a serious problem among young commercial motorbike riders in Douala.}
    \item \eng{When my uncle stopped drinking, he suffered from severe withdrawal symptoms for two weeks.}
    \item \eng{The awareness campaign organised by our school taught us that shisha is not safer than cigarettes.}
    \item \eng{Peer pressure leads many students in Yaoundé to try their first cigarette before the age of fifteen.}
    \item \eng{My grandmother shows early signs of dementia, so our family makes sure she never stays alone.}
\end{enumerate}
\textbf{Point de vigilance :} \eng{addiction} se construit avec \eng{to} (\eng{addiction to drugs}) ; \eng{suffer from} + nom de la maladie.
\end{corrige}

\begin{corrige}[Exercice 9 — Compréhension écrite type Baccalauréat]
\textbf{Barème indicatif (sur 20) :} Q1 : 2 $\times$ 2 pts ; Q2 : 3 pts ; Q3 : 2 $\times$ 3 pts ; Q4 : 3 $\times$ 2 pts ; qualité de la langue : 1 pt.

\textbf{1. True or False}
\begin{enumerate}
    \item \textbf{False.} Quote: \eng{“The user inhales smoke equivalent to 100 cigarettes.”} L'eau ne filtre pas les toxines, contrairement à la croyance de Kwame.
    \item \textbf{False.} Quote: \eng{“enforcement remains weak as smugglers bring it through porous borders.”} L'interdiction existe mais n'est pas appliquée efficacement.
\end{enumerate}
\textbf{2. Multiple Choice : b)} — \eng{“A typical one-hour shisha session involves 200 puffs, compared to 20 puffs for a cigarette.”} C'est le volume de fumée inhalé qui rend la séance dangereuse.

\textbf{3. Short Answers}
\begin{enumerate}
    \item \eng{The burning charcoal produces high levels of carbon monoxide and heavy metals.}
    \item \eng{They recommend stricter licensing for cafes and mandatory health warnings on shisha menus.}
\end{enumerate}
\textbf{4. Vocabulary in Context}
\begin{enumerate}
    \item \eng{has surged} (a fortement augmenté).
    \item \eng{respiratory infections} (infections respiratoires).
    \item \eng{weak} (faible, peu rigoureux — à propos de \eng{enforcement}).
\end{enumerate}
\textbf{Points de vigilance :} citations limitées à 15 mots ; réponses courtes en phrases complètes, avec vos propres mots ; ne recopiez pas tout le paragraphe.
\end{corrige}

\begin{corrige}[Exercice 10 — Traduction type Baccalauréat : Anglais $\rightarrow$ Français]
\textbf{Barème indicatif (sur 20) :} 4 segments $\times$ 4 pts (sens global + vocabulaire) ; fluidité et correction grammaticale du français : 4 pts.

\textbf{Proposition de traduction :}

« Dans les quartiers animés d'Osu et d'East Legon, à Accra, au Ghana, un nouveau rituel social s'est enraciné parmi les étudiants et les jeunes actifs : fumer la chicha. Chaque soir, les cafés se remplissent de nuages de vapeur aux doux parfums — pomme, menthe, raisin — qui masquent une dure réalité. Selon un rapport de 2023 du Service de santé ghanéen, la consommation de chicha chez les jeunes adultes de 18 à 25 ans a bondi de 40 \% en cinq ans. De nombreux consommateurs, comme Kwame, 22 ans, étudiant à l'Université du Ghana, croient que la pipe à eau filtre les toxines. “Ce n'est que de la vapeur aromatisée, affirme-t-il, bien moins dangereuse que les cigarettes que fumait mon père.” Les experts médicaux sont fermement en désaccord. Le Dr Ama Serwah, pneumologue à l'hôpital universitaire de Korle-Bu, explique : “Une séance de chicha typique d'une heure représente 200 bouffées, contre 20 pour une cigarette. Le consommateur inhale une quantité de fumée équivalente à 100 cigarettes.” »

\textbf{Points de vigilance :}
\begin{itemize}
    \item \eng{has taken root} = « s'est enraciné / s'est installé » (passé composé, pas présent).
    \item \eng{young professionals} = « jeunes actifs / jeunes cadres », pas « jeunes professionnels » au sens strict.
    \item \eng{has surged by 40\%} = « a bondi / augmenté de 40 \% » ; ne pas traduire \eng{surge} par « surgir ».
    \item \eng{filters out toxins} = « filtre / élimine les toxines ».
    \item \eng{puffs} = « bouffées » (de fumée), terme consacré.
    \item \eng{strongly disagree} = « sont fermement en désaccord » ; éviter le calque « désapprouvent fortement ».
\end{itemize}
\end{corrige}

\begin{corrige}[Exercice 11 — Expression écrite type Baccalauréat]
\textbf{Barème indicatif (sur 20) :} respect de la tâche et du format : 5 pts ; contenu et idées : 5 pts ; organisation (paragraphes, connecteurs) : 4 pts ; correction grammaticale et vocabulaire : 4 pts ; longueur respectée (120--150 mots) : 2 pts.

\textbf{Modèle 1 — Article for School Magazine :}

\begin{quote}
“Walk through the streets of Yaoundé after school and you will see them: teenagers sharing a cigarette behind a shop, or students entering a shisha bar near the campus. Drug and alcohol abuse is no longer an adult problem; it is knocking at our classroom doors.

The consequences are terrifying. Alcohol slowly destroys the liver and can cause cirrhosis. Cannabis rewires the adolescent brain, damaging memory and concentration. Shisha, wrongly believed to be harmless, delivers as much nicotine as a whole pack of cigarettes in a single session.

What can we do? First, our school should create a health club where students inform their peers. Second, parents must talk openly with their children instead of only punishing them. Finally, sports and cultural activities give us healthy ways to relax.

Our future is in our hands. Let us protect it: say no before it is too late.”
\end{quote}

\textbf{Modèle 2 — Formal Letter :}

\begin{quote}
“Youth Health Club,\\
Government High School Biyem-Assi,\\
P.O. Box 1234, Yaoundé, Cameroon.\\
15th May 2025.\\[4pt]
The Minister of Public Health,\\
Ministry of Public Health,\\
Yaoundé.\\[4pt]
\textbf{Subject: Urgent measures against shisha smoking near school campuses}\\[4pt]
Dear Sir,\\[4pt]
As President of the Youth Health Club, I am writing to express our deep concern about the growing number of shisha bars operating less than two hundred metres from our school. Many students believe shisha is harmless because the smoke passes through water, yet doctors warn that one session equals a hundred cigarettes.\\[4pt]
We respectfully recommend three measures: first, a legal ban on shisha bars within five hundred metres of any school; second, systematic identity checks to prevent minors from entering these establishments; third, funding for educational programmes on tobacco dangers in secondary schools.\\[4pt]
We trust that your Ministry will act swiftly to protect the health of Cameroonian students.\\[4pt]
Yours faithfully,\\[4pt]
Ngo Bell Sandrine,\\
President of the Youth Health Club.”
\end{quote}

\textbf{Modèle 3 — Argumentative Essay (plan développé) :}

\begin{quote}
“\textbf{Introduction:} In Cameroonian tradition, old age was once synonymous with wisdom and respect. Today, however, many elderly people live in neglect, and I strongly believe this situation must change.

\textbf{Development:} Firstly, most retired workers receive a pension so small that it cannot cover food and medicines, and those who worked in the informal sector receive nothing at all. My own grandfather, a former farmer in Bafoussam, depends entirely on his children. Secondly, the traditional family system is breaking down: young people migrate to cities or abroad, leaving their parents alone in villages without care. Moreover, our country has very few geriatric centres, and hospitals are often too expensive for the elderly.

\textbf{Conclusion:} In my opinion, a society is judged by the way it treats its oldest members. The government should create affordable healthcare for the elderly, and families must rediscover their duty of solidarity. Our grandparents built this country; they deserve to age in dignity.”
\end{quote}

\textbf{Points de vigilance :}
\begin{itemize}
    \item Article : titre accrocheur, ton direct (\eng{you}, \eng{we}), paragraphes courts.
    \item Lettre formelle : pas de contractions (\eng{I am}, pas \eng{I'm}) ; formules \eng{Dear Sir,} / \eng{Yours faithfully,} ; objet en gras.
    \item Essai argumentatif : thèse claire en introduction, connecteurs (\eng{firstly, moreover, however, in conclusion}), exemple personnel bienvenu.
    \item Évitez les faux amis : « sensible » = \eng{sensitive} (et non \eng{sensible}) ; « une aide » = \eng{help / assistance} (et non \eng{an aid} au sens courant).
\end{itemize}
\end{corrige}

\newpage

\coursec{Fiche bilan}

\begin{popfigure}
\begin{tikzpicture}[
    mindmap,
    grow cyclic,
    every node/.style={concept, execute at begin node=\hskip0pt},
    root concept/.append style={
        concept color=popblue, minimum size=3.5cm, text width=3.2cm,
        font=\bfseries\large, text=white
    },
    level 1 concept/.append style={
        level distance=4.5cm, sibling angle=360/6,
        minimum size=2.2cm, text width=2.1cm,
        font=\bfseries\normalsize, text=white
    },
    level 2 concept/.append style={
        level distance=3.2cm, sibling angle=360/8,
        minimum size=1.6cm, text width=1.5cm,
        font=\small, text=white
    }
]
\node[root concept, align=center]{Lesson 25\\Listening:\\Health \& Addictions}
    child[concept color=poporange] {node[level 1 concept, align=center]{Theme\\ \& Context}
        child[concept color=poporange!70] {node[level 2 concept, align=center]{Drugs\\ abuse}
        }
        child[concept color=poporange!70] {node[level 2 concept, align=center]{Tobacco\\ \& smoking}
        }
        child[concept color=poporange!70] {node[level 2 concept, align=center]{Alcohol\\ misuse}
        }
        child[concept color=poporange!70] {node[level 2 concept, align=center]{Aging\\ issues}
        }
    }
    child[concept color=popgreen] {node[level 1 concept, align=center]{Key\\ Vocabulary}
        child[concept color=popgreen!70] {node[level 2 concept, align=center]{Addiction\\ withdrawal}
        }
        child[concept color=popgreen!70] {node[level 2 concept, align=center]{Side effects\\ overdose}
        }
        child[concept color=popgreen!70] {node[level 2 concept, align=center]{Secondhand\\ smoke}
        }
        child[concept color=popgreen!70] {node[level 2 concept, align=center]{Rehabilitation\\ prevention}
        }
    }
    child[concept color=poppurple] {node[level 1 concept, align=center]{Listening\\ Strategies}
        child[concept color=poppurple!70] {node[level 2 concept, align=center]{Predicting\\ content}
        }
        child[concept color=poppurple!70] {node[level 2 concept, align=center]{Keywords\\ \& numbers}
        }
        child[concept color=poppurple!70] {node[level 2 concept, align=center]{Speakers'\\ attitude}
        }
        child[concept color=poppurple!70] {node[level 2 concept, align=center]{Note-taking\\ symbols}
        }
    }
    child[concept color=poppink] {node[level 1 concept, align=center]{Grammar\\ Tools}
        child[concept color=poppink!70] {node[level 2 concept, align=center]{Passive\\ voice}
        }
        child[concept color=poppink!70] {node[level 2 concept, align=center]{Reported\\ speech}
        }
        child[concept color=poppink!70] {node[level 2 concept, align=center]{Modals\\ obligation}
        }
        child[concept color=poppink!70] {node[level 2 concept, align=center]{Present\\ Perfect}
        }
    }
    child[concept color=popgold] {node[level 1 concept, align=center]{Question\\ Types}
        child[concept color=popgold!70] {node[level 2 concept, align=center]{True/False\\ justified}
        }
        child[concept color=popgold!70] {node[level 2 concept] {MCQ}
        }
        child[concept color=popgold!70] {node[level 2 concept, align=center]{Short\\ answers}
        }
        child[concept color=popgold!70] {node[level 2 concept] {Gap fill}
        }
    }
    child[concept color=popblue!60] {node[level 1 concept, align=center]{Output\\ Tasks}
        child[concept color=popblue!40] {node[level 2 concept, align=center]{Summary\\ writing}
        }
        child[concept color=popblue!40] {node[level 2 concept, align=center]{Awareness\\ poster}
        }
        child[concept color=popblue!40] {node[level 2 concept, align=center]{Role-play\\ dialogue}
        }
        child[concept color=popblue!40] {node[level 2 concept, align=center]{Debate\\ arguments}
        }
    };
\end{tikzpicture}
\legende{Carte mentale récapitulative : Lesson 25 — Listening: Health, Addictions \& Aging}
\end{popfigure}

\begin{bilan}

\end{bilan}

\courssub{Lexique clé — English / Français}

\begin{center}
\begin{tabularx}{\linewidth}{|X|X|X|}
\hline
\rowcolor{popblueL}
\textbf{English} & \textbf{Français} & \textbf{Collocations / Contexte} \\
\hline
\cle{addiction} & la dépendance, la toxicomanie & \eng{drug addiction}, \eng{to overcome an addiction} \\
\cle{substance abuse} & l'abus de substances & \eng{substance abuse prevention programme} \\
\cle{withdrawal symptoms} & les symptômes de sevrage & \eng{severe withdrawal symptoms} \\
\cle{overdose} & le surdosage & \eng{fatal overdose}, \eng{accidental overdose} \\
\cle{side effects} & les effets secondaires & \eng{harmful side effects}, \eng{long-term side effects} \\
\cle{secondhand smoke} & la fumée passive & \eng{exposure to secondhand smoke} \\
\cle{binge drinking} & la beuverie express / alcoolisation ponctuelle importante & \eng{binge drinking among teenagers} \\
\cle{liver cirrhosis} & la cirrhose du foie & \eng{chronic liver cirrhosis} \\
\cle{rehabilitation} & la réadaptation, la cure de désintoxication & \eng{rehabilitation centre}, \eng{to enter rehab} \\
\cle{peer pressure} & la pression des pairs & \eng{to give in to peer pressure} \\
\cle{dementia} & la démence & \eng{early onset dementia}, \eng{Alzheimer's disease} \\
\cle{life expectancy} & l'espérance de vie & \eng{to increase life expectancy} \\
\cle{awareness campaign} & une campagne de sensibilisation & \eng{to launch an awareness campaign} \\
\cle{prevention} & la prévention & \eng{prevention is better than cure} \\
\hline
\end{tabularx}
\end{center}

\courssub{Points de grammaire indispensables pour le listening (reconnaissance)}

\begin{center}
\begin{tabularx}{\linewidth}{|X|X|X|}
\hline
\rowcolor{popblueL}
\textbf{Notion} & \textbf{Forme / Repérage à l'oral} & \textbf{Fonction dans les textes santé} \\
\hline
\cle{Passive Voice} & \eng{be + V-en} (ex. \eng{are treated}, \eng{was banned}) & Rapports officiels, statistiques : focus sur l'action, pas l'agent. \\
\hline
\cle{Reported Speech} & \eng{said that / reported that / warned that} + \emph{backshift} & Compte-rendu d'études, déclarations d'experts (OMS, ministère). \\
\hline
\cle{Modals} & \eng{must / mustn't / should / have to / can't} & Obligation légale (\eng{must}), conseil fort (\eng{should}), interdiction (\eng{mustn't}). \\
\hline
\cle{Present Perfect} & \eng{have/has + V-en} (ex. \eng{has risen}, \eng{have launched}) & Nouvelles récentes, études publiées, situation actuelle résultant du passé. \\
\hline
\cle{Connectors} & \eng{however, therefore, consequently, whereas, despite} & Structuration du discours : opposition, cause/conséquence, concession. \\
\hline
\cle{Quantifiers} & \eng{a large number of, the majority of, few, little, enough} & Données chiffrées, proportions, gravité du phénomène. \\
\hline
\end{tabularx}
\end{center}

\courssub{Méthode express — Écoute en 3 étapes (Bac)}

\begin{enumerate}
\item \textbf{Anticipation (30'' avant l'écoute)} : Lire le titre, les questions, les images. Cocher le lexique probable (\eng{smoking, liver, elderly, campaign}). Prédire la structure (reportage, interview, PSA\footnote{Public Service Announcement = spot de prévention.}).
\item \textbf{Première écoute (Global)} : Identifier : \textit{Who? What? Where? Tone?} (alarmant, informatif, espoir). Noter \textbf{mots-clés} et \textbf{chiffres} (dates, \%, âges). Ne pas bloquer sur un mot inconnu.
\item \textbf{Deuxième écoute (Détail)} : Répondre aux items précis (V/F justifié, QCM, \textit{Wh-} questions). Vérifier la cohérence. Repérer les \textbf{marqueurs de reformulation} (\eng{in other words, that is to say, meaning that}).
\end{enumerate}

\courssub{Formules utiles pour la production (Speaking / Writing)}

\begin{itemize}
\item \eng{The report highlights the dangers of...} (Le rapport met en lumière les dangers de...)
\item \eng{According to the speaker / the study, ...} (Selon l'orateur / l'étude, ...)
\item \eng{It is crucial / vital / essential to + V} (Il est crucial de...)
\item \eng{Measures must be taken to + V} (Des mesures doivent être prises pour...)
\item \eng{Raising awareness about... is the first step.} (Sensibiliser à... est la première étape.)
\item \eng{Not only ... but also ...} (Non seulement ... mais aussi ... — inversion sujet/aux si en tête).
\end{itemize}



\begin{aretenir}[Checklist avant le Bac — Lesson 25]
$\square$ Je sais orthographier et prononcer les 15 mots de vocabulaire clé (addiction, withdrawal, cirrhosis, dementia, etc.).\\
$\square$ Je reconnais la voix passive à l'oral (\eng{are affected, has been linked to}) et j'en comprends le sens.\\
$\square$ Je repère les modaux d'obligation/interdiction/conseil (\eng{must, should, have to}) dans un discours de prévention.\\
$\square$ J'identifie les connecteurs logiques (\eng{however, therefore, whereas}) pour suivre l'enchaînement des idées.\\
$\square$ J'applique la méthode en 3 étapes : Anticiper $\rightarrow$ Écoute globale $\rightarrow$ Écoute détaillée.\\
$\square$ Je sais prendre des notes efficaces : symboles (\eng{$\uparrow$ $\downarrow$ $\rightarrow$ =}), abréviations (\eng{govt, stats, envt}), mots-clés seulement.\\
$\square$ Je distingue \eng{True / False} justifié par une citation exacte du texte (\eng{"..."}) et non par mon opinion.\\
$\square$ Je sais rédiger un résumé de 50 mots en anglais structuré (Introduction, Points clés, Conclusion).\\
$\square$ Je peux produire un court message de sensibilisation (affiche, slogan, script radio) en réemployant le lexique.\\
$\square$ Je connais les faux-amis pièges : \eng{drug} $\neq$ drogue (médicament aussi), \eng{abuse} $\neq$ abus (mauvais usage), \eng{sensible} $\neq$ sensible (raisonnable).\\
$\square$ Je gère mon stress : si je rate un passage, je passe au suivant et j'infère au contexte global.
\end{aretenir}

\findecours

\end{document}
